% Système poulies-courroie — PCT 3ème — Cours signé OnBuch+
% Programme officiel MINESEC. Compiler avec : tectonic N09-systeme-poulies-courroie.tex
\newcommand{\DOCMATIERE}{PCT}
\newcommand{\DOCNIVEAU}{3ème}
\newcommand{\DOCMODULE}{Module 4 : Technologie et mécanique}
\newcommand{\DOCLECON}{N09}
\newcommand{\DOCTITRE}{Système poulies-courroie}
% OnBuch+ — préambule « cours signés » ENRICHI (compilé avec tectonic / XeLaTeX)
% Système visuel « Pop » d'OnBuch+ : fond crème (jamais blanc), bordures encre
% épaisses, ombres « dures » décalées, titres Archivo Black en capitales.
% Version MATHÉMATIQUES : graphes (pgfplots/tikz), boîtes pédagogiques
% (définition, propriété, méthode, attention, le savais-tu, exercice résolu,
% exercice, TP, bilan), page de garde et signature OnBuch+ ancrée en bas.
%
% Macros attendues dans main.tex AVANT \input{preamble} :
%   \DOCMATIERE  \DOCNIVEAU  \DOCMODULE  \DOCLECON (numéro)  \DOCTITRE
\documentclass[11pt]{article}

\usepackage{fontspec}
\usepackage[french]{babel}
\usepackage[a4paper,margin=2cm,top=3.4cm,bottom=2.8cm,headheight=1.4cm,footskip=1.3cm]{geometry}
\usepackage{amsmath,amssymb,amsfonts}
\usepackage{mathtools} % \xmapsto, \coloneqq... souvent utilisés par les modèles
\usepackage{cancel} % simplifications barrées (bilans de forces)
% \abs{z} (module, valeur absolue) et \norm{u} (norme), utilisés par les modèles
\providecommand{\abs}{}\let\abs\relax\DeclarePairedDelimiter{\abs}{\lvert}{\rvert}
\providecommand{\norm}{}\let\norm\relax\DeclarePairedDelimiter{\norm}{\lVert}{\rVert}
\usepackage[table]{xcolor} % [table] : fournit aussi \rowcolors (alternance de lignes)
\usepackage{pagecolor}
\usepackage{tikz}
\usetikzlibrary{arrows.meta,positioning,calc,shapes.geometric,shapes.misc,shapes.arrows,decorations.pathmorphing,decorations.pathreplacing,decorations.markings,patterns,shadows,mindmap,trees,fit,backgrounds,matrix,intersections,angles,quotes}
\usepackage{pgfplots}
\usepackage[version=4]{mhchem} % équations nucléaires \ce{^{235}_{92}U}
\usepackage[european,siunitx]{circuitikz} % schémas électriques
\pgfplotsset{compat=1.18}
\usepgfplotslibrary{fillbetween,groupplots} % aires entre courbes (calcul intégral)
\usepackage{enumitem}
\usepackage{fancyhdr}
% [most] charge toutes les bibliothèques tcolorbox (listings, minted, raster,
% documentation...) et gonfle inutilement la mémoire TeX sur un document long
% (~300 boîtes) : on ne charge que ce qui est réellement utilisé.
\usepackage[skins,breakable]{tcolorbox}
\usepackage{graphicx}
\usepackage{colortbl}
\usepackage{booktabs}
\usepackage{tabularx}
\usepackage{multirow}
\usepackage{diagbox} % case d'en-tête diagonale des tableaux à double entrée
% Colonnes extensibles centrée / gauche / droite (C, L, R), souvent utilisées par les modèles
\newcolumntype{C}{>{\centering\arraybackslash}X}
\newcolumntype{L}{>{\raggedright\arraybackslash}X}
\newcolumntype{R}{>{\raggedleft\arraybackslash}X}
\usepackage{array}
\usepackage{multicol}
\usepackage{siunitx}
% parse-numbers=false : accepte aussi \SI{2,5 \times 10^{-3}}{...}
\sisetup{locale=FR,output-decimal-marker={,},group-minimum-digits=4,parse-numbers=false}
\DeclareSIUnit\ohm{\ensuremath{\symup{\Omega}}} % U+2126 absent de la police
\DeclareSIUnit\division{div} % réglages d'oscilloscope (ms/div, V/div)
\DeclareSIUnit\tr{tr} % tours (vitesse de rotation, ex. \SI{1500}{\tr\per\minute})
\DeclareSIUnit\molar{M} % concentration molaire (ex. \SI{3}{\molar}, \SI{1}{\micro\molar})
\DeclareSIUnit\an{an} % année (le modèle confond parfois avec \year, un compteur TeX) — ex. \SI{5}{\kilo\meter\per\an}
\DeclareSIUnit\FCFA{FCFA} % franc CFA (ex. \SI{500}{\FCFA\per\kilo\gram})
\DeclareSIUnit\degreeBrix{\textdegree Brix} % teneur en sucre d'un sirop/jus (réfractométrie)
\DeclareSIUnit\month{mois} % le modèle utilise parfois \month, un compteur TeX, comme unité de durée
\DeclareSIUnit\microliter{\micro\liter} % le modèle écrit parfois \microliter en un seul token
\DeclareSIUnit\million{millions} % dénombrement (ex. \SI{15}{\million\per\mL} spermatozoïdes)
\usepackage{qrcode}
% \degree : commande fréquemment utilisée par les modèles pour un angle ou une
% température en dehors de \SI{}{} (non fournie par les paquets chargés).
\providecommand{\degree}{^{\circ}}
% \qed (fin de démonstration), utilisé par les modèles sans amsthm
\providecommand{\qed}{\hfill$\square$}
% \barre : double barre des valeurs interdites dans un tableau de variations (tkz-tab)
\providecommand{\barre}{\ensuremath{\|}}
% environnement proof (amsthm non chargé) utilisé par les modèles
\ifdefined\proof\else\newenvironment{proof}[1][Démonstration]{\par\noindent\textit{#1.}\ }{\qed\par}\fi
\usepackage{lastpage}
\usepackage{hyperref}
\hypersetup{hidelinks}

% ============================================================
% Palette « Pop » OnBuch+ — mêmes hex que la classe P (plus_ui.dart)
% ============================================================
\definecolor{poporange}{HTML}{F59321}
\definecolor{popdark}{HTML}{E07A0C}
\definecolor{popcream}{HTML}{FFF6E8}
\definecolor{popink}{HTML}{1C1714}
\definecolor{poppurple}{HTML}{7A5AE0}
\definecolor{popgreen}{HTML}{2E9E6A}
\definecolor{poppink}{HTML}{E0457B}
\definecolor{popblue}{HTML}{2F7DE1}
\definecolor{popgold}{HTML}{C98A12}
\definecolor{popmuted}{HTML}{8A7F76}
\definecolor{popline}{HTML}{EADFD2}
\definecolor{popgray}{HTML}{8A7F76}
\colorlet{popgrey}{popgray}
% Alias tolérants (le modèle nomme parfois les couleurs différemment)
\colorlet{popwhite}{white}
\colorlet{popblack}{popink}
\colorlet{popred}{poppink}
\colorlet{popyellow}{popgold}
% Teintes claires (fonds de tableaux/encadrés) — le modèle nomme parfois
% les couleurs avec un suffixe L (ex. popblueL) sans les avoir définies.
\colorlet{poporangeL}{poporange!20}
\colorlet{popdarkL}{popdark!20}
\colorlet{poppurpleL}{poppurple!20}
\colorlet{popgreenL}{popgreen!20}
\colorlet{poppinkL}{poppink!20}
\colorlet{popblueL}{popblue!20}
\colorlet{popgoldL}{popgold!20}
\colorlet{popredL}{popred!20}
\colorlet{popgrayL}{popgray!20}
% Teintes claires pour fonds de tableaux / zones
\colorlet{popblueL}{popblue!10!white}
\colorlet{poporangeL}{poporange!14!white}
\colorlet{popgreenL}{popgreen!12!white}
\colorlet{poppurpleL}{poppurple!12!white}
\colorlet{poppinkL}{poppink!12!white}
\colorlet{popgoldL}{popgold!14!white}

\pagecolor{popcream}

% ============================================================
% Typographie
% ============================================================
\defaultfontfeatures{Ligatures=TeX}
\setmainfont{PlusJakartaSans-Medium.ttf}[Path=../../../fonts/,BoldFont=PlusJakartaSans-Bold.ttf]
% Maths en sans-serif (Fira Math) avec chiffres et lettres droites de Plus
% Jakarta, pour que les formules se fondent dans le texte.
\usepackage[math-style=ISO,bold-style=ISO]{unicode-math}
\setmathfont{FiraMath-Regular.otf}[Path=../../../fonts/]
\setmathfont{PlusJakartaSans-Medium.ttf}[Path=../../../fonts/,range={up/num,up/{Latin,latin},\mathup}]
\usepackage{icomma} % virgule décimale sans espace en mode math
% Caractères mathématiques Unicode absents des polices de texte (Plus Jakarta,
% Archivo Black) : rendus par leur équivalent mathématique, en texte comme en
% formule — sinon ils s'affichent en carré vide (ex. « Arithmétique dans ℤ »).
\usepackage{newunicodechar}
\newunicodechar{ℝ}{\ensuremath{\mathbb{R}}}
\newunicodechar{ℤ}{\ensuremath{\mathbb{Z}}}
\newunicodechar{ℂ}{\ensuremath{\mathbb{C}}}
\newunicodechar{ℕ}{\ensuremath{\mathbb{N}}}
\newunicodechar{ℚ}{\ensuremath{\mathbb{Q}}}
\newunicodechar{𝔻}{\ensuremath{\mathbb{D}}}
\newunicodechar{α}{\ensuremath{\alpha}}
\newunicodechar{β}{\ensuremath{\beta}}
\newunicodechar{γ}{\ensuremath{\gamma}}
\newunicodechar{δ}{\ensuremath{\delta}}
\newunicodechar{ε}{\ensuremath{\varepsilon}}
\newunicodechar{θ}{\ensuremath{\theta}}
\newunicodechar{λ}{\ensuremath{\lambda}}
\newunicodechar{μ}{\ensuremath{\mu}}
\newunicodechar{σ}{\ensuremath{\sigma}}
\newunicodechar{τ}{\ensuremath{\tau}}
\newunicodechar{φ}{\ensuremath{\varphi}}
\newunicodechar{ψ}{\ensuremath{\psi}}
\newunicodechar{ω}{\ensuremath{\omega}}
\newunicodechar{Σ}{\ensuremath{\Sigma}}
% majuscules produites par \MakeUppercase dans les titres (α-aminés -> Α)
\newunicodechar{Α}{\ensuremath{\alpha}}
\newunicodechar{Β}{\ensuremath{\beta}}
\newunicodechar{Γ}{\ensuremath{\Gamma}}
\newunicodechar{Δ}{\ensuremath{\Delta}}
\newunicodechar{Ε}{\ensuremath{\varepsilon}}
\newunicodechar{Θ}{\ensuremath{\Theta}}
\newunicodechar{Λ}{\ensuremath{\Lambda}}
\newunicodechar{Μ}{\ensuremath{\mu}}
\newunicodechar{Τ}{\ensuremath{\tau}}
\newunicodechar{Φ}{\ensuremath{\Phi}}
\newunicodechar{Ψ}{\ensuremath{\Psi}}
\newunicodechar{Ω}{\ensuremath{\Omega}}
\newunicodechar{↦}{\ensuremath{\mapsto}}
\newunicodechar{∈}{\ensuremath{\in}}
\newunicodechar{∉}{\ensuremath{\notin}}
\newunicodechar{∧}{\ensuremath{\wedge}}
\newunicodechar{⇔}{\ensuremath{\Leftrightarrow}}
\newunicodechar{⟺}{\ensuremath{\Longleftrightarrow}}
\newunicodechar{⇒}{\ensuremath{\Rightarrow}}
\newunicodechar{⟹}{\ensuremath{\Longrightarrow}}
\newunicodechar{∘}{\ensuremath{\circ}}
\newunicodechar{∪}{\ensuremath{\cup}}
\newunicodechar{∩}{\ensuremath{\cap}}
\newunicodechar{≡}{\ensuremath{\equiv}}
\newunicodechar{⊂}{\ensuremath{\subset}}
\newunicodechar{⊥}{\ensuremath{\perp}}
\newunicodechar{∃}{\ensuremath{\exists}}
\newunicodechar{∀}{\ensuremath{\forall}}
\newunicodechar{∛}{\ensuremath{\sqrt[3]{\,}}}
\newunicodechar{∜}{\ensuremath{\sqrt[4]{\,}}}
\newunicodechar{⃗}{\ensuremath{{}^{\to}}}
\newunicodechar{ⁿ}{\ifmmode^{n}\else\textsuperscript{\ensuremath{n}}\fi}
\newunicodechar{ˣ}{\ifmmode^{x}\else\textsuperscript{\ensuremath{x}}\fi}
\newunicodechar{ᵇ}{\ifmmode^{b}\else\textsuperscript{\ensuremath{b}}\fi}
\newunicodechar{ʳ}{\ifmmode^{r}\else\textsuperscript{\ensuremath{r}}\fi}
\newunicodechar{ᵘ}{\ifmmode^{u}\else\textsuperscript{\ensuremath{u}}\fi}
\newunicodechar{ʸ}{\ifmmode^{y}\else\textsuperscript{\ensuremath{y}}\fi}
\newunicodechar{ᵃ}{\ifmmode^{a}\else\textsuperscript{\ensuremath{a}}\fi}
\newunicodechar{ᵉ}{\ifmmode^{e}\else\textsuperscript{\ensuremath{e}}\fi}
\newunicodechar{ᵏ}{\ifmmode^{k}\else\textsuperscript{\ensuremath{k}}\fi}
\newunicodechar{ᵖ}{\ifmmode^{p}\else\textsuperscript{\ensuremath{p}}\fi}
\newunicodechar{ᴺ}{\ifmmode^{N}\else\textsuperscript{\ensuremath{N}}\fi}
\newunicodechar{ᶜ}{\ifmmode^{c}\else\textsuperscript{\ensuremath{c}}\fi}
\newunicodechar{ʰ}{\ifmmode^{h}\else\textsuperscript{\ensuremath{h}}\fi}
\newunicodechar{ᵠ}{\ifmmode^{q}\else\textsuperscript{\ensuremath{q}}\fi}
\newunicodechar{ₙ}{\ifmmode_{n}\else\textsubscript{\ensuremath{n}}\fi}
\newunicodechar{ₐ}{\ifmmode_{a}\else\textsubscript{\ensuremath{a}}\fi}
\newunicodechar{ᵢ}{\ifmmode_{i}\else\textsubscript{\ensuremath{i}}\fi}
\newunicodechar{ᵣ}{\ifmmode_{r}\else\textsubscript{\ensuremath{r}}\fi}
\newunicodechar{ₖ}{\ifmmode_{k}\else\textsubscript{\ensuremath{k}}\fi}
\newunicodechar{ᵦ}{\ifmmode_{\beta}\else\textsubscript{\ensuremath{\beta}}\fi}
\newunicodechar{ₓ}{\ifmmode_{x}\else\textsubscript{\ensuremath{x}}\fi}
\newfontfamily\popdisplay{ArchivoBlack-Regular.ttf}[Path=../../../fonts/]
\newfontfamily\poplogo{SpaceGrotesk-Bold.ttf}[Path=../../../fonts/]
\newfontfamily\popsemi{PlusJakartaSans-SemiBold.ttf}[Path=../../../fonts/]
\newfontfamily\popxbold{PlusJakartaSans-ExtraBold.ttf}[Path=../../../fonts/]
\newfontfamily\popxboldls{PlusJakartaSans-ExtraBold.ttf}[Path=../../../fonts/,LetterSpace=3.0]

\setlist[enumerate]{leftmargin=1.7em,itemsep=5pt,topsep=4pt}
\setlist[itemize]{leftmargin=1.7em,itemsep=4pt,topsep=4pt,label={\color{poporange}\textbullet}}
\setlength{\parindent}{0pt}
\setlength{\parskip}{5pt}
\renewcommand{\arraystretch}{1.25}

% pgfplots : style Pop par défaut
\pgfplotsset{
  every axis/.append style={axis line style={popink,line width=1pt},tick style={popink},
    grid style={popline,line width=0.5pt},label style={font=\small},tick label style={font=\footnotesize}},
  popcurve/.style={poporange,line width=1.8pt,no markers,smooth},
}
% Même style disponible dans un \draw TikZ simple (hors pgfplots)
\tikzset{popcurve/.style={poporange,line width=1.8pt}}
\tikzset{popgrid/.style={popline,very thin}}
\colorlet{popcurve}{poporange} % le nom du style est parfois utilisé comme couleur

% ============================================================
% Pill / squiggle
% ============================================================
\newcommand{\poppill}[3]{%
  \tikz[baseline=-0.55ex]\node[draw=popink,line width=1.2pt,rounded corners=2.6pt,%
    fill=#2,inner xsep=6.5pt,inner ysep=3pt]{%
    \popxboldls\fontsize{7}{7}\selectfont\color{#3}\MakeUppercase{#1}};%
}
\newcommand{\popsquiggle}[1]{%
  \tikz[baseline=-0.4ex]{
    \draw[#1,line width=2.6pt,line cap=round]
      plot [smooth] coordinates {(0,0.09) (0.34,0.01) (0.68,0.21) (1.02,0.01) (1.36,0.10)};
  }%
}
% Mot-clé surligné (marqueur orange sous le texte)
\newcommand{\cle}[1]{{\popxbold\color{popdark}#1}}

% ============================================================
% En-tête / pied
% ============================================================
\pagestyle{fancy}
\fancyhf{}
\renewcommand{\headrulewidth}{0pt}
\renewcommand{\footrulewidth}{0pt}
\lhead{\raisebox{-0.15ex}{\poplogo\fontsize{13}{13}\selectfont\color{popink}OnBuch\color{poporange}+}}
\rhead{\poppill{\DOCMATIERE\ \textbullet\ \DOCNIVEAU\ \textbullet\ Leçon \DOCLECON}{white}{popink}}
\lfoot{\popxbold\fontsize{7.6}{7.6}\selectfont\color{popmuted}\MakeUppercase{Cours signé OnBuch+}\ \textbullet\ {\popsemi\DOCTITRE}}
\rfoot{\tikz[baseline=-0.6ex]\node[rounded corners=8.5pt,fill=popink,minimum height=17pt,minimum width=17pt,inner xsep=5pt,inner sep=0]{\popxbold\fontsize{8}{8}\selectfont\color{popcream}\thepage\,/\,\pageref*{LastPage}};}

% ============================================================
% Titres
% ============================================================
\newcommand{\chaptitre}[2]{%
  \begin{tcolorbox}[enhanced,colback=poporange,colframe=popink,boxrule=2.6pt,arc=11pt,
      drop shadow=popink,left=15pt,right=15pt,top=12pt,bottom=11pt,before skip=3pt,after skip=18pt]
    \poppill{#2}{popink}{popcream}\par\vspace{9pt}
    {\raggedright\hyphenpenalty=10000\exhyphenpenalty=10000\popdisplay\fontsize{22}{24}\selectfont\color{popcream}\MakeUppercase{#1}\par}
  \end{tcolorbox}
}
% Section numérotée : pastille orange avec le numéro + titre Archivo
\newcounter{popsec}
\newcommand{\coursec}[1]{%
  \stepcounter{popsec}\setcounter{popsub}{0}%
  \par\vspace{16pt}\noindent
  \begin{minipage}{\linewidth}
  \tikz[baseline=-0.7ex]\node[draw=popink,line width=1.6pt,fill=poporange,rounded corners=4pt,minimum size=22pt,inner sep=0]{\popdisplay\fontsize{12}{12}\selectfont\color{popcream}\thepopsec};%
  \hspace{8pt}{\popdisplay\fontsize{15}{17}\selectfont\color{popink}\MakeUppercase{#1}}\par
  \vspace{4pt}\noindent{\color{poporange}\rule{36pt}{3.6pt}}
  \end{minipage}\par\nopagebreak\vspace{8pt}
}
\newcounter{popsub}
\newcommand{\courssub}[1]{%
  \stepcounter{popsub}%
  \par\vspace{9pt}\noindent{\popxbold\fontsize{11.5}{13}\selectfont\color{popink}{\color{poporange}\thepopsec.\thepopsub}\hspace{6pt}#1}\par\nopagebreak\vspace{4pt}
}

% ============================================================
% Boîtes pédagogiques — même recette : fond blanc, bordure épaisse colorée,
% ombre dure encre, badge plein. Titre optionnel [#1].
% ============================================================
\tcbset{popbase/.style={breakable,enhanced,colback=white,boxrule=2.3pt,arc=9pt,
  shadow={3pt}{-3pt}{0pt}{popink},left=14pt,right=14pt,top=11pt,bottom=11pt,before skip=11pt,after skip=15pt,
  fontupper=\popsemi\fontsize{10.6}{14.5}\selectfont\color{popink},
  fontlower=\popsemi\fontsize{10.6}{14.5}\selectfont\color{popink}}}
% Coche dessinée (le glyphe ✓ n'existe pas dans les polices du document)
\newcommand{\popcheck}[1]{\tikz[baseline=-0.5ex]\draw[#1,line width=1.6pt,line cap=round,line join=round](-0.13,0.0)--(-0.04,-0.09)--(0.13,0.1);}
\providecommand{\coche}{\popcheck{popgreen}}
\providecommand{\resultat}[1]{\par\smallskip\noindent\fcolorbox{popgreen}{popgreenL}{\parbox{\dimexpr\linewidth-2\fboxsep-2\fboxrule\relax}{\textbf{Résultat.} #1}}\par\smallskip}
\newcommand{\popboxhead}[3]{\poppill{#1}{#2}{white}\ifx\relax#3\relax\else\hspace{6pt}{\popxbold\fontsize{10.6}{12}\selectfont\color{popink}#3}\fi\par\vspace{7pt}}

\newtcolorbox{aretenir}[1][]{popbase,colframe=popgreen,before upper={\popboxhead{À retenir}{popgreen}{#1}}}
\newtcolorbox{exemplebox}[1][]{popbase,colframe=poporange,before upper={\popboxhead{Exemple}{poporange}{#1}}}
\newtcolorbox{definition}[1][]{popbase,colframe=popblue,before upper={\popboxhead{Définition}{popblue}{#1}}}
\newtcolorbox{propriete}[1][]{popbase,colframe=poppurple,before upper={\popboxhead{Propriété}{poppurple}{#1}}}
\newtcolorbox{methode}[1][]{popbase,colframe=popgold,before upper={\popboxhead{Méthode}{popgold}{#1}}}
\newtcolorbox{attention}[1][]{popbase,colframe=poppink,before upper={\popboxhead{Attention}{poppink}{#1}}}
\newtcolorbox{experience}[1][]{popbase,colframe=popblue,colback=popblueL,before upper={\popboxhead{Expérience}{popblue}{#1}}}
% « Le savais-tu ? » : carte violette pleine (contexte, culture, Cameroun)
\newtcolorbox{savaistu}[1][]{popbase,colback=poppurple,colframe=popink,
  fontupper=\popsemi\fontsize{10.4}{14.2}\selectfont\color{white},
  before upper={\poppill{Le savais-tu\,?}{white}{poppurple}\ifx\relax#1\relax\else\hspace{6pt}{\popxbold\color{white}#1}\fi\par\vspace{7pt}}}
% Exercice résolu : énoncé en haut, \tcblower puis la solution sur fond crème
\newcounter{popexo}\newcounter{popexr}
\newtcolorbox{exoresolu}[1][]{popbase,colframe=poporange,colbacklower=popcream,
  segmentation style={popink,line width=1.2pt,dashed},
  before upper={\stepcounter{popexr}\popboxhead{Exercice résolu \thepopexr}{poporange}{#1}},
  before lower={\poppill{Solution}{popink}{popcream}\par\vspace{6pt}}}
% Exercice (non résolu en place) — la correction est dans la partie Corrigés
\newtcolorbox{exercice}[1][]{popbase,colframe=popink,
  before upper={\stepcounter{popexo}\popboxhead{Exercice \thepopexo}{popink}{#1}}}
% Corrigé
\newtcolorbox{corrige}[1][]{popbase,colframe=popgreen,colback=popgreenL,
  before upper={\popboxhead{Corrigé}{popgreen}{#1}}}
% Objectifs / prérequis / bilan
\newtcolorbox{objectifs}{popbase,colframe=popink,colback=poporangeL,
  before upper={\popboxhead{Objectifs de la leçon}{popink}{}}}
\newtcolorbox{prerequis}{popbase,colframe=popink,colback=popblueL,
  before upper={\popboxhead{Prérequis}{popblue}{}}}
\newtcolorbox{bilan}{popbase,colframe=popink,colback=white,boxrule=2.8pt,
  before upper={\popboxhead{Fiche bilan}{poporange}{L'essentiel à connaître pour l'examen}}}
% Figure encadrée avec légende
\newtcolorbox{popfigure}[1][]{enhanced,breakable=false,colback=white,colframe=popline,boxrule=1.4pt,arc=8pt,
  left=10pt,right=10pt,top=10pt,bottom=8pt,before skip=10pt,after skip=12pt,halign=center,
  lower separated=false,halign lower=center,
  fontlower=\popsemi\fontsize{9}{11.5}\selectfont\color{popmuted},
  #1}
\newcommand{\legende}[1]{\par\vspace{6pt}{\popsemi\fontsize{9}{11.5}\selectfont\color{popmuted}#1}}

% ============================================================
% Signature OnBuch+
% ============================================================
\newcommand{\poplogoinline}[1]{{\poplogo\fontsize{#1}{#1}\selectfont\color{popink}OnBuch\color{poporange}+}}
\newcommand{\signatureonbuch}{%
  \begin{tcolorbox}[enhanced,colback=popink,colframe=popink,boxrule=2.2pt,arc=10pt,
      drop shadow=poporange,left=14pt,right=14pt,top=10pt,bottom=10pt,before skip=0pt,after skip=0pt]
    \tikz[baseline=-0.7ex]\node[circle,fill=popgreen,draw=popcream,line width=1.3pt,minimum size=22pt,inner sep=0]{\popcheck{white}};%
    \hspace{9pt}\begin{minipage}[c]{0.62\linewidth}
      {\popxbold\fontsize{12}{13}\selectfont\color{popcream}Signé\ \textbullet\ {\poplogo OnBuch\color{poporange}+}}\par\vspace{2pt}
      {\raggedright\popsemi\fontsize{8}{10.5}\selectfont\color{popline}\MakeUppercase{Équipe pédagogique OnBuch+}\\\MakeUppercase{Conforme au programme officiel MINESEC}\par}
    \end{minipage}\hfill
    {\poplogo\fontsize{22}{22}\selectfont\color{popcream}OnBuch\color{poporange}+}
  \end{tcolorbox}%
}

% ============================================================
% Page de garde
% #1 titre · #2 matière · #3 niveau · #4 module · #5 n° leçon · #6 durée ·
% #7 description / objectifs (texte ou itemize)
% ============================================================
\newcommand{\pagegarde}[7]{%
  \thispagestyle{empty}
  \newgeometry{a4paper,margin=1.7cm,top=1.6cm,bottom=1.4cm}%
  \begin{tikzpicture}[remember picture,overlay]
    \fill[poporange!18] ([xshift=-3.2cm,yshift=-3.6cm]current page.north east) circle (4.6cm);
    \fill[poppurple!14] ([xshift=2.4cm,yshift=7.6cm]current page.south west) circle (3.8cm);
  \end{tikzpicture}%
  {\poplogo\fontsize{30}{30}\selectfont\color{popink}OnBuch\color{poporange}+}\hfill
  \raisebox{0.6ex}{\poppill{Cours signé \textbullet\ Premium}{popink}{popcream}}\par
  \vspace{1pt}\popsquiggle{poppurple}\par
  \vspace{8pt}
  \begin{tcolorbox}[enhanced,colback=poporange,colframe=popink,boxrule=2.8pt,arc=13pt,
      drop shadow=popink,left=18pt,right=18pt,top=12pt,bottom=13pt,after skip=10pt]
    \poppill{#2 \textbullet\ #3}{popink}{popcream}\hspace{5pt}\poppill{#4}{white}{popink}\par\vspace{8pt}
    {\popdisplay\fontsize{14}{14}\selectfont\color{popink}LEÇON #5}\par\vspace{4pt}
    {\raggedright\hyphenpenalty=10000\exhyphenpenalty=10000\popdisplay\fontsize{25}{27}\selectfont\color{popcream}\MakeUppercase{#1}\par}
    \vspace{8pt}
    {\popxbold\fontsize{9.5}{9.5}\selectfont\color{popink}\MakeUppercase{Durée conseillée : #6}}
  \end{tcolorbox}
  \begin{tcolorbox}[enhanced,colback=white,colframe=popink,boxrule=2.2pt,arc=10pt,
      drop shadow=popink,left=16pt,right=16pt,top=10pt,bottom=10pt,after skip=8pt]
    \poppill{Dans cette leçon}{poppurple}{white}\par\vspace{5pt}
    \popsemi\fontsize{9.3}{12.6}\selectfont\color{popink}#7
  \end{tcolorbox}
  \noindent\begin{minipage}[t]{0.487\linewidth}
    \begin{tcolorbox}[enhanced,colback=popgreen,colframe=popink,boxrule=2.2pt,arc=10pt,
        drop shadow=popink,left=11pt,right=11pt,top=10pt,bottom=10pt,height=3.1cm,valign=center]
      \tcbox[colback=white,colframe=popink,boxrule=1.3pt,arc=5pt,boxsep=0pt,nobeforeafter,
        left=3pt,right=3pt,top=3pt,bottom=3pt,valign=center]{\qrcode[height=1.9cm]{https://play.google.com/store/apps/details?id=com.luvvix.onbuch&hl=fr}}%
      \hspace{8pt}\begin{minipage}[c]{0.5\linewidth}
        {\popdisplay\fontsize{11}{12.5}\selectfont\color{white}\MakeUppercase{Télécharge OnBuch}}\par\vspace{4pt}
        {\raggedright\popsemi\fontsize{8.4}{11}\selectfont\color{white}Gratuit \textbullet\ Google Play\\Tuteur IA, annales, concours, résultats\par}
      \end{minipage}
    \end{tcolorbox}
  \end{minipage}\hfill
  \begin{minipage}[t]{0.487\linewidth}
    \begin{tcolorbox}[enhanced,colback=poppurple,colframe=popink,boxrule=2.2pt,arc=10pt,
        drop shadow=popink,left=11pt,right=11pt,top=10pt,bottom=10pt,height=3.1cm,valign=center]
      \tikz[baseline=-0.7ex]\node[circle,fill=white,draw=popink,line width=1.3pt,minimum size=1.6cm,inner sep=0]{\popdisplay\fontsize{24}{24}\selectfont\color{poppurple}+};%
      \hspace{8pt}\begin{minipage}[c]{0.6\linewidth}
        {\popdisplay\fontsize{11}{12.5}\selectfont\color{white}\MakeUppercase{Passe à OnBuch+}}\par\vspace{4pt}
        {\raggedright\popsemi\fontsize{8.4}{11}\selectfont\color{white}Tous les cours signés, Tuteur IA illimité, fiches PDF — onglet OnBuch+ de l'app\par}
      \end{minipage}
    \end{tcolorbox}
  \end{minipage}
  \vspace{8pt}
  \popcontenu
  \vfill
  \signatureonbuch
  \restoregeometry
  \newpage
}

% Rangée « Ce cours contient » (page de garde)
\newcommand{\poptile}[3]{% #1 couleur #2 gros texte #3 légende
  \begin{tcolorbox}[enhanced,colback=#1,colframe=popink,boxrule=2pt,arc=9pt,shadow={2.5pt}{-2.5pt}{0pt}{popink},
      left=6pt,right=6pt,top=9pt,bottom=9pt,halign=center,valign=center,height=1.75cm,width=\linewidth,nobeforeafter]
    {\popdisplay\fontsize{12.5}{13}\selectfont\color{white}\MakeUppercase{#2}}\par\vspace{3pt}
    {\centering\popsemi\fontsize{7.8}{9.5}\selectfont\color{white}#3\par}
  \end{tcolorbox}}
\newcommand{\popcontenu}{%
  \par\noindent{\popxboldls\fontsize{7.5}{7.5}\selectfont\color{popmuted}CE COURS SIGNÉ CONTIENT}\par\vspace{7pt}
  \noindent\begin{minipage}[t]{0.235\linewidth}\poptile{popblue}{Cours}{complet, illustré, pas à pas}\end{minipage}\hfill
  \begin{minipage}[t]{0.235\linewidth}\poptile{poporange}{Méthodes}{et exercices résolus}\end{minipage}\hfill
  \begin{minipage}[t]{0.235\linewidth}\poptile{poppink}{Exercices}{type examen + corrigés}\end{minipage}\hfill
  \begin{minipage}[t]{0.235\linewidth}\poptile{popgreen}{Fiche bilan}{l'essentiel à retenir}\end{minipage}\par}

% Sommaire
\newtcolorbox{sommaire}{enhanced,colback=white,colframe=popink,boxrule=2.2pt,arc=10pt,
  drop shadow=popink,left=18pt,right=18pt,top=13pt,bottom=13pt,before skip=6pt,after skip=20pt,
  before upper={\poppill{Sommaire}{poppurple}{white}\par\vspace{9pt}\popsemi\fontsize{10.6}{14.5}\selectfont\color{popink}}}

% Signature de fin de cours — poussée tout en bas de la dernière page
\newcommand{\findecours}{%
  \par\vfill\nopagebreak
  \begin{center}\popsquiggle{poporange}\end{center}\vspace{4pt}
  \signatureonbuch
}
\AtBeginDocument{\def\llbracket{\mathopen{[\mkern-3.5mu[}}\def\rrbracket{\mathclose{]\mkern-3.5mu]}}} % intervalles d'entiers (glyphe ⟦ absent de la police)
% Glyphes absents de Fira Math : redéfinis à partir de glyphes disponibles
\AtBeginDocument{%
  \def\setminus{\mathbin{\backslash}}\let\smallsetminus\setminus
  \def\vdots{\mathord{\vbox{\baselineskip3.2pt\lineskiplimit0pt\kern4pt\hbox{.}\hbox{.}\hbox{.}}}}%
  \def\ddots{\mathinner{\mkern1mu\raise6.5pt\hbox{.}\mkern2mu\raise3.6pt\hbox{.}\mkern2mu\raise0.7pt\hbox{.}\mkern1mu}}%
  \def\triangle{\mathord{\scalebox{0.9}{$\symup{\Delta}$}}}%
  \def\nabla{\mathord{\raisebox{\depth}{\scalebox{1}[-1]{$\symup{\Delta}$}}}}%
  \def\checkmark{\popcheck{popdark}}%
  \def\star{\mathbin{\ast}}\def\bigstar{\mathord{\ast}}%
  \def\diamond{\mathbin{\scalebox{0.6}{$\Diamond$}}}%
  \def\bigcup{\mathop{\mathchoice{\raisebox{-0.3em}{\scalebox{1.7}{$\cup$}}}{\scalebox{1.25}{$\cup$}}{\cup}{\cup}}\displaylimits}%
  \def\bigcap{\mathop{\mathchoice{\raisebox{-0.3em}{\scalebox{1.7}{$\cap$}}}{\scalebox{1.25}{$\cap$}}{\cap}{\cap}}\displaylimits}%
  \def\lesseqqgtr{\mathrel{\lessgtr}}\def\bigoplus{\mathop{\oplus}\displaylimits}%
}
\providecommand{\ssi}{si et seulement si} % abréviation fréquente
\AtBeginDocument{\providecommand{\wideparen}{\overparen}} % arc de cercle

%% FIGFIX-BEGIN v5 — correctifs globaux des figures (docs/onbuch-generation/tools/figfix)
\usepackage{adjustbox}\usepackage{etoolbox}\tcbuselibrary{raster}
% toute figure TikZ/pgfplots plus large que la ligne est réduite à la largeur disponible
\BeforeBeginEnvironment{tikzpicture}{\begin{adjustbox}{max width=\linewidth}}
\AfterEndEnvironment{tikzpicture}{\end{adjustbox}}
% légendes pgfplots lisibles par-dessus les courbes
\pgfplotsset{every axis/.append style={legend style={fill=white,fill opacity=0.92,text opacity=1,draw=black!15,inner sep=3pt}}}
% étiquettes d'arêtes (syntaxe edge["..."]) avec fond blanc
\tikzset{every edge quotes/.append style={fill=white,inner sep=1.2pt}}
%% FIGFIX-END
\begin{document}

\pagegarde{Système poulies-courroie}{PCT}{3ème}{Module 4 : Technologie et mécanique}{N09}{1 séance (1 h chacune)}{Cette leçon présente le système poulies-courroie, un mécanisme très répandu pour transmettre un mouvement de rotation entre deux arbres éloignés. L'élève apprendra à le définir, à en identifier les avantages et les inconvénients, et à connaître les méthodes de correction des glissements.
\vspace{4pt}
\begin{multicols}{2}\small
\begin{itemize}
\item À la fin de cette leçon, je sais définir un système poulies-courroie et nommer ses constituants.
\item À la fin de cette leçon, je sais identifier un système poulies-courroie dans des objets de la vie courante (moulin, groupe électrogène, machine à coudre, ventilateur de voiture).
\item À la fin de cette leçon, je sais décrire le principe de fonctionnement : transmission du mouvement par adhérence.
\item À la fin de cette leçon, je sais citer les avantages et les inconvénients du système poulies-courroie.
\item À la fin de cette leçon, je sais expliquer ce qu'est le glissement d'une courroie et ses causes.
\item À la fin de cette leçon, je sais proposer les méthodes de correction des glissements (tension, galet tendeur, courroie crantée).
\end{itemize}\end{multicols}}

\begin{sommaire}
\begin{multicols}{2}
\begin{enumerate}
\item Observation et définition du système poulies-courroie
\item Fonctionnement et sens de rotation
\item Avantages du système poulies-courroie
\item Inconvénients du système poulies-courroie
\item Méthodes de correction des glissements
\item Sécurité et entretien
\item Activité d'intégration
\item Méthodes et exercices résolus
\item Exercices
\item Corrigés des exercices
\item Fiche bilan
\end{enumerate}
\end{multicols}
\end{sommaire}

\begin{objectifs}
\begin{itemize}
\item À la fin de cette leçon, je sais définir un système poulies-courroie et nommer ses constituants.
\item À la fin de cette leçon, je sais identifier un système poulies-courroie dans des objets de la vie courante (moulin, groupe électrogène, machine à coudre, ventilateur de voiture).
\item À la fin de cette leçon, je sais décrire le principe de fonctionnement : transmission du mouvement par adhérence.
\item À la fin de cette leçon, je sais citer les avantages et les inconvénients du système poulies-courroie.
\item À la fin de cette leçon, je sais expliquer ce qu'est le glissement d'une courroie et ses causes.
\item À la fin de cette leçon, je sais proposer les méthodes de correction des glissements (tension, galet tendeur, courroie crantée).
\item À la fin de cette leçon, je sais interpréter un schéma de transmission poulies-courroie et déterminer le sens de rotation des poulies.
\item À la fin de cette leçon, je sais réaliser un calcul simple de rapport de transmission à partir des diamètres des poulies.
\item À la fin de cette leçon, je sais appliquer les consignes de sécurité liées à l'utilisation des machines à courroie.
\end{itemize}
\end{objectifs}

\begin{prerequis}
\begin{itemize}
\item Notion de mouvement de rotation et de vitesse de rotation (tours par minute)
\item Notion de transmission de mouvement (leçons de technologie de 4ème : roues de friction, engrenages simples)
\item Notion de force et d'action mécanique
\item Lecture et interprétation d'un schéma technique simple
\item Calculs de proportionnalité (rapports, produits en croix)
\end{itemize}
\end{prerequis}

\newpage

\coursec{Observation et définition du système poulies-courroie}

Au quartier, le menuisier fait tourner sa scie circulaire grâce à un moteur électrique placé à côté : entre les deux, une simple courroie en caoutchouc relie deux poulies. Quand la courroie est usée ou détendue, la lame ralentit dès qu'on pousse le bois : on dit qu'elle « patine ». Comment ce système transmet-il le mouvement ? Pourquoi glisse-t-il parfois, et comment y remédier ? C'est ce que nous allons découvrir dans cette leçon.

\courssub{Mise en situation : observons autour de nous}

\textbf{Une observation guidée.}

Regarde bien les trois machines suivantes, que l'on rencontre facilement au Cameroun :

\begin{itemize}
\item le \cle{moulin à maïs} du quartier : un moteur (thermique ou électrique) fait tourner la meule qui écrase les grains ;
\item l'\cle{alternateur} d'une voiture ou d'une moto-taxi : il est entraîné par le moteur pour produire l'électricité qui recharge la batterie ;
\item la \cle{machine à coudre} à pédale du tailleur : le mouvement du pied fait tourner l'aiguille.
\end{itemize}

\begin{experience}[Observation d'une maquette ou d'images]
\textbf{Matériel :} une maquette de transmission (ou des images de moulin à maïs, d'alternateur de voiture, de machine à coudre), une fiche d'observation.

\textbf{Protocole :}
\begin{enumerate}
\item Observer chaque machine et repérer les organes qui tournent.
\item Identifier la source d'énergie (moteur, pédale).
\item Suivre le « chemin » du mouvement, de la source jusqu'à l'outil (meule, alternateur, aiguille).
\item Noter les organes communs aux trois machines.
\end{enumerate}

\textbf{Observations :} dans les trois cas, on trouve :
\begin{itemize}
\item une roue fixée sur l'arbre du moteur (ou de la pédale) ;
\item une autre roue fixée sur l'arbre de l'outil ;
\item une bande souple et fermée (en caoutchouc ou en cuir) qui passe sur les deux roues.
\end{itemize}

\textbf{Interprétation :} quand le moteur tourne, sa roue entraîne la bande souple, et la bande entraîne la seconde roue : le mouvement de rotation est \emph{transmis} d'un arbre à l'autre, alors que les deux arbres sont \emph{éloignés} l'un de l'autre.

\textbf{Conclusion :} les trois machines utilisent le même mécanisme de transmission : le \cle{système poulies-courroie}.
\end{experience}

\begin{attention}[Ne pas confondre]
Ne confonds pas la \cle{poulie} (la roue, souvent en métal, avec une gorge ou un rebord) et la \cle{courroie} (la bande souple qui relie les poulies). La poulie est un organe rigide ; la courroie est un organe flexible. Au BEPC, inverser les deux termes dans une définition coûte des points !
\end{attention}

\courssub{Les constituants du système poulies-courroie}

L'observation précédente nous permet de dresser la liste des organes qui composent tout système poulies-courroie.

\begin{aretenir}[Les quatre constituants]
Un système poulies-courroie comprend toujours :
\begin{enumerate}
\item une \cle{poulie menante} (ou \cle{poulie motrice}) : elle est fixée sur l'arbre du moteur ; c'est elle qui \emph{donne} le mouvement ;
\item une \cle{poulie menée} (ou \cle{poulie réceptrice}) : elle est fixée sur l'arbre de l'outil ; c'est elle qui \emph{reçoit} le mouvement ;
\item une \cle{courroie} : bande souple et fermée (en boucle) qui passe sur les deux poulies et relie leurs mouvements ;
\item deux \cle{arbres} (axes) \emph{éloignés} l'un de l'autre, sur lesquels les poulies sont fixées (clavetées ou serrées) pour tourner solidairement avec eux.
\end{enumerate}
\end{aretenir}

\textbf{Rôle de chaque élément : la chaîne du mouvement.}

Le mouvement suit toujours le même itinéraire, comme une chaîne de solidarité :
\[
\text{moteur} \;\longrightarrow\; \text{poulie menante} \;\longrightarrow\; \text{courroie} \;\longrightarrow\; \text{poulie menée} \;\longrightarrow\; \text{outil}.
\]
La poulie menante, entraînée par le moteur, \emph{entraîne la courroie} par frottement ; la courroie, en se déplaçant, \emph{entraîne la poulie menée}, qui fait tourner l'arbre de l'outil (la meule du moulin, le rotor de l'alternateur, le volant de la machine à coudre).

\begin{popfigure}
\begin{center}
\begin{tikzpicture}[scale=1.0, every node/.style={font=\small}]
% Poulie menante (petite, à gauche)
\draw[fill=popblueL, thick, popblue] (0,0) circle (0.8);
\draw[fill=white, thick] (0,0) circle (0.12);
\fill (0,0) circle (0.04);
% Poulie menée (grande, à droite)
\draw[fill=poporangeL, thick, poporange] (4.6,0) circle (1.4);
\draw[fill=white, thick] (4.6,0) circle (0.12);
\fill (4.6,0) circle (0.04);
% Courroie : brins (tangentes approximatives)
\draw[thick, popdark] (0,0.8) -- (4.6,1.4);
\draw[thick, popdark] (0,-0.8) -- (4.6,-1.4);
% Flèches de rotation
\draw[->, thick, popblue] (-0.55,0.75) arc (130:60:0.85);
\draw[->, thick, poporange] (5.9,0.6) arc (50:-30:1.55);
% Flèche de déplacement du brin tendu
\draw[->, thick, popgreen] (1.1,0.98) -- (2.3,1.16);
% Légendes
\node[popblue, below] at (0,-1.1) {\textbf{Poulie menante}};
\node[popblue, below] at (0,-1.5) {(motrice, côté moteur)};
\node[poporange, below] at (4.6,-1.7) {\textbf{Poulie menée}};
\node[poporange, below] at (4.6,-2.1) {(réceptrice, côté outil)};
\node[popgreen, above] at (2.3,1.35) {\textbf{Brin tendu}};
\node[popmuted, below] at (2.3,-1.35) {\textbf{Brin mou}};
\node[popdark] at (6.6,0.4) {\textbf{Courroie}};
\draw[->, popdark] (6.3,0.25) -- (5.2,0.9);
% Arbres
\node[popmuted, above] at (0,0.95) {arbre moteur};
\node[popmuted, above] at (4.6,1.55) {arbre de l'outil};
\end{tikzpicture}
\end{center}
\legende{Schéma d'un système poulies-courroie. La poulie menante (à gauche) tourne dans le sens de la flèche bleue : elle tire le brin supérieur de la courroie, qui devient le \textbf{brin tendu}, tandis que le brin inférieur est le \textbf{brin mou}. La courroie entraîne alors la poulie menée (à droite) dans le même sens de rotation.}
\end{popfigure}

\courssub{Définition du système poulies-courroie}

\begin{definition}[Système poulies-courroie]
Un \cle{système poulies-courroie} est un mécanisme qui \textbf{transmet un mouvement de rotation entre deux arbres éloignés} au moyen d'une \textbf{courroie souple} passant sur \textbf{deux poulies}.
\end{definition}

Décortiquons cette définition pour bien la retenir :

\begin{itemize}
\item « \emph{transmet un mouvement de rotation} » : le système ne crée pas le mouvement (c'est le moteur qui le crée), il le \emph{transporte} d'un arbre à un autre ;
\item « \emph{entre deux arbres éloignés} » : c'est la grande utilité du système ! Le moteur et l'outil n'ont pas besoin d'être côte à côte ; la courroie peut mesurer plusieurs mètres ;
\item « \emph{au moyen d'une courroie souple} » : la souplesse de la courroie permet d'absorber les chocs et les vibrations, contrairement à une chaîne ou à des engrenages ;
\item « \emph{passant sur deux poulies} » : les poulies guident la courroie et la tendent.
\end{itemize}

\begin{methode}[Identifier les organes d'une transmission sur un schéma]
Face à un schéma ou à une machine réelle, procède toujours dans cet ordre :
\begin{enumerate}
\item \textbf{Repérer la source d'énergie} : cherche le moteur (électrique ou thermique), la pédale ou la manivelle ;
\item \textbf{Identifier la poulie menante} : c'est la poulie fixée sur l'arbre de cette source d'énergie ;
\item \textbf{Suivre la courroie} : la bande souple fermée qui part de la poulie menante ;
\item \textbf{Identifier la poulie menée} : c'est la poulie que la courroie atteint en second, fixée sur l'arbre de l'outil ;
\item \textbf{Vérifier} : le mouvement doit suivre le chemin moteur $\rightarrow$ poulie menante $\rightarrow$ courroie $\rightarrow$ poulie menée $\rightarrow$ outil.
\end{enumerate}
\end{methode}

\begin{exemplebox}[Des exemples de la vie courante]
\begin{itemize}
\item \textbf{Le groupe électrogène du quartier} : le moteur thermique entraîne l'alternateur par une courroie ; c'est cet alternateur qui produit le courant électrique pendant les délestages.
\item \textbf{La machine à coudre à pédale} : le pied du tailleur fait tourner une grande poulie (poulie menante) ; une courroie en cuir entraîne la petite poulie du volant (poulie menée), qui actionne l'aiguille.
\item \textbf{Le ventilateur du moteur de voiture} : sur les anciens véhicules, le moteur entraîne par courroie le ventilateur qui refroidit le radiateur, ainsi que la pompe à eau et l'alternateur.
\end{itemize}
Dans chaque cas, retrouve les quatre constituants : source d'énergie, poulie menante, courroie, poulie menée.
\end{exemplebox}

\courssub{Le principe de fonctionnement : la transmission par adhérence}

\textbf{Pourquoi la courroie entraîne-t-elle la poulie ?}

Posons une question simple : la courroie n'est ni collée, ni attachée aux poulies. Alors pourquoi ne glisse-t-elle pas en permanence ? La réponse tient en un mot : le \cle{frottement}.

\begin{experience}[Mettre en évidence l'adhérence]
\textbf{Matériel :} une maquette poulies-courroie (ou deux bobines et un élastique large), un marqueur.

\textbf{Protocole :}
\begin{enumerate}
\item Marquer un point blanc sur la poulie menante et un point sur la courroie, face à face.
\item Tendre correctement la courroie, puis faire tourner lentement la poulie menante d'un tour complet.
\item Observer les deux points : le point de la courroie accompagne-t-il celui de la poulie ?
\item Détendre fortement la courroie et recommencer.
\end{enumerate}

\textbf{Observations :}
\begin{itemize}
\item Courroie bien tendue : les deux points restent ensemble ; la poulie menée tourne régulièrement.
\item Courroie détendue : le point de la poulie « file » devant celui de la courroie ; la poulie menée tourne lentement ou s'arrête, alors que la poulie menante continue de tourner.
\end{itemize}

\textbf{Interprétation :} quand la courroie est bien tendue, elle est fortement plaquée contre la gorge des poulies. Les forces de \emph{frottement} (appelées ici \emph{adhérence}) entre le caoutchouc de la courroie et le métal de la poulie empêchent le glissement : la courroie et les poulies avancent ensemble, « solidairement ». Quand la courroie est détendue, la pression de contact diminue, l'adhérence devient insuffisante et la courroie \emph{glisse} : c'est le « patinage » observé chez le menuisier.

\textbf{Conclusion :} la transmission du mouvement dans un système poulies-courroie se fait \textbf{par adhérence (frottement)} entre la courroie et les poulies.
\end{experience}

\begin{propriete}[Principe de la transmission par adhérence]
Dans un système poulies-courroie, le mouvement est transmis \textbf{par adhérence} : ce sont les forces de frottement entre la courroie et la gorge des poulies qui entraînent la courroie, puis la poulie menée. Cette adhérence n'est suffisante que si la courroie est \textbf{bien tendue} et en \textbf{bon état} (ni usée, ni graisseuse).
\end{propriete}

Ce principe explique à la fois la souplesse du système (pas de choc brutal, la courroie peut glisser légèrement en cas de surcharge, ce qui protège le moteur) et sa faiblesse (le glissement devient un défaut quand il est permanent : on parle de \cle{glissement} ou de « patinage »).

\courssub{Vocabulaire : brin tendu, brin mou, poulie motrice et poulie réceptrice}

Pour décrire correctement un système poulies-courroie (notamment dans une copie de BEPC), il faut maîtriser un vocabulaire précis.

\begin{definition}[Brin tendu et brin mou]
La courroie forme une boucle fermée : elle possède donc \textbf{deux brins}, c'est-à-dire deux portions rectilignes entre les deux poulies.
\begin{itemize}
\item Le \cle{brin tendu} est le brin que la poulie menante \emph{tire} : il est fortement tendu, car c'est lui qui transmet l'effort ;
\item le \cle{brin mou} est le brin qui \emph{revient} vers la poulie menante : il est moins tendu, parfois légèrement flottant.
\end{itemize}
\end{definition}

\textbf{Comment repérer le brin tendu ?} C'est très simple : regarde le sens de rotation de la poulie menante. Le brin tendu est toujours celui \emph{vers lequel} la surface de la poulie menante se dirige en quittant la courroie... Autre astuce : c'est le brin qui « part » du côté où la poulie menante entraîne la courroie. Sur notre schéma, la poulie menante tourne dans le sens des aiguilles d'une montre : c'est donc le brin supérieur qui est tiré, il est tendu ; le brin inférieur est mou.

\begin{aretenir}[Les synonymes à connaître]
\begin{center}
\begin{tabularx}{\linewidth}{|X|X|}
\hline
\rowcolor{popblueL} \textbf{Poulie qui donne le mouvement} & \textbf{Poulie qui reçoit le mouvement} \\
\hline
poulie \textbf{menante} & poulie \textbf{menée} \\
\hline
poulie \textbf{motrice} & poulie \textbf{réceptrice} \\
\hline
poulie \textbf{conductrice} & poulie \textbf{conduite} \\
\hline
\end{tabularx}
\end{center}
Ces trois paires de termes désignent exactement les mêmes organes. Au BEPC, tu peux utiliser indifféremment « menante/menée » ou « motrice/réceptrice ».
\end{aretenir}

\textbf{La gorge de la poulie et la courroie trapézoïdale.}

Pour améliorer l'adhérence, la plupart des poulies ne sont pas plates : elles possèdent une \cle{gorge} en forme de V, dans laquelle vient se loger une \cle{courroie trapézoïdale} (sa coupe a la forme d'un trapèze). Coincée dans la gorge, la courroie frotte sur \emph{deux flancs} au lieu d'un seul : l'adhérence est bien meilleure et la courroie ne risque pas de sortir de la poulie.

\begin{popfigure}
\begin{center}
\begin{tikzpicture}[scale=1.6]
% Corps de la poulie (coupe)
\fill[popblueL] (-2.2,-0.9) rectangle (2.2,0.9);
\draw[thick, popblue] (-2.2,-0.9) rectangle (2.2,0.9);
% Gorge en V évidée
\fill[white] (-0.9,0.9) -- (-0.25,-0.35) -- (0.25,-0.35) -- (0.9,0.9) -- cycle;
\draw[thick, popblue] (-0.9,0.9) -- (-0.25,-0.35) -- (0.25,-0.35) -- (0.9,0.9);
% Courroie trapézoïdale en place
\fill[poporangeL] (-0.72,0.9) -- (-0.22,-0.13) -- (0.22,-0.13) -- (0.72,0.9) -- cycle;
\draw[thick, poporange] (-0.72,0.9) -- (-0.22,-0.13) -- (0.22,-0.13) -- (0.72,0.9) -- cycle;
% Axe de la poulie
\draw[dashed, popmuted] (-2.6,-0.9) -- (2.6,-0.9);
\node[popmuted, right, font=\small] at (2.6,-0.9) {axe};
% Légendes
\node[poporange, font=\small, above] at (0,0.95) {\textbf{Courroie trapézoïdale}};
\node[popblue, font=\small] at (-1.7,0.1) {\textbf{Gorge}};
\draw[->, popblue] (-1.35,0.05) -- (-0.55,0.35);
\node[popblue, font=\small, below] at (0,-0.95) {\textbf{Poulie (coupe)}};
% Flancs en contact
\draw[<->, popgreen, thick] (-0.62,0.55) -- (-0.32,0.02);
\node[popgreen, font=\footnotesize, left] at (-0.75,0.3) {flanc};
\end{tikzpicture}
\end{center}
\legende{Coupe d'une poulie à gorge avec sa courroie trapézoïdale en place. La courroie (en orange) se loge dans la gorge en V de la poulie (en bleu) et frotte sur ses deux flancs : l'adhérence est renforcée. Attention : le fond de la courroie ne doit \emph{jamais} toucher le fond de la gorge, sinon elle ne frotte plus sur les flancs et patine.}
\end{popfigure}

\begin{attention}[Pièges fréquents]
\begin{itemize}
\item Dire que « la courroie est collée à la poulie » est \textbf{faux} : c'est le \emph{frottement} (l'adhérence), et non la colle, qui transmet le mouvement.
\item Le brin tendu n'est pas forcément le brin du dessus : tout dépend du \emph{sens de rotation} de la poulie menante. Repère-le toujours à partir de la flèche de rotation.
\item Une courroie trapézoïdale trop usée touche le fond de la gorge : elle ne frotte plus sur les flancs et patine, même bien tendue. C'est pour cela qu'on la remplace.
\end{itemize}
\end{attention}

\courssub{Sens de rotation}

\begin{propriete}[Sens de rotation (courroie non croisée)]
Dans un système poulies-courroie à courroie \textbf{non croisée} (montage standard), la poulie menante et la poulie menée tournent \textbf{dans le même sens} (toutes deux dans le sens des aiguilles d'une montre, ou toutes deux dans le sens inverse).
\end{propriete}

\begin{attention}[Courroie croisée]
Si la courroie est croisée (en huit), les deux poulies tournent en \textbf{sens opposés}. Ce montage est rare car il use la courroie plus vite au point de croisement. Au BEPC, sauf mention « courroie croisée », on considère toujours le montage standard (non croisé).
\end{attention}

\courssub{Avantages du système poulies-courroie}

\begin{aretenir}[Avantages principaux]
\begin{enumerate}
\item \textbf{Entraxe variable :} permet de relier deux arbres \emph{éloignés} de plusieurs mètres (ex : moteur au sol, meule en hauteur).
\item \textbf{Absorption des chocs et vibrations :} l'élasticité de la courroie amortit les à-coups du démarrage et les irrégularités du moteur (protection de l'outil et du moteur).
\item \textbf{Silencieux :} pas de cliquetis métallique comme avec une chaîne ou des engrenages.
\item \textbf{Coût modéré et montage simple :} pas d'alignement micrométrique requis, pièces peu onéreuses.
\item \textbf{Sécurité par glissement :} en cas de surcharge brutale (bourrage de la scie), la courroie patine au lieu de casser l'arbre ou de griller le moteur (fusible mécanique).
\end{enumerate}
\end{aretenir}

\courssub{Inconvénients du système poulies-courroie}

\begin{aretenir}[Inconvénients principaux]
\begin{enumerate}
\item \textbf{Glissement (patinage) :} la transmission n'est pas \emph{rigoureusement synchrone} ; le rapport de vitesse n'est pas constant (perte de 1 à 3 \% environ). Interdit pour la distribution moteur (sauf courroie crantée).
\item \textbf{Usure et vieillissement :} la courroie s'allonge, se craquelle, durcit (chaleur, huile, ozone) ; elle doit être surveillée et remplacée périodiquement.
\item \textbf{Entretien de la tension :} il faut vérifier et retendre régulièrement (tendeur ou déplacement du moteur).
\item \textbf{Encombrement radial :} les poulies et la courroie occupent de la place autour des arbres.
\item \textbf{Sensibilité à l'environnement :} huile, graisse, eau, poussière dégradent l'adhérence.
\end{enumerate}
\end{aretenir}

\courssub{Méthodes de correction des glissements}

Le glissement (patinage) est le défaut majeur. Voici comment le prévenir ou le corriger :

\begin{methode}[Corriger et prévenir le glissement]
\begin{enumerate}
\item \textbf{Tendre correctement la courroie :} utiliser un \cle{tendeur} (poulie folle sur le brin mou) ou déplacer le moteur sur ses glissières. Règle pratique : au milieu du brin, la flèche de flexion ne doit pas dépasser \SI{10}{mm} par mètre d'entraxe (ou selon préconisation constructeur).
\item \textbf{Utiliser des courroies trapézoïdales et gorges en V :} le coinçage sur les flancs multiplie l'adhérence par rapport à une courroie plate.
\item \textbf{Nettoyer les surfaces de contact :} \emph{jamais} de graisse, d'huile ou de peinture sur la courroie ni dans les gorges. Nettoyer au dégraissant (essence, acétone) et sécher avant montage.
\item \textbf{Remplacer la courroie usée :} si elle est brillante (vitrifiée), craquelée, effilochée ou si elle touche le fond de la gorge $\rightarrow$ changer \emph{immédiatement}. Toujours changer par un modèle de même référence (longueur, profil, nombre de crans).
\item \textbf{Augmenter l'angle d'enroulement :} si possible, ajouter une \cle{poulie folle} (tendeur) sur le brin mou pour augmenter le contact sur la poulie menante.
\item \textbf{Choisir une courroie adaptée :} courroie crantée (synchrone) si le glissement est interdit (ex : distribution automobile, imprimante 3D).
\end{enumerate}
\end{methode}

\courssub{Rapport de transmission (calcul simple)}

Même s'il y a un léger glissement, on définit le \cle{rapport de transmission} (ou rapport de vitesse) par le rapport des diamètres (ou des rayons) des poulies, en supposant « pas de glissement » pour le calcul théorique.

\begin{propriete}[Rapport de vitesse]
Soit $D_1$ le diamètre de la poulie menante et $D_2$ celui de la poulie menée. Soit $N_1$ et $N_2$ leurs vitesses de rotation (en tr/min ou tr/s).
\[
\frac{N_2}{N_1} = \frac{D_1}{D_2} \qquad \text{ou} \qquad N_1 \cdot D_1 = N_2 \cdot D_2
\]
\textbf{Conséquence :} la \textbf{petite poulie tourne plus vite} que la grande. C'est pour cela que sur le moulin à maïs, la poulie du moteur est petite et celle de la meule est grande : la meule tourne moins vite mais avec plus de force (couple).
\end{propriete}

\begin{exoresolu}[Calcul de la vitesse de la meule]
Sur un moulin à maïs, le moteur tourne à $N_1 = \SI{1500}{tr/min}$. La poulie menante (moteur) a un diamètre $D_1 = \SI{80}{mm}$. La poulie menée (meule) a un diamètre $D_2 = \SI{240}{mm}$. Calculer la vitesse de rotation de la meule $N_2$ (théorique, sans glissement).
\tcblower
\textbf{Solution détaillée.}

On applique la formule du rapport de transmission :
\[
N_1 \cdot D_1 = N_2 \cdot D_2
\]
On isole $N_2$ :
\[
N_2 = N_1 \times \frac{D_1}{D_2}
\]
On remplace par les valeurs :
\[
N_2 = 1500 \times \frac{80}{240} = 1500 \times \frac{1}{3} = \SI{500}{tr/min}
\]
\textbf{Réponse :} La meule tourne théoriquement à \SI{500}{tr/min}. En réalité, à cause du glissement, elle tournera un peu moins vite (environ \SI{485}{tr/min} à \SI{495}{tr/min}).
\end{exoresolu}

\courssub{Sécurité et entretien}

\begin{aretenir}[Règles de sécurité et d'entretien]
\begin{itemize}
\item \textbf{Carter de protection :} toute transmission par courroie accessible doit être \emph{carénée} (cache-courroie) pour éviter l'arrachement de vêtements, cheveux ou doigts. C'est obligatoire sur les machines-outils (scie, perceuse, moulin).
\item \textbf{Arrêt machine :} \emph{Jamais} intervenir (tendre, changer, nettoyer) sur une courroie en mouvement. Couper l'alimentation (clé, disjoncteur, bougie) et attendre l'arrêt complet.
\item \textbf{Vérifications périodiques (tous les mois ou 500 h) :}
      \begin{enumerate}
      \item Tension de la courroie (flèche au milieu du brin).
      \item État de surface : pas de craquelures, pas de brillance excessive (vitrification), pas de morceaux arrachés.
      \item Alignement des poulies : les deux poulies doivent être dans le même plan (règle métallique posée sur les faces extérieures). Un désalignement use la courroie en biais et la fait sauter.
      \item Propreté : pas d'huile, de graisse, de sciure ou de poussière dans les gorges.
      \end{enumerate}
\item \textbf{Lubrification :} on \emph{ne graisse jamais la courroie} ! On graisse uniquement les \emph{roulements des arbres} (paliers) selon les préconisations.
\item \textbf{Remplacement :} changer la courroie au premier signe d'usure avancée ou tous les 2 ans / 2000 h de fonctionnement (selon constructeur). Sur un groupe électrogène de quartier, prévoir une courroie de rechange dans la boîte à outils.
\end{itemize}
\end{aretenir}

\begin{exoresolu}[Identifier les organes d'un moulin à maïs]
Sur le moulin à maïs du quartier, le moteur thermique est fixé au sol. Une courroie relie une petite poulie fixée sur l'arbre du moteur à une grande poulie fixée sur l'arbre de la meule. Le moteur tourne dans le sens des aiguilles d'une montre (vu de face), le brin supérieur de la courroie étant celui qui avance vers la meule.

1) Identifier la poulie menante et la poulie menée.

2) Quel est le rôle de la courroie ?

3) Quel brin est tendu lorsque le moulin fonctionne ?

4) Par quel phénomène le mouvement est-il transmis ?
\tcblower
\textbf{Solution détaillée.}

1) La poulie \textbf{menante} (motrice) est la \emph{petite poulie fixée sur l'arbre du moteur}, car c'est le moteur qui fournit le mouvement. La poulie \textbf{menée} (réceptrice) est la \emph{grande poulie fixée sur l'arbre de la meule}, car elle reçoit le mouvement.

2) La courroie est l'organe souple qui \emph{relie les deux poulies} et \emph{transmet le mouvement de rotation} de la poulie menante à la poulie menée, entre les deux arbres éloignés.

3) Le moteur tourne dans le sens des aiguilles d'une montre : la poulie menante \emph{tire} le brin supérieur (celui qui avance vers la meule). Le \textbf{brin supérieur est donc le brin tendu} ; le brin inférieur est le brin mou.

4) Le mouvement est transmis \textbf{par adhérence} (frottement) entre la courroie et les gorges des poulies, à condition que la courroie soit bien tendue et en bon état.
\end{exoresolu}

\begin{savaistu}[La courroie de distribution : une courroie qui peut coûter un moteur]
Dans une voiture, la \cle{courroie de distribution} synchronise le vilebrequin et l'arbre à cames, qui commande l'ouverture des soupapes. C'est une \emph{courroie crantée} (dentée) : pas de glissement possible, le calage reste parfait. Elle s'use progressivement et doit être remplacée environ tous les \SI{100000}{km} (selon les constructeurs). Si elle casse en roulant, les soupapes peuvent percuter les pistons : le moteur est alors \emph{détruit}, et la réparation coûte souvent plus cher qu'une vidange complète du porte-monnaie ! C'est pourquoi les garagistes sérieux de Douala ou de Yaoundé insistent toujours pour la changer à temps. Une simple courroie en caoutchouc armé protège donc tout un moteur.
\end{savaistu}

\begin{exercice}[À toi de jouer]
1) Donne la définition d'un système poulies-courroie en citant ses quatre constituants.

2) Sur la machine à coudre à pédale du tailleur, identifie la source d'énergie, la poulie motrice et la poulie réceptrice.

3) Explique, en deux phrases, pourquoi la scie du menuisier « patine » lorsque sa courroie est détendue.

4) Dessine un schéma simple d'un système poulies-courroie dont la poulie menante tourne dans le sens inverse des aiguilles d'une montre, et indique le brin tendu.

5) Cite trois avantages et deux inconvénients du système poulies-courroie par rapport à une transmission par engrenages.

6) Un moteur de groupe électrogène tourne à \SI{3000}{tr/min}. Sa poulie (menante) a un diamètre de \SI{100}{mm}. L'alternateur (poulie menée) doit tourner à \SI{1500}{tr/min}. Quel doit être le diamètre de la poulie de l'alternateur ? (On néglige le glissement).

7) Donne trois actions d'entretien préventif pour éviter le glissement d'une courroie trapézoïdale.
\end{exercice}

\begin{corrige}[Exercice n]
1) Un système poulies-courroie est un mécanisme qui transmet un mouvement de rotation entre deux arbres éloignés au moyen d'une courroie souple passant sur deux poulies. Les quatre constituants : poulie menante (motrice), poulie menée (réceptrice), courroie, deux arbres éloignés.

2) Source d'énergie : le pied du tailleur (pédale). Poulie motrice : la grande poulie fixée sur l'arbre de la pédale (volant d'inertie). Poulie réceptrice : la petite poulie fixée sur l'arbre du mécanisme d'aiguille.

3) Quand la courroie est détendue, la pression de contact entre la courroie et les gorges des poulies diminue. Les forces de frottement (adhérence) deviennent insuffisantes pour entraîner la poulie menée : la courroie glisse sur la poulie menante, c'est le patinage.

4) Schéma : deux cercles (poulies) reliés par deux lignes parallèles (courroie). Flèche sur la poulie de gauche (menante) dans le sens inverse des aiguilles d'une montre (trigonométrique). Le brin tendu est celui que la poulie menante tire : c'est le brin \emph{inférieur} (qui part vers la droite vers la poulie menée). Le brin supérieur est le brin mou. Les deux poulies tournent dans le même sens (inverse des aiguilles d'une montre).

5) Avantages : entraxe variable (distance), absorption des chocs/vibrations, silencieux, coût faible, sécurité par glissement en surcharge. (Trois attendus).
   Inconvénients : glissement (rapport non constant), usure/vieillissement de la courroie, nécessité de retendre régulièrement, encombrement radial, sensibilité à l'huile/poussière. (Deux attendus).

6) Formule : $N_1 D_1 = N_2 D_2 \Rightarrow D_2 = \frac{N_1 D_1}{N_2} = \frac{3000 \times 100}{1500} = \SI{200}{mm}$.
   La poulie de l'alternateur doit avoir un diamètre de \SI{200}{mm}.

7) Actions préventives : (1) Vérifier et régler la tension régulièrement (flèche normale). (2) Nettoyer les gorges et la courroie (pas d'huile, de graisse, de poussière). (3) Vérifier l'alignement des poulies (même plan). (4) Remplacer la courroie si usée (craquelures, fond de gorge touché, brillance). (Trois attendus).
\end{corrige}

\coursec{Observation et définition du système poulies-courroie}

\courssub{Situation déclenchante : le moteur de la moto-taxi}

Imagine un moto-taxi stationné au carrefour \og Warda \fg{} à Yaoundé. Le moteur tourne au ralenti, mais la roue arrière ne bouge pas. Le conducteur actionne la poignée d'accélérateur : le moteur monte en régime, et soudain la roue arrière se met à tourner, propulsant la moto vers l'avant. Comment le mouvement de rotation du moteur (arbre de sortie de boîte de vitesses) est-il transmis à la roue arrière ? Si tu regardes sous le carter gauche, tu vois deux \cle{poulies} (disques rainurés) reliées par une \cle{courroie} en caoutchouc renforcé. C'est un \textbf{système poulies-courroie}.

\begin{experience}[Découverte du système poulies-courroie sur moto ou maquette]
\textbf{Matériel :} moto-taxi (ou photo détaillée du carter gauche), ou maquette de technologie (deux poulies sur axes parallèles, courroie), feutre, règle.

\textbf{Protocole :}
\begin{enumerate}
\item Identifier les deux poulies : laquelle est reliée au moteur (sortie de boîte) ? Laquelle est reliée à la roue arrière ?
\item Observer la courroie : est-elle plate, trapézoïdale (en V), dentée ? Quel est son état (neuve, usée, craquelée) ?
\item Vérifier si la courroie est \emph{croisée} (forme de 8) ou \emph{non croisée} (parallèle).
\item Faire tourner la roue arrière à la main (moteur arrêté) et observer si le moteur tourne aussi (point mort engagé).
\end{enumerate}

\textbf{Observations :}
\begin{itemize}
\item Deux poulies de diamètres différents : la petite sur le moteur, la grande sur la roue.
\item Une courroie trapézoïdale (en V) non croisée, tendue.
\item Quand on tourne la roue, la poulie moteur tourne aussi : la transmission est \textbf{réversible}.
\end{itemize}

\textbf{Interprétation :} Le système permet de \textbf{transmettre un mouvement de rotation} entre deux axes parallèles éloignés, sans contact direct entre les dents (contrairement aux engrenages). La courroie assure la liaison par \textbf{adhérence} (frottement) dans les gorges des poulies.

\textbf{Conclusion :} On distingue la \cle{poulie menante} (entraînée par le moteur, ici la petite) et la \cle{poulie menée} (qui reçoit le mouvement, ici la grande). La courroie est l'élément intermédiaire flexible.
\end{experience}

\begin{definition}[Système poulies-courroie]
Un \textbf{système poulies-courroie} est un dispositif de transmission de mouvement de rotation entre deux axes parallèles, composé de :
\begin{itemize}
\item deux \textbf{poulies} (disques à gorge) solidaires de leurs arbres respectifs ;
\item une \textbf{courroie} (bande flexible, plate ou trapézoïdale) enroulée sur les deux poulies.
\end{itemize}
La poulie qui reçoit l'énergie du moteur est la \textbf{poulie menante} ; l'autre est la \textbf{poulie menée}.
\end{definition}

\begin{savaistu}[Courroies au Cameroun : de la SONARA à l'atelier]
À la raffinerie \textbf{SONARA} (Limbe), d'énormes courroies trapézoïdales entraînent des pompes centrifuges sur des kilomètres de tuyauterie. Dans ton quartier, le \textbf{soudeur} utilise une meuleuse d'angle : son moteur électrique entraîne le disque par une toute petite courroie plate. À la maison, le \textbf{ventilateur} de plafond ou le \textbf{mélangeur} (blender) fonctionnent aussi par courroie. C'est une technologie universelle, simple et robuste.
\end{savaistu}

\coursec{Fonctionnement et sens de rotation}

Dans la section précédente, nous avons observé le système poulies-courroie et donné sa définition : deux poulies, reliées par une courroie, permettent de transmettre un mouvement de rotation d'un arbre à un autre, sans contact direct entre les arbres. Nous avons distingué la poulie \cle{menante} (celle qui reçoit le mouvement du moteur) et la poulie \cle{menée} (celle qui est entraînée). Une question se pose alors naturellement : quand la poulie menante tourne, \emph{comment} tourne la poulie menée ? Dans le même sens ? En sens inverse ? À la même vitesse ? C'est exactement ce que nous allons découvrir dans cette section, d'abord par l'expérience, puis par le raisonnement, et enfin par le calcul.

\courssub{Expérience de classe : la maquette à deux poulies}

\begin{experience}[Sens de rotation avec une maquette poulies-courroie]
\textbf{Matériel :} une maquette constituée de deux poulies (ou deux roues de vélo, ou deux couvercles circulaires) montées sur des axes fixes parallèles, un élastique large (ou un lacet fermé) jouant le rôle de courroie, un feutre pour marquer un repère sur chaque poulie.

\textbf{Protocole :}
\begin{enumerate}
\item Marquer un point au feutre sur le bord de chaque poulie pour bien voir le sens de rotation.
\item Placer l'élastique \emph{sans le croiser} autour des deux poulies (montage dit \og courroie non croisée \fg{}).
\item Faire tourner lentement la poulie menante dans le sens des aiguilles d'une montre. Observer le sens de rotation de la poulie menée.
\item Retirer l'élastique, le croiser une fois en forme de 8 (montage dit \og courroie croisée \fg{}), puis recommencer la manipulation.
\item Recommencer en diminuant volontairement la tension de l'élastique (poulies plus écartées ou élastique détendu) et observer ce qui se passe.
\end{enumerate}

\textbf{Observations :}
\begin{itemize}
\item Avec la courroie \textbf{non croisée} : la poulie menée tourne dans le \textbf{même sens} que la poulie menante.
\item Avec la courroie \textbf{croisée} : la poulie menée tourne en \textbf{sens inverse} de la poulie menante.
\item Avec un élastique détendu : la poulie menante tourne, mais la poulie menée tourne mal, s'arrête ou ne tourne pas du tout : l'élastique \cle{glisse} sur la poulie.
\end{itemize}

\textbf{Interprétation :} la courroie entraîne la poulie menée grâce à l'\cle{adhérence} entre la courroie et la gorge de la poulie. Si l'adhérence est insuffisante (courroie détendue, gorge lisse ou graissée), la courroie glisse et le mouvement n'est plus transmis correctement.

\textbf{Conclusion :} le sens de rotation de la poulie menée dépend de la façon dont la courroie est montée : non croisée ou croisée.
\end{experience}

\courssub{Constat expérimental : courroie non croisée et courroie croisée}

L'expérience précédente nous donne un résultat fondamental, que tout élève de 3ème doit connaître et savoir utiliser sur un schéma.

\begin{propriete}[Sens de rotation dans un système poulies-courroie]
Propriété établie expérimentalement :
\begin{itemize}
\item Si la courroie est \cle{non croisée} (montage normal, en parallèle), la poulie menante et la poulie menée tournent dans le \textbf{même sens de rotation}.
\item Si la courroie est \cle{croisée} (montage en forme de 8), la poulie menante et la poulie menée tournent en \textbf{sens inverses}.
\end{itemize}
Cette propriété s'applique quel que soit le diamètre des poulies : croiser la courroie change le sens de rotation, mais pas le rapport des vitesses.
\end{propriete}

\begin{popfigure}
\begin{tikzpicture}[scale=1]
% ===== Courroie non croisée =====
\begin{scope}
\draw[thick,popblue,fill=popblueL] (0,0) circle (0.9);
\draw[thick,popblue,fill=popblueL] (3.2,0) circle (0.9);
\fill[popdark] (0,0) circle (0.06);
\fill[popdark] (3.2,0) circle (0.06);
% courroie non croisée
\draw[very thick,poporange] (-0.9,0) -- (2.3,0);
\draw[very thick,poporange] (0.9,0) -- (4.1,0);
\draw[very thick,poporange] (-0.9,0) arc (180:360:0.9);
\draw[very thick,poporange] (0.9,0) arc (0:180:0.9);
\draw[very thick,poporange] (2.3,0) arc (180:360:0.9);
\draw[very thick,poporange] (4.1,0) arc (0:180:0.9);
% flèches de rotation (même sens : trigonométrique)
\draw[-{Stealth[length=3mm]},very thick,popgreen] (0,1.25) arc (90:150:1.25);
\draw[-{Stealth[length=3mm]},very thick,popgreen] (3.2,1.25) arc (90:150:1.25);
\node[popdark,font=\bfseries] at (1.6,-1.7) {Courroie non croisée : même sens};
\end{scope}
% ===== Courroie croisée =====
\begin{scope}[xshift=7.2cm]
\draw[thick,popblue,fill=popblueL] (0,0) circle (0.9);
\draw[thick,popblue,fill=popblueL] (3.2,0) circle (0.9);
\fill[popdark] (0,0) circle (0.06);
\fill[popdark] (3.2,0) circle (0.06);
% courroie croisée (en 8) - points de tangence à 45 deg (rayon 0.9 -> 0.636)
\draw[very thick,poporange] (-0.64,0.64) -- (3.84,-0.64);
\draw[very thick,poporange] (-0.64,-0.64) -- (3.84,0.64);
\draw[very thick,poporange] (-0.64,0.64) arc (135:225:0.9);
\draw[very thick,poporange] (-0.64,-0.64) arc (225:315:0.9);
\draw[very thick,poporange] (3.84,0.64) arc (45:-45:0.9);
\draw[very thick,poporange] (3.84,-0.64) arc (-45:-135:0.9);
% flèches de rotation (sens inverses)
\draw[-{Stealth[length=3mm]},very thick,popgreen] (0,1.25) arc (90:150:1.25); % Menante trigonométrique
\draw[-{Stealth[length=3mm]},very thick,popgreen] (3.2,1.25) arc (90:30:1.25);  % Menée horaire
\node[popdark,font=\bfseries] at (1.6,-1.7) {Courroie croisée : sens inverses};
\end{scope}
\end{tikzpicture}
\legende{Comparaison des deux montages. À gauche, la courroie non croisée : les deux poulies tournent dans le même sens (flèches vertes identiques, sens trigonométrique). À droite, la courroie croisée en forme de 8 : les poulies tournent en sens inverses.}
\end{popfigure}

\courssub{Interprétation : le rôle de l'adhérence}

Pourquoi la courroie entraîne-t-elle la poulie menée ? Faisons un raisonnement simple. Quand la poulie menante tourne, sa gorge frotte contre la courroie. Grâce à l'\cle{adhérence} (le frottement entre la courroie et la gorge), la courroie est entraînée et se déplace comme un tapis roulant. À son tour, la courroie en mouvement frotte contre la gorge de la poulie menée et l'entraîne en rotation.

\begin{aretenir}[La condition indispensable : l'adhérence]
La transmission du mouvement dans un système poulies-courroie repose sur l'\textbf{adhérence} entre la courroie et les gorges des poulies. Si l'adhérence est insuffisante (courroie détendue, usée, mouillée ou graissée), la courroie \cle{glisse} sur les poulies : la poulie menante tourne \og dans le vide \fg{} et la poulie menée ralentit ou s'arrête. C'est le défaut appelé \textbf{glissement}, que nous étudierons en détail dans la section 5.
\end{aretenir}

\begin{savaistu}[L'adhérence dans la vie courante]
Sur une moto-taxi, quand la courroie (ou la chaîne) est détendue, le moteur rugit mais la roue arrière peine à tourner : c'est le glissement. De même, sur un groupe électrogène, une courroie usée fait un bruit de sifflement aigu au démarrage : elle patine dans la gorge de la poulie. Le mécanicien la remplace ou la retend pour rétablir l'adhérence.
\end{savaistu}

\begin{methode}[Déterminer le sens de rotation sur un schéma]
Pour trouver le sens de rotation de la poulie menée sur un schéma, suis ces étapes :
\begin{enumerate}
\item Repérer la \textbf{poulie menante} (celle reliée au moteur) et son sens de rotation (souvent indiqué par une flèche).
\item Repérer le \textbf{brin tendu} de la courroie, c'est-à-dire le côté de la courroie qui part de la poulie menante dans le sens du mouvement.
\item Suivre ce brin avec le doigt jusqu'à la poulie menée : la courroie y arrive par un certain côté (haut ou bas).
\item En déduire le sens de rotation de la poulie menée : le point de la poulie touché par le brin tendu se déplace dans le même sens que la courroie.
\item Vérifier avec la règle : courroie non croisée $\Rightarrow$ même sens ; courroie croisée $\Rightarrow$ sens inverses.
\end{enumerate}
\end{methode}

\courssub{Tailles des poulies et vitesse de rotation}

Reprenons notre maquette, mais cette fois avec deux poulies de \textbf{diamètres différents}. Marquons un repère sur chaque poulie et faisons faire exactement \textbf{un tour} à la petite poulie menante. Que constate-t-on ? La grande poulie menée n'a pas encore terminé son tour ! En effet, la courroie avance d'une longueur égale au tour de la petite poulie ; mais comme la grande poulie a une circonférence plus grande, cette même longueur de courroie ne lui fait parcourir qu'une fraction de tour.

\begin{propriete}[Réduction et multiplication de vitesse]
\begin{itemize}
\item Une \textbf{petite poulie menante} entraînant une \textbf{grande poulie menée} produit une \cle{réduction de vitesse} : la poulie menée tourne plus lentement que la poulie menante.
\item Une \textbf{grande poulie menante} entraînant une \textbf{petite poulie menée} produit une \cle{multiplication de vitesse} : la poulie menée tourne plus vite que la poulie menante.
\end{itemize}
Phrase-clé à retenir : \textbf{la grande poulie tourne toujours plus lentement que la petite.}
\end{propriete}

\begin{popfigure}
\begin{tikzpicture}[scale=1]
% petite poulie menante D1 = 50 mm -> rayon 0.75
\draw[thick,popblue,fill=popblueL] (0,0) circle (0.75);
\fill[popdark] (0,0) circle (0.06);
\draw[thick,popdark,<->] (-0.75,-1.05) -- (0.75,-1.05);
\node[popdark] at (0,-1.35) {$D_1 = \SI{50}{mm}$};
\node[popdark,font=\bfseries] at (0,1.15) {Menante};
% grande poulie menée D2 = 100 mm -> rayon 1.5
\draw[thick,popblue,fill=popblueL] (4.25,0) circle (1.5);
\fill[popdark] (4.25,0) circle (0.06);
\draw[thick,popdark,<->] (2.75,-1.85) -- (5.75,-1.85);
\node[popdark] at (4.25,-2.15) {$D_2 = \SI{100}{mm}$};
\node[popdark,font=\bfseries] at (4.25,1.9) {Menée};
% courroie
\draw[very thick,poporange] (-0.75,0) -- (2.75,0);
\draw[very thick,poporange] (0.75,0) -- (5.75,0);
\draw[very thick,poporange] (-0.75,0) arc (180:360:0.75);
\draw[very thick,poporange] (0.75,0) arc (0:180:0.75);
\draw[very thick,poporange] (2.75,0) arc (180:360:1.5);
\draw[very thick,poporange] (5.75,0) arc (0:180:1.5);
% flèches de rotation
\draw[-{Stealth[length=3mm]},very thick,popgreen] (0,1.0) arc (90:150:1.0);
\draw[-{Stealth[length=3mm]},very thick,popgreen] (4.25,1.75) arc (90:150:1.75);
% indication vitesses
\node[poppurple,font=\bfseries] at (0,0.45) {$N_1$};
\node[poppurple,font=\bfseries] at (4.25,0.5) {$N_2 = \dfrac{N_1}{2}$};
\end{tikzpicture}
\legende{Réduction de vitesse : la petite poulie menante ($D_1 = \SI{50}{mm}$) entraîne la grande poulie menée ($D_2 = \SI{100}{mm}$). Le diamètre est doublé, donc la vitesse est divisée par deux : la grande poulie tourne deux fois moins vite.}
\end{popfigure}

\courssub{Le rapport de transmission}

Pour calculer la vitesse de la poulie menée, on utilise une grandeur appelée \cle{rapport de transmission}, noté $r$.

\begin{definition}[Rapport de transmission]
Le \textbf{rapport de transmission} d'un système poulies-courroie est le quotient :
\[ r = \dfrac{D_1}{D_2} = \dfrac{N_2}{N_1} \]
où :
\begin{itemize}
\item $D_1$ est le diamètre de la poulie \textbf{menante} (en mm) ;
\item $D_2$ est le diamètre de la poulie \textbf{menée} (en mm) ;
\item $N_1$ est la vitesse de rotation de la poulie \textbf{menante} (en tr/min) ;
\item $N_2$ est la vitesse de rotation de la poulie \textbf{menée} (en tr/min).
\end{itemize}
On en tire la formule pratique : $N_2 = \dfrac{D_1}{D_2} \times N_1$.
\end{definition}

\begin{attention}[L'erreur classique : inverser les diamètres]
L'erreur la plus fréquente au BEPC est d'inverser le rapport des diamètres, c'est-à-dire d'écrire $N_2 = \dfrac{D_2}{D_1} \times N_1$. Pour ne jamais te tromper, utilise ce contrôle simple : \textbf{la grande poulie tourne toujours plus lentement que la petite}. Si la poulie menée est plus grande que la menante, ton résultat $N_2$ doit être \emph{plus petit} que $N_1$. Si tu trouves une vitesse plus grande, c'est que tu as inversé le rapport : reprends ton calcul. Autre piège : les deux diamètres doivent être exprimés dans la \textbf{même unité} avant le calcul.
\end{attention}

\begin{exemplebox}[Calcul d'une vitesse de poulie menée]
Une poulie menante de diamètre $D_1 = \SI{60}{mm}$ tourne à $N_1 = \SI{1200}{tr/min}$. Elle entraîne une poulie menée de diamètre $D_2 = \SI{120}{mm}$. Calculons la vitesse $N_2$ de la poulie menée :
\begin{align*}
N_2 &= \dfrac{D_1}{D_2} \times N_1 \\
N_2 &= \dfrac{60}{120} \times 1200 \\
N_2 &= 0,5 \times 1200 = \SI{600}{tr/min}
\end{align*}
Vérification : la poulie menée est deux fois plus grande que la menante, elle tourne donc deux fois moins vite : $1200 \div 2 = \SI{600}{tr/min}$. Le résultat est cohérent.
\end{exemplebox}

\begin{exoresolu}[Courroie d'un groupe électrogène]
Sur un groupe électrogène, le moteur entraîne une poulie menante de diamètre $D_1 = \SI{80}{mm}$ tournant à $N_1 = \SI{3000}{tr/min}$. La courroie, non croisée, entraîne la poulie de l'alternateur de diamètre $D_2 = \SI{200}{mm}$.
\begin{enumerate}
\item Calculer le rapport de transmission $r$.
\item Calculer la vitesse de rotation $N_2$ de l'alternateur.
\item Les deux poulies tournent-elles dans le même sens ? Justifier.
\end{enumerate}
\tcblower
\textbf{Solution détaillée.}
\begin{enumerate}
\item Le rapport de transmission est :
\[ r = \dfrac{D_1}{D_2} = \dfrac{80}{200} = 0,4 \]
\item La vitesse de la poulie menée est :
\begin{align*}
N_2 &= r \times N_1 = \dfrac{D_1}{D_2} \times N_1 \\
N_2 &= 0,4 \times 3000 = \SI{1200}{tr/min}
\end{align*}
Vérification : la poulie menée est plus grande que la menante, elle doit tourner plus lentement. Or $\SI{1200}{tr/min} < \SI{3000}{tr/min}$ : le résultat est cohérent.
\item La courroie est \textbf{non croisée} : d'après la propriété établie expérimentalement, les deux poulies tournent donc dans le \textbf{même sens} de rotation.
\end{enumerate}
\end{exoresolu}

\begin{exercice}[À toi de jouer]
La poulie du moteur d'une machine à laver a un diamètre $D_1 = \SI{40}{mm}$ et tourne à $N_1 = \SI{1400}{tr/min}$. Elle entraîne, par une courroie croisée, la poulie du tambour de diamètre $D_2 = \SI{280}{mm}$.
\begin{enumerate}
\item Calculer la vitesse de rotation $N_2$ du tambour.
\item Indiquer, en le justifiant, si le tambour tourne dans le même sens que le moteur.
\end{enumerate}
\end{exercice}

\begin{corrige}[Exercice : machine à laver]
\begin{enumerate}
\item $N_2 = \dfrac{D_1}{D_2} \times N_1 = \dfrac{40}{280} \times 1400 = \dfrac{1}{7} \times 1400 = \SI{200}{tr/min}$. Le tambour tourne à \SI{200}{tr/min} (vérification : poulie menée plus grande $\Rightarrow$ vitesse plus faible, c'est cohérent).
\item La courroie est \textbf{croisée} : le tambour tourne donc en \textbf{sens inverse} du moteur.
\end{enumerate}
\end{corrige}

\begin{aretenir}[L'essentiel de la section]
\begin{itemize}
\item Courroie \textbf{non croisée} : les deux poulies tournent dans le \textbf{même sens}.
\item Courroie \textbf{croisée} : les deux poulies tournent en \textbf{sens inverses}.
\item Le mouvement est transmis grâce à l'\textbf{adhérence} entre la courroie et les gorges ; sans adhérence, la courroie \textbf{glisse}.
\item La \textbf{grande} poulie tourne toujours plus \textbf{lentement} que la petite.
\item Rapport de transmission : $r = \dfrac{D_1}{D_2} = \dfrac{N_2}{N_1}$, d'où $N_2 = \dfrac{D_1}{D_2} \times N_1$.
\end{itemize}
\end{aretenir}

\coursec{Avantages du système poulies-courroie}

Pourquoi choisir une courroie plutôt que des engrenages (pignons) ou une chaîne ? Observons les situations concrètes.

\courssub{Expérience comparative : courroie vs engrenages}

\begin{experience}[Comparaison courroie / engrenages]
\textbf{Matériel :} deux maquettes : l'une avec deux engrenages en prise directe, l'autre avec deux poulies-courroie (entraxe variable). Dynamomètre ou masse accrochée à la poulie menée pour tester le couple.

\textbf{Protocole :}
\begin{enumerate}
\item Sur la maquette engrenages : tourner la menante. Observer le bruit, les vibrations, la distance fixe entre axes.
\item Sur la maquette poulies-courroie : écarter les poulies (augmenter l'entraxe), retendre la courroie. Tourner la menante. Observer le silence, la possibilité de changer l'entraxe.
\item Essayer de bloquer brutalement la poulie menée (freinage main) sur les deux systèmes. Observer ce qui se passe (patinage courroie vs blocage brutal engrenages).
\end{enumerate}

\textbf{Observations :}
\begin{itemize}
\item Engrenages : bruit de cliquetis, axes fixes, transmission rigide (casse possible en surcharge).
\item Poulies-courroie : silencieux, axes mobiles (entraxe variable), la courroie patine en surcharge (protection du moteur).
\end{itemize}

\textbf{Conclusion :} La courroie offre souplesse, silence et sécurité par rapport aux engrenages rigides.
\end{experience}

\begin{propriete}[Avantages principaux du système poulies-courroie]
\begin{enumerate}
\item \textbf{Entraxe variable :} la distance entre les deux axes peut être grande (jusqu'à plusieurs mètres) et ajustable. Idéal pour moto (moteur fixe, roue mobile avec suspension) ou machines-outils.
\item \textbf{Fonctionnement silencieux :} pas de choc dent contre dent comme chez les engrenages. Confort acoustique (ventilateur, machine à laver).
\item \textbf{Amortissement des chocs et vibrations :} l'élasticité de la courroie (caoutchouc) absorbe les à-coups du moteur (ex: monocylindre moto) et protège la transmission.
\item \textbf{Protection par glissement (fusible mécanique) :} en cas de surcharge ou blocage de la menée, la courroie glisse au lieu de casser l'arbre ou le moteur. C'est une sécurité passive.
\item \textbf{Entretien simple et coût faible :} pas de lubrification obligatoire (contrairement aux chaînes), remplacement facile et peu coûteux.
\item \textbf{Vitesse élevée possible :} les courroies trapézoïdales ou dentées supportent des vitesses périphériques élevées ($> \SI{30}{m/s}$).
\end{enumerate}
\end{propriete}

\begin{savaistu}[La courroie trapézoïdale (courroie en V) : standard industriel]
Au Cameroun, comme partout, la \textbf{courroie trapézoïdale} (profil V) a remplacé la courroie plate sur la plupart des machines (groupes électrogènes, pompes à eau, compresseurs, alternateurs de voiture). Sa section en V augmente la surface de contact dans la gorge : l'adhérence est multipliée par l'effet de coin (\og coinage \fg{}), ce qui permet de transmettre plus de puissance sans glisser, avec moins de tension.
\end{savaistu}

\coursec{Inconvénients du système poulies-courroie}

Malgré ses qualités, le système a des limites qu'un technicien doit connaître pour bien choisir sa transmission.

\begin{propriete}[Inconvénients principaux]
\begin{enumerate}
\item \textbf{Glissement inévitable (fluage) :} même bien tendue, la courroie glisse légèrement (fluage élastique). Le rapport de vitesse \textbf{n'est pas constant} (contrairement aux engrenages ou chaînes). Précision insuffisante pour du positionnement précis (imprimante 3D, commande numérique).
\item \textbf{Usure et vieillissement :} le caoutchouc se craquelle, durcit, s'allonge (fluage plastique) sous l'effet de la chaleur, de l'huile, de l'ozone, des UV. Durée de vie limitée (quelques années).
\item \textbf{Encombrement radial :} les poulies ont un diamètre minimal (pour ne pas trop courber la courroie), ce qui prend de la place radialement.
\item \textbf{Sensibilité à l'environnement :} l'huile, le gras, l'eau, la poussière détruisent l'adhérence et dégradent la gomme. Interdit en milieu gras sans protection.
\item \textbf{Puissance limitée :} pour de très fortes puissances (gros moteurs industriels, camions), on préfère les chaînes ou les engrenages, plus compacts et sans glissement.
\end{enumerate}
\end{propriete}

\begin{attention}[Ne pas confondre glissement et fluage]
\begin{itemize}
\item Le \textbf{glissement} (patinage) : la courroie tourne sur la poulie (adhérence rompue). Défaut grave, bruit, échauffement, destruction rapide.
\item Le \textbf{fluage} (glissement élastique) : micro-déformation cyclique de la courroie dans la gorge due à la différence de tension entre brin tendu et brin mou. Phénomène \textbf{normal, inévitable}, de l'ordre de 1 à 2 \%. Il fait que $N_2$ réel est légèrement inférieur à $N_2$ théorique.
\end{itemize}
Au BEPC, on utilise la formule théorique $N_2 = \frac{D_1}{D_2} N_1$ (sans fluage), sauf si l'énoncé donne un \og rendement \fg{} ou un \og coefficient de glissement \fg{}.
\end{attention}

\coursec{Méthodes de correction des glissements}

Le glissement (patinage) est l'ennemi n°1 de la courroie. Voici comment l'éviter ou le corriger, du plus simple au plus efficace.

\begin{methode}[Corriger / prévenir le glissement d'une courroie]
\begin{enumerate}
\item \textbf{Vérifier et régler la tension (la cause n°1).} Une courroie détendue glisse. Règle pratique : au milieu de l'entraxe, on doit pouvoir fléchir la courroie d'environ \SI{10}{mm} à \SI{15}{mm} par mètre d'entraxe avec le pouce (force modérée). Sur moto : jeu de \SI{10}{mm} à \SI{15}{mm} sur le brin inférieur. Utiliser un \textbf{tendeur} (galet tendeur automatique ou excentrique manuel).
\item \textbf{Nettoyer les gorges et la courroie.} Dégraisser au nettoyant freins ou essence (pas d'huile, pas de graisse, pas de silicone !). Une gorge graissée = glissement garanti.
\item \textbf{Remplacer la courroie usée.} Si elle est craquelée, durcie, brillante (vitrifiée), élargie (fond de gorge usé) : elle n'adhère plus. Changer \textbf{toujours par la même référence} (longueur, profil, nombre de crans si dentée).
\item \textbf{Utiliser une courroie trapézoïdale (profil V) au lieu de plate.} L'effet de coin dans la gorge multiplie l'adhérence par $1/\sin(\alpha/2)$ ($\alpha \approx 36^\circ \Rightarrow$ coefficient $\approx 3,2$). C'est la solution standard.
\item \textbf{Augmenter l'angle d'enroulement.} Si possible, ajouter un \textbf{galet tendeur} (galet fou) sur le brin mou pour augmenter l'enroulement sur la petite poulie (la plus critique). Passer de $150^\circ$ à $180^\circ$ ou plus améliore fortement l'adhérence.
\item \textbf{Aligner les poulies.} Désalignement (axes non parallèles, poulies décalées) = usure latérale + glissement local. Vérifier à la règle ou au laser.
\item \textbf{Choisir une courroie dentée (synchrone) si glissement inacceptable.} Dents en caoutchouc + âme en fibre de verre/aramide + poulies à dents. Transmission \textbf{sans glissement} (rapport constant), mais plus cher, plus fragile aux chocs, bruit plus fort.
\end{enumerate}
\end{methode}

\begin{experience}[TP : Mesurer l'effet de la tension sur le glissement]
\textbf{Matériel :} maquette poulies-courroie avec poulie menante motorisée (moteur CC + variateur), poulie menée freinée par un frein à poudre ou masse pendue, tachymètre optique (ou chronomètre + repère), tendeur à vis.

\textbf{Protocole :}
\begin{enumerate}
\item Régler tension faible (courroie molle). Mesurer $N_1$ et $N_2$. Calculer $r_{mes} = N_2/N_1$. Noter le bruit.
\item Augmenter la tension progressivement. Relever $N_1, N_2$ à chaque palier.
\item Tracer $r_{mes}$ en fonction de la tension (ou de la déflexion).
\end{enumerate}

\textbf{Interprétation attendue :} $r_{mes}$ augmente avec la tension et se rapproche de $D_1/D_2$. Trop tendu $\rightarrow$ roulements qui chauffent, courroie qui casse. Il y a un optimum.
\end{experience}

\coursec{Sécurité et entretien}

Un bon technicien ne se contente pas de calculer : il assure la sécurité et la durabilité.

\begin{aretenir}[Règles de sécurité impératives (BEPC \& vie pro)]
\begin{itemize}
\item \textbf{Arrêt total \& consignation} avant toute intervention (couper le courant, retirer la clé, verrouiller le disjoncteur). Zéro tolérance.
\item \textbf{Carter de protection obligatoire} sur tout système poulies-courroie accessible (moteur, alternateur, ventilateur, machine-outil). Risque majeur : \textbf{prise en main} (doigts, cheveux, vêtements aspirés dans l'engrenage poulie/courroie). Amputation possible.
\item \textbf{Ne jamais tendre/retirer une courroie à la force (levier, marteau)} : risque de casser la poulie (fonte fragile), d'endommager les roulements, de se blesser. Desserrer le tendeur ou écarter le moteur (sur glissières).
\item \textbf{Éloigner les produits inflammables} (essence, solvants) des courroies qui chauffent (échauffement par frottement/hystérésis).
\end{itemize}
\end{aretenir}

\begin{methode}[Entretien préventif d'une transmission par courroie (check-list mensuelle)]
\begin{enumerate}
\item \textbf{Visuel :} inspecter la courroie (fissures transversales, bords effilochés, fond de gorge brillant/vitrifié, morceaux manquants). $\rightarrow$ Remplacer si défaut.
\item \textbf{Tension :} vérifier la déflexion au pouce (ou au tensiomètre électronique pour les pros). Retendre si jeu excessif.
\item \textbf{Alignement :} poser une règle métallique sur les faces des deux poulies : elles doivent être dans le même plan (écart $< \SI{0,5}{mm}$). Corriger si décalage (clés, entretoises, déplacement moteur).
\item \textbf{Propreté :} essuyer gorges et courroie avec un chiffon sec. Pas de produit "resserrant courroie" (bombe aérosol) : ça colle la poussière et détruit le caoutchouc à terme.
\item \textbf{Roulements :} écouter les bruits anormaux (ronronnement, grattement) au niveau des paliers de poulies. Graisser si graisseurs, sinon surveiller.
\item \textbf{Galets tendeurs :} vérifier leur rotation libre (pas de grippage) et leur usure (profil de gorge).
\end{enumerate}
\end{methode}

\begin{savaistu}[Le "truc" du mécanicien camerounais : la courroie de secours]
Tout bon chauffeur de taxi-moto ou de car (transport interurbain) garde \textbf{une courroie de rechange neuve} dans sa boîte à outils, ainsi qu'une \textbf{clé à pipe de 14 ou 17} (écrous de tendeur/moteur). En panne sur la route (courroie cassée), il la change en 15 min sous le manguier. C'est ça, la technologie au service de l'autonomie ! De même, sur un groupe électrogène de quartier (le "groupe" qui alimente le quartier en cas de coupure ENEO), le chef de quartier vérifie la tension de la courroie d'alternateur \textbf{tous les lundis} : c'est la condition pour avoir la lumière le soir.
\end{savaistu}

\begin{exoresolu}[Diagnostic de panne sur pompe à eau]
Une pompe à eau de forage (pompe immergée entraînée par moteur thermique en surface via poulies-courroie) ne débite plus assez. Le moteur tourne à fond, la poulie menante tourne vite, mais la poulie menée (arbre pompe) tourne lentement. La courroie est neuve, bien tendue, non graissée.
\begin{enumerate}
\item Quelle est la cause probable du glissement malgré une courroie neuve et tendue ?
\item Proposer deux solutions techniques pour augmenter l'adhérence sans changer les diamètres des poulies.
\end{enumerate}
\tcblower
\textbf{Solution détaillée.}
\begin{enumerate}
\item Cause probable : \textbf{angle d'enroulement insuffisant sur la petite poulie menante}. Si l'entraxe est court, la courroie n'enroule que $\approx 120^\circ$ sur la petite poulie. La différence de tension entre brin tendu et brin mou ne génère pas assez de force de frottement pour transmettre le couple requis (pompe dure à démarrer, charge élevée).
\item Solutions :
      \begin{enumerate}
      \item Ajouter un \textbf{galet tendeur (galet fou)} sur le brin mou, positionné pour augmenter l'enroulement sur la petite poulie (viser $> 150^\circ$, idéalement $180^\circ$).
      \item Remplacer la courroie plate/trapézoïdale simple par une \textbf{courroie trapézoïdale à crans (courroie dentée intérieure)} ou une \textbf{courroie poly-V (multi-gorges)} : coefficient d'adhérence bien supérieur, plus flexible pour petit rayon.
      \end{enumerate}
\end{enumerate}
\end{exoresolu}

\begin{exercice}[Synthèse : Le ventilateur d'extraction d'une cuisine collective]
Un ventilateur d'extraction (hotte professionnelle) doit tourner à $N_2 = \SI{900}{tr/min}$. Le moteur électrique disponible tourne à $N_1 = \SI{2800}{tr/min}$. On dispose de poulies de diamètres standards : $\SI{50}{mm}$, $\SI{75}{mm}$, $\SI{100}{mm}$, $\SI{150}{mm}$, $\SI{200}{mm}$.
\begin{enumerate}
\item Choisir un couple de poulies $(D_1, D_2)$ parmi les standards pour obtenir $N_2 \approx \SI{900}{tr/min}$. Calculer la vitesse exacte obtenue.
\item La courroie sera non croisée. Si le moteur tourne dans le sens horaire (vue de face), dans quel sens tourne le ventilateur ?
\item Citer deux avantages du système poulies-courroie pour cette application (ventilateur en toiture, moteur en local technique).
\item Indiquer deux vérifications d'entretien mensuel spécifiques à cette installation (poussière, graisse, vibrations).
\end{enumerate}
\end{exercice}

\begin{corrige}[Exercice : Ventilateur d'extraction]
\begin{enumerate}
\item Il faut une réduction de vitesse : $D_1 < D_2$. Rapport visé $r = N_2/N_1 = 900/2800 \approx 0,32$.
      Essais avec standards :
      $D_1=50, D_2=150 \Rightarrow r=0,333 \Rightarrow N_2 = 933 \text{ tr/min}$ (proche).
      $D_1=75, D_2=200 \Rightarrow r=0,375 \Rightarrow N_2 = 1050 \text{ tr/min}$ (trop rapide).
      Choix optimal : \textbf{$D_1 = \SI{50}{mm}$ (menante), $D_2 = \SI{150}{mm}$ (menée)}.
      Vitesse exacte : $N_2 = \frac{50}{150} \times 2800 = \frac{2800}{3} \approx \SI{933}{tr/min}$.
\item Courroie non croisée $\Rightarrow$ \textbf{même sens} que le moteur : \textbf{horaire} (vue de face).
\item Avantages :
      \begin{enumerate}
      \item \textbf{Entraxe long} : moteur au sol (local technique), ventilateur sur le toit (distance plusieurs mètres) : courroie idéale, arbre long impossible.
      \item \textbf{Silence \& amortissement} : pas de transmission de vibrations du moteur vers la hotte (confort cuisine), démarrage progressif (pas de choc sur pales).
      \end{enumerate}
\item Vérifications mensuelles :
      \begin{enumerate}
      \item \textbf{Propreté gorges/courroie} : cuisine = graisses en suspension (aérosols). Dégraissage impératif (nettoyant dégraissant non agressif pour caoutchouc).
      \item \textbf{Tension \& alignement} : vibrations du ventilateur (déséquilibre pales) + dilatation thermique (toit chaud) $\rightarrow$ détente courroie + désalignement poulies.
      \end{enumerate}
\end{enumerate}
\end{corrige}

\begin{aretenir}[Bilan complet de la leçon : Système poulies-courroie]
\begin{itemize}
\item \textbf{Définition :} 2 poules + 1 courroie = transmission rotation entre axes parallèles. Menante (moteur) $\rightarrow$ Menée (utilisateur).
\item \textbf{Sens de rotation :} Non croisée = même sens. Croisée = sens inverses.
\item \textbf{Vitesse :} Grande poulie = lente. Petite poulie = rapide. $N_2 = \frac{D_1}{D_2} N_1$. Rapport $r = D_1/D_2$.
\item \textbf{Adhérence :} Condition sine qua non. Glissement = panne. Fluage = normal (1-2\%).
\item \textbf{Avantages :} Entraze variable, silencieux, amorti chocs, fusible (glissement), pas de graisse, peu cher.
\item \textbf{Inconvénients :} Glissement (fluage), usure caoutchouc, sensible huile/chaleur, encombrement radial, puissance limitée.
\item \textbf{Correction glissement :} Tension (règle du pouce), propreté, courroie trapézoïdale/V, galet tendeur (augmenter enroulement), alignement, courroie dentée si synchrone requis.
\item \textbf{Sécurité :} Carter obligatoire (prise en main), arrêt machine avant intervention, pas de levier forcé.
\item \textbf{Entretien :} Visuel (fissures), tension, alignement, propreté, roulements, galets.
\end{itemize}
\end{aretenir}

\coursec{Avantages du système poulies-courroie}

Dans la section précédente, nous avons vu comment le système poulies-courroie fonctionne : une poulie motrice entraîne une courroie, laquelle entraîne à son tour une poulie réceptrice, avec un sens de rotation identique ou inversé selon que la courroie est ouverte ou croisée. Mais pourquoi ce système est-il si répandu ? Pourquoi le mécanicien de quartier, le constructeur de motos et l'ingénieur de la SONARA l'utilisent-ils tous ? C'est ce que nous allons découvrir maintenant : les \cle{avantages} du système poulies-courroie.

\courssub{Fonctionnement silencieux}

\begin{experience}[Comparaison du bruit : engrenages contre courroie]
\textbf{Matériel :} un vieux réveil mécanique à engrenages (ou une boîte de vitesses de démonstration), et un ventilateur électrique dont le moteur entraîne une petite courroie (ou une machine à coudre).

\textbf{Protocole :}
\begin{enumerate}
\item Approche ton oreille du réveil mécanique en marche et écoute attentivement.
\item Mets ensuite le ventilateur en marche et écoute le bruit produit par sa transmission à courroie.
\item Compare les deux niveaux sonores.
\end{enumerate}

\textbf{Observations :} le réveil à engrenages produit un tic-tac sec et métallique, caractéristique du choc des dents les unes contre les autres. La transmission par courroie, elle, est presque inaudible : on entend à peine un léger frottement.

\textbf{Interprétation :} dans un engrenage, les dents en métal entrent en contact direct et brutal : métal contre métal. Chaque choc produit une vibration sonore. Dans le système poulies-courroie, la courroie est en \cle{caoutchouc} ou en matière plastique souple : elle \emph{absorbe} les vibrations au lieu de les transmettre. Le contact entre la courroie et la poulie est un contact souple, sans choc métallique.

\textbf{Conclusion :} le système poulies-courroie fonctionne de manière \emph{silencieuse}, contrairement aux engrenages.
\end{experience}

\begin{aretenir}[Avantage n°1 : le silence]
La courroie, souple et élastique, \cle{absorbe les vibrations} et évite les chocs métal contre métal. Le système poulies-courroie est donc beaucoup plus \textbf{silencieux} qu'un système d'engrenages ou qu'une chaîne.
\end{aretenir}

\begin{exemplebox}[La courroie de distribution des voitures modernes]
C'est précisément pour son \textbf{silence} que la courroie est préférée à la chaîne dans les moteurs des voitures modernes. La \cle{courroie de distribution} synchronise le mouvement du vilebrequin et celui de l'arbre à cames (qui commande l'ouverture des soupapes). Autrefois, on utilisait une chaîne métallique : le moteur était bruyant. Aujourd'hui, la courroie crantée en caoutchouc armé rend le moteur bien plus discret, ce qui améliore le confort des passagers. Écoute le moteur d'une vieille voiture et celui d'une voiture récente : la différence de bruit vient en grande partie de ce choix technologique.
\end{exemplebox}

\courssub{Transmission possible entre arbres éloignés}

Imagine que tu veuilles transmettre le mouvement du moteur d'un groupe électrogène vers une pompe à eau située à un mètre de distance. Avec des engrenages, il faudrait une longue série de roues dentées intermédiaires : lourd, coûteux et compliqué ! Avec une courroie, le problème disparaît.

\begin{aretenir}[Avantage n°2 : la distance]
Le système poulies-courroie permet de transmettre un mouvement de rotation entre deux arbres \textbf{éloignés l'un de l'autre}. Il suffit d'adapter la \cle{longueur de la courroie} à l'\cle{écartement} (distance entre les axes des deux poulies). Plus les arbres sont éloignés, plus la courroie est longue : c'est aussi simple que cela.
\end{aretenir}

C'est pourquoi on trouve ce système partout où le moteur ne peut pas être collé à la machine entraînée : le moteur de la machine à coudre et son mécanisme de couture, le groupe électrogène et son alternateur, le moteur de la moto et la roue arrière (sur les scooters à transmission par courroie).

\begin{attention}[Distance raisonnable]
La courroie permet de relier des arbres éloignés, mais pas à l'infini ! Si la courroie est trop longue, elle \emph{vibre} et risque de sauter de la poulie. Dans la pratique, on limite l'écartement ou on ajoute un galet tendeur pour guider la courroie.
\end{attention}

\courssub{Simplicité et faible coût}

\begin{aretenir}[Avantage n°3 : simplicité et économie]
Le système poulies-courroie est \textbf{simple} et \textbf{peu coûteux} :
\begin{itemize}
\item les poulies sont des pièces faciles à fabriquer (fonte, aluminium ou même plastique) ;
\item la courroie est un article courant, vendu partout, même dans les marchés de pièces détachées ;
\item le remplacement d'une courroie usée est \textbf{rapide et facile} : on détend, on retire l'ancienne, on pose la nouvelle, on retend. Nul besoin de démonter toute la machine.
\end{itemize}
\end{aretenir}

Compare avec un engrenage : si une dent casse, il faut souvent remplacer la roue entière, usinée avec précision, donc chère. Avec une courroie, le mécanicien de quartier à Yaoundé ou à Douala peut réparer ta machine en quelques minutes avec une courroie neuve achetée à bas prix.

\begin{savaistu}[Des courroies qui durent]
Certaines courroies de motos et de scooters modernes durent plus de \SI{20000}{km} sans aucun entretien ! C'est un avantage considérable par rapport à la chaîne de transmission, qu'il faut \emph{graisser régulièrement} (tous les \SI{500}{km} environ), nettoyer et retendre, sous peine de l'user rapidement. C'est pourquoi de nombreux scooters destinés aux villes africaines, où les pistes sont poussiéreuses, utilisent une transmission finale par courroie : moins d'entretien, moins de pannes, plus de tranquillité pour le moto-taximan.
\end{savaistu}

\courssub{Souplesse de fonctionnement}

\begin{experience}[Démarrage progressif]
Branche un ventilateur dont la transmission se fait par courroie (ou observe une machine à laver au démarrage de l'essorage). Observe attentivement la mise en route.

\textbf{Observation :} la machine ne passe pas brutalement de l'arrêt à la pleine vitesse. Le tambour ou l'hélice accélère \emph{progressivement}.

\textbf{Interprétation :} la courroie, légèrement élastique, s'étire un peu au démarrage et transmet l'effort en douceur. Elle absorbe les \cle{à-coups} (variations brusques de l'effort). Si la machine rencontre une résistance soudaine, la courroie amortit le choc au lieu de le transmettre directement au moteur.
\end{experience}

\begin{aretenir}[Avantage n°4 : la souplesse]
Grâce à l'\cle{élasticité} de la courroie, le système poulies-courroie assure un \textbf{démarrage progressif} et \textbf{absorbe les à-coups} de fonctionnement. Le mouvement transmis est doux et régulier, ce qui ménage à la fois le moteur et la machine entraînée.
\end{aretenir}

\courssub{Sécurité : la courroie, un fusible mécanique}

Voici l'un des avantages les plus remarquables du système, et il peut sembler paradoxal : le \emph{glissement}, que nous étudierons comme un défaut dans la section 4, est ici un \textbf{atout de sécurité}.

Imagine qu'une pierre se coince dans la pompe à eau entraînée par ton groupe électrogène. La pompe se bloque brutalement. Que se passe-t-il ?

\begin{itemize}
\item Avec des engrenages rigides : le moteur continue de forcer, les dents cassent, l'arbre se tord, le moteur grille. La panne est grave et coûteuse.
\item Avec une courroie : la courroie \emph{glisse} sur la poulie motrice, ou finit par \emph{casser}. Le moteur continue de tourner à vide, sans forcer. Il suffit de retirer la pierre et de remettre (ou remplacer) la courroie.
\end{itemize}

\begin{aretenir}[Avantage n°5 : sécurité — fusible mécanique]
En cas de \cle{surcharge} ou de \cle{blocage} de la machine entraînée, la courroie \textbf{glisse} sur les poulies ou \textbf{casse}. Elle interrompt ainsi la transmission avant que le moteur ou la machine ne soient endommagés. On dit que la courroie joue le rôle de \cle{fusible mécanique} : comme le fusible électrique qui fond pour protéger ton installation électrique, la courroie « sacrifie » une pièce peu coûteuse pour protéger les pièces coûteuses du système.
\end{aretenir}

\begin{exemplebox}[Le fusible électrique et le fusible mécanique]
Dans ta maison, quand un court-circuit se produit, le fusible fond et coupe le courant : on remplace un petit fusible à quelques francs au lieu de racheter un téléviseur. De même, quand la machine se bloque, la courroie glisse ou casse : on remplace une courroie bon marché au lieu de réparer un moteur. Le principe de protection est exactement le même, l'un en électricité, l'autre en mécanique.
\end{exemplebox}

\courssub{Grande gamme de vitesses transmissibles}

Enfin, le système poulies-courroie peut transmettre des vitesses de rotation très variées. En choisissant simplement les \textbf{diamètres} des poulies, on obtient le rapport de transmission souhaité (relation théorique, sans glissement) :
\[
\frac{N_{r}}{N_{m}} = \frac{D_{m}}{D_{r}}
\]
où $N_{m}$ et $D_{m}$ sont la vitesse et le diamètre de la poulie motrice, $N_{r}$ et $D_{r}$ ceux de la poulie réceptrice. Une petite poulie motrice et une grande poulie réceptrice donnent une forte réduction de vitesse ; l'inverse donne une multiplication. De plus, la courroie supporte des vitesses de rotation élevées, là où une chaîne devient bruyante et dangereuse.

\begin{aretenir}[Avantage n°6 : la gamme de vitesses]
En jouant sur les diamètres des poulies, on obtient facilement le rapport de vitesse désiré, et la courroie convient à des vitesses de rotation \textbf{élevées}. Le système offre donc une \textbf{grande gamme de vitesses transmissibles}.
\end{aretenir}

\begin{exoresolu}[Choisir les poulies d'un groupe électrogène]
Le moteur d'un groupe électrogène tourne à $N_{m} = \SI{3000}{tr/min}$. L'alternateur qu'il entraîne par courroie doit tourner à $N_{r} = \SI{1500}{tr/min}$. La poulie motrice a un diamètre $D_{m} = \SI{10}{cm}$. Calcule le diamètre $D_{r}$ de la poulie réceptrice, puis cite deux avantages du choix d'une courroie pour cette application.
\tcblower
\textbf{Solution détaillée.}

On applique la relation du rapport de transmission :
\[
\frac{N_{r}}{N_{m}} = \frac{D_{m}}{D_{r}}
\quad\Longrightarrow\quad
D_{r} = D_{m} \times \frac{N_{m}}{N_{r}}
\]
Application numérique :
\begin{align*}
D_{r} &= 10 \times \frac{3000}{1500}\\
&= 10 \times 2\\
&= \SI{20}{cm}
\end{align*}
La poulie réceptrice doit avoir un diamètre de \SI{20}{cm} : comme elle est deux fois plus grande que la poulie motrice, elle tourne deux fois moins vite.

\textbf{Avantages de la courroie ici :} le fonctionnement est \emph{silencieux} (important pour un groupe électrogène utilisé près des habitations) et, en cas de blocage de l'alternateur, la courroie \emph{glisse ou casse}, protégeant le moteur (rôle de fusible mécanique). On peut ajouter la simplicité d'entretien et le faible coût de remplacement.
\end{exoresolu}

\courssub{Bilan : tableau récapitulatif en construction}

Au fur et à mesure de la leçon, nous complétons le tableau comparatif du système poulies-courroie. Voici la colonne des avantages que nous venons d'établir ; la colonne des inconvénients sera remplie dans la section suivante.

\begin{popfigure}
\begin{tikzpicture}[font=\small]
\fill[popgreenL] (0,0) rectangle (4.1,0.8);
\fill[poporangeL] (4.3,0) rectangle (8.4,0.8);
\node at (2.05,0.4) {\textbf{Avantages}};
\node at (6.35,0.4) {\textbf{Inconvénients}};
\draw[popline] (0,0) rectangle (4.1,-5.6);
\draw[popline] (4.3,0) rectangle (8.4,-5.6);
\node[text width=3.7cm, anchor=north west, align=left] at (0.15,-1.0)
{\textbullet\ Fonctionnement silencieux\\[2pt]
\textbullet\ Transmission entre arbres éloignés\\[2pt]
\textbullet\ Simplicité et faible coût\\[2pt]
\textbullet\ Souplesse : démarrage progressif, absorption des à-coups\\[2pt]
\textbullet\ Sécurité : fusible mécanique en cas de surcharge\\[2pt]
\textbullet\ Grande gamme de vitesses};
\node[text width=3.7cm, anchor=north west, align=left, text=popmuted] at (4.45,-1.0)
{\emph{À compléter dans la section suivante...}};
\end{tikzpicture}
\legende{Tableau récapitulatif du système poulies-courroie : les six avantages étudiés dans cette section. La colonne des inconvénients sera complétée à la section 4.}
\end{popfigure}

\begin{attention}[Ne confonds pas avantage et perfection !]
Un piège fréquent au BEPC : croire que le système poulies-courroie n'a que des qualités. Chaque avantage a sa contrepartie : le silence vient de la souplesse, mais cette souplesse provoque aussi du \emph{glissement} ; le glissement protège le moteur, mais il fait perdre de la précision sur la vitesse transmise. Retiens cette idée : en technologie, un même phénomène peut être un avantage dans une situation et un inconvénient dans une autre. C'est exactement ce que nous verrons dans la section 4.
\end{attention}

\begin{exercice}[Pour t'entraîner]
Recopie et complète les phrases suivantes avec les mots : \emph{silencieux}, \emph{fusible mécanique}, \emph{vibrations}, \emph{éloignés}, \emph{coûteux}.
\begin{enumerate}
\item La courroie souple absorbe les \ldots : le système est donc \ldots.
\item La longueur de la courroie s'adapte à l'écartement : on peut relier des arbres \ldots.
\item En cas de blocage, la courroie glisse ou casse : elle joue le rôle de \ldots.
\item Le remplacement d'une courroie est facile et peu \ldots.
\end{enumerate}
\end{exercice}

\begin{corrige}[Exercice : Pour t'entraîner]
\begin{enumerate}
\item La courroie souple absorbe les \textbf{vibrations} : le système est donc \textbf{silencieux}.
\item La longueur de la courroie s'adapte à l'écartement : on peut relier des arbres \textbf{éloignés}.
\item En cas de blocage, la courroie glisse ou casse : elle joue le rôle de \textbf{fusible mécanique}.
\item Le remplacement d'une courroie est facile et peu \textbf{coûteux}.
\end{enumerate}
\end{corrige}

Tu maîtrises maintenant les six grands avantages du système poulies-courroie : silence, distance, économie, souplesse, sécurité et gamme de vitesses. Mais aucun système n'est parfait : dans la section suivante, nous examinerons honnêtement ses inconvénients, à commencer par ce glissement qui, nous l'avons vu, est à la fois un ami et un ennemi.

\coursec{Inconvénients du système poulies-courroie}

Après avoir vu les atouts indéniables du système poulies-courroie — silence de fonctionnement, amortissement des chocs, simplicité d'installation et coût modéré — nous devons maintenant examiner ses limites. En technologie, comme dans la vie courante, il n'existe pas de solution parfaite : chaque choix technique est un compromis. Comprendre les inconvénients, c'est savoir \emph{où} et \emph{quand} ce système reste pertinent, et \emph{comment} le maintenir pour en tirer le meilleur parti. C'est ce que nous allons faire maintenant, en nous appuyant sur des observations concrètes, vécues au quotidien dans nos ateliers, nos maisons et sur nos routes camerounaises.

\courssub{Le glissement : l'ennemi n\textdegree 1 de la transmission exacte}

Commençons par le phénomène le plus connu et le plus redouté : le glissement. Vous avez certainement déjà entendu ce sifflement aigu, ce \og{} chant \fg{} caractéristique d'une courroie qui patine sur une poulie. Cela arrive souvent au démarrage d'un groupe électrogène, quand on accélère brutalement une moto-taxi, ou lorsque le moulin à maïs du quartier force sur une charge trop importante.

\begin{experience}[Observation du glissement sur un banc d'essai ou un moulin à maïs]
\textbf{Matériel :} Un système poulies-courroie monté sur bâti (ou le moulin à maïs du quartier), une charge variable (poids ou frein à poudre), un tachymètre (ou une application smartphone pour mesurer la vitesse de rotation), craie ou marqueur.
\textbf{Protocole :}
\begin{enumerate}
    \item Repérer un point sur la poulie menante (moteur) et un point sur la poulie menée (récepteur) avec de la craie.
    \item Faire tourner le système à vide, mesurer les vitesses de rotation $N_1$ (menante) et $N_2$ (menée). Calculer le rapport réel $N_2/N_1$ et le comparer au rapport théorique $D_1/D_2$.
    \item Augmenter progressivement la charge sur la poulie menée. Noter l'apparition d'un bruit de sifflement.
    \item Mesurer à nouveau $N_2$ sous charge. Observer l'écart avec la valeur théorique.
    \item Relâcher la tension de la courroie (si le montage le permet) et répéter l'essai sous la même charge.
\end{enumerate}
\textbf{Observations :}
\begin{itemize}
    \item À vide, le rapport réel est très proche du rapport théorique (différence $< 1\%$).
    \item Sous charge, la vitesse $N_2$ diminue : la poulie menée tourne moins vite que prévu. On entend un sifflement.
    \item Si la courroie est détendue, le glissement apparaît dès des charges plus faibles et est plus important.
\end{itemize}
\textbf{Interprétation :} Le frottement entre la courroie et les poulies n'est pas suffisant pour transmettre l'effort tangentiel requis par la charge. La courroie \emph{patine} sur la poulie menante (brin tendu qui glisse) et/ou sur la poulie menée (brin mou qui glisse). La transmission n'est plus \emph{rigide} : le rapport de transmission n'est plus constant.
\legende{Expérience mettant en évidence le glissement : la vitesse menée chute sous l'effet de la charge et de la détente.}
\end{experience}

\begin{definition}[Glissement (ou patinage)]
Le \cle{glissement} est le déplacement relatif de la courroie par rapport à la surface de la poulie sur laquelle elle devrait adhérer parfaitement. La poulie menée tourne alors moins vite que la vitesse imposée par le rapport de diamètres ; dans les cas extrêmes, elle peut même s'arrêter alors que la poulie menante continue de tourner.
\end{definition}

\noindent
Le glissement se quantifie par le \cle{taux de glissement} $g$ (en \%), défini par :
\[
g = \frac{N_{2,\text{théorique}} - N_{2,\text{réel}}}{N_{2,\text{théorique}}} \times 100
\]
où $N_{2,\text{théorique}} = N_1 \times \frac{D_1}{D_2}$.
Un glissement de \SIrange{1}{2}{\percent} est souvent toléré en fonctionnement normal. Au-delà de \SIrange{3}{4}{\percent}, la transmission devient inefficace : la courroie s'échauffe, s'use prématurément et la précision est perdue.

\begin{attention}[Piège fréquent : la détente, fausse amie]
Ne croyez \textbf{jamais} qu'une courroie un peu \og{} molle \fg{} ou détendue transmet mieux parce qu'elle \og{} accroche \fg{} plus. C'est \textbf{l'inverse} : la détente est la \textbf{première cause de glissement}. Une courroie détendue réduit l'angle d'enroulement et la pression de contact sur les poulies, diminuant drastiquement la force de frottement maximale transmissible ($F_{\text{frott}} = \mu \cdot N$). \textbf{Toujours tendre correctement} selon les préconisations du constructeur (flèche de \SIrange{10}{15}{mm} par mètre d'entraxe pour une courroie trapézoïdale classique).
\end{attention}

\courssub{Usure et vieillissement : une maintenance incontournable}

Une courroie n'est pas éternelle. C'est un élément \emph{consommable}, fait d'élastomères (caoutchouc naturel ou synthétique) renforcés de fibres textiles (coton, polyester, aramide) ou métalliques. Au Cameroun, avec nos variations de température (chaleur intense en saison sèche, humidité en saison des pluies) et la poussière rouge latéritique omniprésente, le vieillissement est accéléré.

\begin{itemize}
    \item \textbf{Détente (allongement permanent) :} Sous l'effet de la tension continue et des cycles de charge, les fibres de renfort s'allongent irréversiblement. La courroie \og{} fait le ventre \fg{}, le brin mou s'affaisse. C'est le premier signe : il faut retendre (voir Section 6).
    \item \textbf{Fissuration (craquelures) :} L'ozone, les UV du soleil, la chaleur du moteur et les huiles attaquent le caoutchouc. Des micro-fissures apparaissent sur les flancs et le dos de la courroie, perpendiculairement au sens de roulement. Elles s'approfondissent jusqu'à la rupture.
    \item \textbf{Effilochage / Délaminage :} Les bords s'effilochent (mauvais alignement des poulies) ou les plis se décollent (contamination par l'huile/graisse).
    \item \textbf{Rupture :} C'est la fin brutale. La machine s'arrête net. Sur un moteur de voiture (courroie de distribution), c'est souvent la casse moteur (pistons contre soupapes). Sur un groupe électrogène, c'est la panne de courant.
\end{itemize}

\begin{popfigure}
\begin{tikzpicture}[scale=0.9, every node/.style={font=\small}]
    % Courroie correctement tendue
    \begin{scope}[xshift=-7cm]
        \draw[popblue, thick] (0,0) circle (1.5);
        \draw[popblue, thick] (5,0) circle (1);
        % Brin tendu supérieur
        \draw[poporange, very thick] (1.5,0) -- (4,0);
        % Brin tendu inférieur
        \draw[poporange, very thick] (1.5,0) -- (4,0);
        % Flèches rotation
        \draw[->, popdark, thick] (0,1.8) arc (90:450:1.8);
        \draw[->, popdark, thick] (5,1.3) arc (90:450:1.3);
        \node[popblue] at (0,-2.2) {\textbf{Correctement tendue}};
        \node[poporange] at (2.5,0.4) {Brins tendus};
        \draw[<->, popmuted] (1.5,-1.5) -- (4,-1.5) node[midway,below] {Entraxe $E$};
    \end{scope}
    
    % Courroie détendue (brin mou affaissé)
    \begin{scope}[xshift=5cm]
        \draw[popblue, thick] (0,0) circle (1.5);
        \draw[popblue, thick] (5,0) circle (1);
        % Brin tendu supérieur (tendu)
        \draw[poporange, very thick] (1.5,0) -- (4,0);
        % Brin MOU inférieur (affaissé - courbe de chaînette simplifiée)
        \draw[poppink, very thick] (1.5,0) .. controls (2.5,-0.8) and (4,-0.8) .. (4,0);
        % Flèches rotation (menée plus lente ou arrêtée)
        \draw[->, popdark, thick] (0,1.8) arc (90:450:1.8);
        \draw[->, poppink, thick, dashed] (5,1.3) arc (90:200:1.3);
        \node[popblue] at (0,-2.2) {\textbf{Détendue (glissement)}};
        \node[poporange] at (2.5,0.4) {Brin tendu};
        \node[poppink] at (2.5,-1.1) {Brin \textbf{mou} (affaissé)};
        \draw[<->, popmuted] (1.5,-1.5) -- (4,-1.5) node[midway,below] {Entraxe $E$};
        \node[align=center, popdark, font=\tiny] at (5, -2.2) {Angle d'enroulement\\réduit sur poulie menée};
    \end{scope}
\end{tikzpicture}
\legende{Comparaison : courroie correctement tendue (brins rectilignes, bon enroulement) vs courroie détendue (brin mou affaissé, angle d'enroulement réduit sur la poulie menée, glissement accru).}
\end{popfigure}

\begin{exemplebox}[Le moulin à maïs du quartier]
Maman \emph{Ngo} gère un moulin à maïs à \emph{Mvog-Ada} (Yaoundé). Un matin, la courroie du moteur Diesel \emph{Lister} (poulie menante $\varnothing \SI{150}{mm}$) vers la meule (poulie menée $\varnothing \SI{400}{mm}$) se met à siffler. La meule ralentit, le maïs n'est plus broyé finement. Maman \emph{Ngo} arrête le moteur, regarde : la courroie a \og{} fait le ventre \fg{}, le brin inférieur touche presque le châssis. Elle a 3 ans, des fissures sur les flancs. \textbf{Diagnostic :} Usure + détente $\rightarrow$ glissement. \textbf{Solution :} Remplacement immédiat (courroie neuve type SPA ou SPB selon gorge) + alignement poulies + tension correcte. Coût : \SIrange{3500}{5000}{FCFA}. Économie : évite la casse du moteur ou une journée de travail perdue.
\end{exemplebox}

\courssub{Rendement inférieur aux engrenages : des pertes énergétiques réelles}

Le rendement $\eta$ d'une transmission est le rapport de la puissance utile (sortie) sur la puissance fournie (entrée) : $\eta = P_{\text{sortie}} / P_{\text{entrée}}$. Pour un système poulies-courroie bien entretenu, $\eta$ se situe généralement entre \textbf{\SIrange{92}{97}{\percent}}. C'est bon, mais inférieur aux engrenages bien lubrifiés (\SIrange{98}{99}{\percent}) ou aux chaînes (\SIrange{97}{99}{\percent}).

Où part l'énergie perdue (les \SIrange{3}{8}{\percent}) ?
\begin{enumerate}
    \item \textbf{Frottement interne dans la courroie (hystérésis) :} La courroie se déforme en entrant dans la poulie (flexion) et reprend sa forme en sortant. Ce cycle de déformation/réforme à chaque tour dissipe de l'énergie sous forme de chaleur dans le caoutchouc. C'est la \textbf{perte principale} (surtout pour les grosses sections et vitesses élevées).
    \item \textbf{Frottement de glissement :} Même sans glissement macroscopique, il y a un \emph{microslippage} permanent (différence de vitesse infinitésimale) dû à l'élasticité de la courroie (glissement élastique). En cas de glissement macroscopique, les pertes deviennent énormes (échauffement violent).
    \item \textbf{Frottement dans les roulements :} La tension de la courroie crée une charge radiale sur les arbres des poulies, augmentant le frottement dans les paliers/roulements.
    \item \textbf{Ventilation (pertes aérodynamiques) :} À haute vitesse périphérique ($> \SIrange{25}{30}{m/s}$), la courroie brasse l'air (effet ventilateur).
\end{enumerate}

\begin{propriete}[Rendement typique selon le type de courroie (valeurs indicatives, bien tendues, charge nominale)]
\begin{center}
\begin{tabularx}{\linewidth}{|l|X|X|}\hline
\textbf{Type de courroie} & \textbf{Rendement $\eta$ typique} & \textbf{Commentaire} \\ \hline
Trapézoïdale classique (Z, A, B, C, D, E) & \SIrange{93}{96}{\percent} & Standard, bon compromis coût/perf. \\ \hline
Trapézoïdale étroite (SPZ, SPA, SPB, SPC) & \SIrange{95}{97}{\percent} & Meilleure flexion, plus de puissance. \\ \hline
Dentée (synchrone) & \SIrange{98}{99}{\percent} & Pas de glissement (dents), mais bruit + coût. \\ \hline
Plate (cuir, caoutchouc, synthétique) & \SIrange{90}{95}{\percent} & Ancien, gros diamètres, glissement élevé. \\ \hline
\end{tabularx}
\end{center}
\legende{Le choix du type de courroie influence directement le rendement. Pour une moto-taxi (variateur CVT), on utilise une courroie trapézoïdale étroite renforcée aramide pour limiter les pertes.}
\end{propriete}

\courssub{Sensibilité aux conditions extérieures : chaleur, huile, poussière}

La courroie est un organe \og{} à nu \fg{}, souvent exposé. Son matériau (élastomère) est sensible à son environnement.

\begin{itemize}
    \item \textbf{Chaleur excessive ($> \SIrange{60}{70}{\celsius}$ en continu) :} Durcit le caoutchouc (perte d'élasticité $\rightarrow$ fissures), accélère la détente, réduit le coefficient de frottement $\mu$ $\rightarrow$ glissement. \emph{Exemple :} Courroie d'alternateur ou de pompe à eau juste derrière le radiateur d'un camion chaud, ou courroie de compresseur de climatisation.
    \item \textbf{Huile, graisse, hydrocarbures (essence, gasoil) :} Le caoutchouc gonfle, se ramollit, perd son adhérence ($\mu$ chute drastiquement). Glissement immédiat et irréversible. \emph{Exemple :} Fuite de joint de vilebrequin sur courroie de distribution (catastrophe) ; huile de chaîne de moto projetée sur courroie de variateur ; graissage maladroit des poulies (\textbf{interdit !}).
    \item \textbf{Poussière abrasive (sable, latérite) :} Pénètre dans les gorges, use prématurément les flancs de la courroie \textbf{et} les gorges des poulies (les poulies s'usent aussi !). Une poulie usée (gorge évasée) détruit une courroie neuve en quelques heures.
    \item \textbf{Produits chimiques (acides, bases, solvants) :} Attaque chimique du caoutchouc.
\end{itemize}

\begin{attention}[Erreur de maintenance : graisser la courroie !]
\textbf{JAMAIS} on ne met d'huile, de graisse, de savon, de \emph{WD-40} ou de \og{} produit miracle \fg{} sur une courroie qui siffle pour \og{} la faire accrocher \fg{}. Cela détruit l'adhérence ($\mu \approx 0$) et le caoutchouc. Si elle siffle : \textbf{1. Arrêter. 2. Nettoyer poulies et courroie (chiffon sec, dégraissant volatil si huile, puis sécher). 3. Vérifier usure poulies. 4. Retendre ou remplacer.}
\end{attention}

\courssub{Encombrement et puissance limitée : les contraintes géométriques et mécaniques}

Deux derniers inconvénients majeurs conditionnent le choix de cette technologie face aux engrenages ou aux chaînes.

\paragraph{1. Encombrement (entraxe et diamètres).}
Pour transmettre une puissance donnée sans glissement, il faut un \emph{angle d'enroulement} suffisant sur la petite poulie (généralement $> 120^\circ$, idéalement $> 150^\circ$). Cela impose un \textbf{entraxe minimum} entre les arbres. De plus, pour des rapports de réduction importants (ex: $1/5$, $1/10$), la grande poulie devient \textbf{très encombrante} (diamètre $5$ à $10$ fois celui de la petite).
\begin{itemize}
    \item \emph{Exemple :} Sur un compresseur d'atelier entraîné par moteur électrique \SI{3000}{tr/min} $\rightarrow$ compresseur \SI{600}{tr/min} (rapport $1/5$). Si poulie menante $\varnothing \SI{80}{mm}$, poulie menée $\varnothing \SI{400}{mm}$. Cela prend de la place sur le bâti.
    \item \emph{Contre-exemple :} Un réducteur à engrenages (train d'engrenages) fait le même rapport dans un boîtier compact, lubrifié, sans tension à surveiller.
\end{itemize}

\paragraph{2. Puissance transmissible limitée.}
La puissance transmissible $P$ (en kW) dépend de la vitesse, des diamètres, de l'angle d'enroulement, du coefficient de frottement et de la section de la courroie. Elle est limitée par :
\begin{itemize}
    \item La résistance à la traction des fibres (rupture).
    \item La résistance au frottement (glissement).
    \item La résistance à la fatigue par flexion (durée de vie).
\end{itemize}
Pour de \textbf{grosses puissances} ($> \SIrange{50}{100}{kW}$ par courroie, soit $\sim \SIrange{70}{130}{ch}$), on doit utiliser \textbf{plusieurs courroies en parallèle} (jeux de $2$, $3$, $4$, $5$ courroies \emph{appariées} / \emph{matched set}) sur des poulies à gorges multiples. Cela devient lourd, large, cher à maintenir (il faut changer tout le jeu si l'une casse). Au-delà, les \textbf{engrenages} (boîtes de vitesse camions, réducteurs industriels) ou les \textbf{chaînes} (transmissions primaires moto, vélos, convoyeurs lourds) s'imposent.

\begin{exoresolu}[Choix technique : pompe d'irrigation motorisée]
\textbf{Énoncé :} Un agriculteur à \emph{Yagoua} (Extrême-Nord) veut motoriser sa pompe centrifuge (puissance absorbée $P = \SI{7,5}{kW} \approx \SI{10}{ch}$, vitesse pompe $N_2 = \SI{1500}{tr/min}$) avec un moteur Diesel $N_1 = \SI{3000}{tr/min}$. Entraxe disponible $\approx \SI{600}{mm}$. Environnement : poussiéreux, chaud (\SI{40}{\celsius}). Proposer une solution poulies-courroie adaptée et justifier.
\textbf{Solution détaillée :}
\begin{enumerate}
    \item \textbf{Rapport de transmission :} $i = N_1/N_2 = 3000/1500 = 2$. Soit $D_2/D_1 = 2$.
    \item \textbf{Choix des diamètres (standardisation) :} Pour entraxe $E = \SI{600}{mm}$ et angle d'enroulement correct sur petite poulie, choisir $D_1 = \SI{160}{mm}$ (menante), $D_2 = \SI{315}{mm}$ (menée). Vérification angle : $\sin \alpha \approx (D_2-D_1)/(2E) = 155/1200 \approx 0,13 \Rightarrow \alpha \approx 7,5^\circ$. Angle sur petite poulie $\approx 180^\circ - 2\alpha \approx 165^\circ$ (Excellent $> 150^\circ$).
    \item \textbf{Type de courroie :} Puissance \SI{7,5}{kW}, vitesse périphérique $v = \pi D_1 N_1 / 60\,000 \approx 3,14 \times 160 \times 3000 / 60000 \approx \SI{25}{m/s}$. Limite courroie trapézoïdale classique ($\sim \SIrange{25}{30}{m/s}$). Choisir \textbf{courroie trapézoïdale étroite section SPA} (ou SPB pour marge). Résiste mieux à la chaleur et à la flexion.
    \item \textbf{Nombre de courroies :} Puissance par courroie SPA $\varnothing \SI{160}{mm}$, $N_1=\SI{3000}{tr/min}$, $i=2$ $\approx \SI{4,5}{kW}$ (tableaux constructeur). Il faut \textbf{2 courroies SPA en parallèle} (puissance $\approx \SI{9}{kW}$ avec coefficient de charge). Poulie menante à 2 gorges SPA, poulie menée à 2 gorges SPA.
    \item \textbf{Protection :} Carter grillagé (laisse passer l'air, retient pierres/poussière gros grains). Tendeur à ressort (compense dilatation thermique + détente).
    \item \textbf{Alternative chaîne :} Possible (puissance OK, moins sensible à la poussière si graissée), mais plus bruyante, lubrification contraignante en zone poussiéreuse (pâte abrasive), pas d'amortissement chocs (pompe = charge régulière). \textbf{Courroies SPA x2 retenues.}
\end{enumerate}
\legende{Le dimensionnement réel utilise les catalogues constructeurs (tableaux puissance/vitesse). Ici, 2 courroies SPA assurent la fiabilité dans des conditions difficiles.}
\end{exoresolu}

\begin{aretenir}[Bilan des inconvénients du système poulies-courroie]
\begin{itemize}
    \item \textbf{Glissement :} Perte de rapport de transmission, échauffement, usure. Cause principale : détente, surcharge, huile, poulies usées.
    \item \textbf{Usure / Maintenance :} Élément consommable (caoutchouc). Fissures, détente, rupture. Remplacement périodique obligatoire (préventif).
    \item \textbf{Rendement :} \SIrange{92}{97}{\percent} (pertes par hystérésis, microslippage, ventilation). Inférieur aux engrenages/chaînes.
    \item \textbf{Sensibilité environnement :} Chaleur, huile/graisse (interdits !), poussière abrasive, ozone/UV.
    \item \textbf{Encombrement :} Grand entraxe nécessaire, grosses poulies pour forts rapports.
    \item \textbf{Puissance limitée :} Nécessite jeux de courroies multiples au-delà $\sim \SIrange{50}{100}{kW}$.
\end{itemize}
\emph{Ces inconvénients ne sont pas des défauts rédhibitoires, mais des \textbf{contraintes de conception et d'exploitation} à maîtriser.}
\end{aretenir}

\courssub{Synthèse : quand choisir (ou pas) le système poulies-courroie ?}

\begin{center}
\begin{tabularx}{\linewidth}{|l|X|X|}\hline
\textbf{Critère} & \textbf{Favorable aux poulies-courroie} & \textbf{Défavorable (préférer engrenages/chaînes)} \\ \hline
\textbf{Précision de rapport} & Non critique (ventilateurs, pompes, alternateurs, compresseurs, machines-outils légères) & Critique (distribution moteur, horlogerie, robotique, fraiseuses CN) \\ \hline
\textbf{Entraxe} & Grand ($> 2 \times D_{\text{grand}}$) ou variable & Très court (logement compact) \\ \hline
\textbf{Amortissement chocs/vibrations} & Requis (moteur thermique $\rightarrow$ machine) & Indésirable (rigidité maximale) \\ \hline
\textbf{Environnement} & Propre, tempéré, sec & Très huileux, très chaud ($> \SI{80}{\celsius}$), très abrasif non protégé \\ \hline
\textbf{Puissance} & Moyenne ($< \SI{50}{kW}$ par brin) & Très élevée ($> \SI{100}{kW}$) ou couple de démarrage énorme \\ \hline
\textbf{Maintenance} & Accessible, surveillance visuelle facile & Inaccessible, \og{} fit and forget \fg{} souhaité \\ \hline
\end{tabularx}
\end{center}

\noindent
Dans la pratique camerounaise, le système poulies-courroie reste \textbf{roi} pour : les groupes électrogènes (moteur $\rightarrow$ alternateur), les pompes à eau (moteur thermique ou électrique $\rightarrow$ pompe), les compresseurs d'atelier, les moulins à maïs/manioc, les scies à bois, les ventilateurs de séchoirs, les alternateurs de voitures/camions, les variateurs de moto-taxi (CVT). C'est une technologie \textbf{robuste, réparable localement, peu coûteuse}, à condition de respecter les règles de tension, d'alignement et de propreté que nous verrons en Section 6.

\begin{methode}[Diagnostiquer un problème sur un système poulies-courroie en 5 étapes]
\begin{enumerate}
    \item \textbf{ÉCOUTER :} Sifflement = glissement (détente, huile, charge, usure). Claquement régulier = jointure courroie ou gorge abîmée. Bruit roulement = palier poulie/moteur.
    \item \textbf{REGARDER (arrêté) :} État courroie (fissures, effilochage, brillance/glacage = glissement ancien, huile). État gorges poulies (usure, rouille, huile). Alignement (règle/laser sur flancs poulies). Tension (flèche brin mou).
    \item \textbf{TOUCHER (arrêté, moteur froid, clé sur contact coupé) :} Tension manuelle (dureté brins). Jeu radial poulies (roulements). Température poulies après marche (surchauffe = glissement/frottement).
    \item \textbf{MESURER (si possible) :} Vitesse menante $N_1$ et menée $N_2$ (tachymètre). Calculer glissement $g$. Mesurer entraxe.
    \item \textbf{AGIR :} Nettoyer (sec) $\rightarrow$ Retendre (norme) $\rightarrow$ Realigner $\rightarrow$ Remplacer (courroie + vérifier poulies) $\rightarrow$ Protéger (carter, supprimer fuite huile).
\end{enumerate}
\end{methode}

Nous avons maintenant dressé le tableau complet des forces et faiblesses du système poulies-courroie. La section suivante, \og{} Méthodes de correction des glissements \fg{}, nous donnera les clés pratiques pour maintenir ce système au top de sa performance, jour après jour, sous nos latitudes.

\coursec{Méthodes de correction des glissements}

Nous venons de voir, dans la section précédente, que le système poulies-courroie possède un défaut majeur : la courroie peut \cle{glisser} sur les poulies, ce qui diminue la puissance transmise, use la courroie et fait perdre du rendement. La bonne nouvelle, c'est que ce glissement n'est pas une fatalité ! Les techniciens disposent de plusieurs méthodes simples et efficaces pour le corriger. C'est l'objet de cette section : comprendre \emph{pourquoi} la courroie glisse, puis apprendre \emph{comment} y remédier, comme le fait chaque jour le mécanicien de moto-taxi ou le réparateur de groupe électrogène.

\courssub{Rappel : les causes du glissement}

Avant de soigner une maladie, il faut en connaître la cause. De même, avant de corriger un glissement, il faut identifier son origine. Une courroie transmet le mouvement par \cle{adhérence} : c'est le frottement entre la courroie et la gorge de la poulie qui « accroche » et entraîne la poulie réceptrice. Si cette adhérence devient insuffisante, la courroie patine.

\begin{experience}[Observation : pourquoi la courroie du groupe électrogène patine-t-elle ?]
\textbf{Matériel :} un groupe électrogène dont la courroie glisse (ou la vidéo d'un tel cas), une clé, un chiffon.

\textbf{Protocole :}
\begin{enumerate}
\item Faire tourner le moteur et observer : la poulie motrice tourne vite, mais la poulie réceptrice tourne lentement ou s'arrête, avec un sifflement aigu.
\item Arrêter le moteur, puis appuyer avec le pouce au milieu du brin de courroie : elle s'enfonce de plusieurs centimètres.
\item Toucher la courroie et les gorges des poulies : on sent une pellicule d'huile et de graisse.
\item Examiner les flancs de la courroie : ils sont lisses, brillants, craquelés.
\end{enumerate}

\textbf{Observations :} la courroie est détendue, grasse et usée.

\textbf{Interprétation :} la tension insuffisante réduit la pression de la courroie contre la poulie ; la graisse supprime le frottement ; l'usure a poli les surfaces de contact. Dans les trois cas, l'adhérence devient trop faible : la courroie glisse.

\textbf{Conclusion :} le glissement a quatre causes principales.
\end{experience}

\begin{aretenir}[Les quatre causes principales du glissement]
\begin{enumerate}
\item \textbf{Tension insuffisante} de la courroie (elle est trop lâche) ;
\item \textbf{Usure} de la courroie (flancs polis, craquelures, allongement) ;
\item \textbf{Présence de graisse ou d'huile} sur la courroie ou les poulies ;
\item \textbf{Surcharge} : l'effort demandé à la courroie dépasse ce que l'adhérence peut transmettre.
\end{enumerate}
À chaque cause correspond une méthode de correction. Étudions-les une à une.
\end{aretenir}

\courssub{Méthode 1 : retendre la courroie en écartant les axes des poulies}

\textbf{Intuition.} Imagine un élastique tendu entre tes deux mains. Si tu rapproches les mains, l'élastique devient mou et pend ; si tu écartes les mains, il se tend. La courroie se comporte exactement pareil : pour la tendre, il suffit d'\emph{écarter les deux poulies}.

\textbf{Principe.} Sur la plupart des machines (groupe électrogène, pompe à eau, moulin à maïs), le moteur est fixé sur des \cle{rails de fixation} percés de trous oblongs. Pour retendre la courroie, on desserre les boulons de fixation du moteur, on fait glisser le moteur pour l'écarter de la poulie réceptrice, puis on resserre les boulons. La distance entre les axes des poulies (l'\emph{entraxe}) augmente, et la courroie se tend.

\begin{methode}[Retendre une courroie en déplaçant le moteur]
\begin{enumerate}
\item Arrêter la machine et couper l'alimentation (sécurité d'abord !).
\item Desserrer légèrement les boulons de fixation du moteur sur ses rails.
\item Écarter le moteur (à la main ou avec un levier) jusqu'à obtenir la bonne tension.
\item Vérifier la tension : en appuyant avec le pouce au milieu du brin le plus long, la courroie doit s'enfoncer d'environ \SI{1}{cm} à \SI{1,5}{cm}.
\item Resserrer fermement les boulons et refaire un essai de fonctionnement.
\end{enumerate}
\end{methode}

\begin{attention}[Ni trop lâche, ni trop tendue !]
Une courroie trop lâche glisse ; mais une courroie \textbf{trop tendue} use prématurément les roulements des arbres et peut casser. La bonne tension est un juste milieu : environ \SI{1}{cm} d'enfoncement sous le pouce. Piège classique au BEPC : ne jamais écrire qu'il faut tendre la courroie « le plus possible ».
\end{attention}

\courssub{Méthode 2 : utiliser un galet tendeur}

Sur certaines machines, on ne peut pas déplacer le moteur, ou la tension se relâche régulièrement avec l'usure. On utilise alors un dispositif astucieux : le \cle{galet tendeur}.

\begin{definition}[Galet tendeur]
Un \textbf{galet tendeur} est un rouleau monté sur un bras articulé et rappelé par un \textbf{ressort}, qui appuie en permanence sur le \textbf{brin mou} de la courroie pour en maintenir la tension automatiquement.
\end{definition}

Le galet presse le \emph{brin mou} (le brin qui n'est pas tiré), ce qui tend l'autre brin. Grâce au ressort, la pression se règle toute seule : même si la courroie s'allonge légèrement avec le temps, le galet la suit et la tension reste correcte. C'est une correction \emph{automatique} et \emph{permanente}.

\begin{popfigure}
\begin{tikzpicture}[scale=1.0, every node/.style={font=\small}]
% Poulie motrice (grande)
\draw[thick, fill=popblueL] (0,0) circle (1.5);
\draw[fill=popdark] (0,0) circle (0.08);
\node at (0,-2.0) {Poulie motrice};
% Poulie réceptrice (petite)
\draw[thick, fill=poporangeL] (5,0) circle (0.9);
\draw[fill=popdark] (5,0) circle (0.08);
\node at (5,-1.6) {Poulie r\'eceptrice};
% Courroie : brin tendu (haut)
\draw[very thick, popdark] (0,1.5) -- (5,0.9);
% Brin mou (bas) avec enroulement sur galet tendeur
\draw[very thick, popdark] (0,-1.5) -- (3.45,-1.5);
% Contact galet (tangence bas)
\draw[very thick, popdark] (3.45,-1.5) -- (4.0,-2.05);
\draw[very thick, popdark] (4.0,-2.65) -- (4.55,-1.5);
\draw[very thick, popdark] (4.55,-1.5) -- (5,-0.9);
% Galet tendeur
\draw[thick, fill=popgreenL] (4.0,-2.1) circle (0.55);
\draw[fill=popdark] (4.0,-2.1) circle (0.06);
% Ressort
\draw[thick, popgreen, decorate, decoration={coil, aspect=0.5, segment length=3pt, amplitude=3pt}] (4.0,-2.65) -- (4.0,-3.3);
\draw[thick, popgreen] (3.6,-3.3) -- (4.4,-3.3);
\node[popgreen, align=center] at (5.5,-2.9) {Ressort de\\rappel};
% Flèches sens rotation
\draw[->, thick, popblue] (0.9,1.1) arc (50:120:1.1);
\draw[->, thick, poporange] (5.55,0.65) arc (50:120:0.65);
% Légendes brins
\node[align=center] at (1.5,1.8) {Brin tendu};
\node[align=center] at (1.5,-2.5) {Brin mou\\(poussé par galet)};
\end{tikzpicture}
\legende{Schéma d'un système poulies-courroie équipé d'un galet tendeur : le rouleau vert, poussé par un ressort, appuie sur le brin mou de la courroie et maintient sa tension en permanence.}
\end{popfigure}

\begin{exemplebox}[Le tendeur automatique des voitures]
Dans les voitures, la courroie d'accessoires entraîne l'\textbf{alternateur} (qui produit l'électricité), la pompe à eau et parfois le compresseur de \textbf{climatisation}. Un \textbf{tendeur automatique} (galet monté sur ressort) maintient la tension de cette courroie en permanence, sans aucun réglage. Quand ce tendeur est fatigué, on entend un sifflement au démarrage : c'est le signe que la courroie glisse !
\end{exemplebox}

\courssub{Méthode 3 : remplacer la courroie usée}

Une courroie qui a beaucoup travaillé s'\emph{allonge}, ses flancs deviennent \emph{lisses et brillants}, des craquelures apparaissent. Même bien tendue, elle n'adhère plus correctement. La seule solution est alors de la \textbf{remplacer par une courroie neuve de bonne dimension}.

\begin{attention}[Choisir la bonne courroie de remplacement]
La courroie neuve doit avoir :
\begin{itemize}
\item le même \textbf{profil} (section trapézoïdale, plate ou ronde) ;
\item la même \textbf{longueur} (ou la longueur prévue par le constructeur) ;
\item la même \textbf{largeur}.
\end{itemize}
Une courroie trop longue sera impossible à tendre ; une courroie trop courte sera trop tendue et cassera ou usera les roulements. Piège fréquent : croire qu'on peut réparer durablement une courroie usée en la retendant — cela ne fait que retarder la panne.
\end{attention}

\courssub{Méthode 4 : nettoyer les poulies et la courroie}

L'huile et la graisse sont les ennemies de l'adhérence : elles forment un film glissant qui supprime le frottement entre la courroie et la gorge de la poulie. Sur un moteur, une fuite d'huile peut facilement souiller la courroie.

\begin{methode}[Nettoyer une transmission souillée]
\begin{enumerate}
\item Arrêter la machine et couper l'alimentation.
\item Nettoyer les gorges des poulies avec un chiffon propre et un peu d'essence ou de dégraissant.
\item Nettoyer la courroie ; si elle est imprégnée d'huile en profondeur, la remplacer (le caoutchouc gorgé d'huile ne sèche jamais vraiment).
\item Réparer la fuite d'huile à l'origine de la souillure, sinon le problème reviendra.
\end{enumerate}
\end{methode}

\begin{attention}[Ne jamais graisser une courroie !]
Contrairement à une chaîne de bicyclette, une courroie ne doit \textbf{jamais} être graissée ni huilée. C'est une erreur grave et très fréquente : la graisse détruit l'adhérence, qui est pourtant le principe même de la transmission par courroie.
\end{attention}

\courssub{Méthode 5 : employer une courroie crantée (synchrone)}

Quand le glissement est absolument interdit — par exemple pour entraîner l'arbre à cames d'un moteur, qui doit tourner \emph{exactement} en phase avec le vilebrequin — on change de technologie : on utilise une \cle{courroie crantée}.

\textbf{Principe.} La face intérieure de la courroie porte des \textbf{dents} (des crans) qui s'engagent dans les \textbf{rainures} d'une poulie dentée. La transmission ne se fait plus par adhérence mais par \emph{engrènement}, comme entre deux engrenages : les dents de la courroie poussent mécaniquement les dents de la poulie. Le glissement devient \textbf{impossible}.

\begin{popfigure}
\begin{tikzpicture}[scale=1.1, every node/.style={font=\small}]
% Poulie dentée (vue partielle)
\draw[thick, fill=popblueL] (0,0) circle (1.6);
\draw[fill=popdark] (0,0) circle (0.08);
% Dents de la poulie sur le haut (arc 30 à 150)
\foreach \a in {30,50,70,90,110,130,150}{
  \draw[thick, fill=popblue] (\a:1.6) -- (\a:1.9) -- (\a+7:1.9) -- (\a+7:1.6) -- cycle;
}
% Courroie crantée au-dessus (brin horizontal)
\draw[very thick, popdark] (-2.6,2.15) -- (2.6,2.15); % fond
\draw[very thick, popdark] (-2.6,2.55) -- (2.6,2.55); % haut
% Crans de la courroie (dents vers le bas)
\foreach \x in {-2.2,-1.5,-0.8,-0.1,0.6,1.3,2.0}{
  \draw[thick, fill=poporange] (\x,2.15) rectangle (\x+0.35,1.82);
}
\node[poporange] at (3.6,2.35) {Courroie crant\'ee};
\node[popblue] at (3.3,1.1) {Poulie dent\'ee};
\draw[->, popmuted] (3.3,1.35) -- (1.2,1.75);
% Engagement (flèche double)
\draw[<->, popgreen, thick] (0.28,1.9) -- (0.28,2.12);
\node[popgreen] at (-1.7,3.0) {Les dents s'engagent};
\draw[->, popgreen] (-1.7,2.8) -- (0.1,2.3);
% Sens de rotation
\draw[->, thick, popblue] (1.0,1.15) arc (45:120:1.15);
\end{tikzpicture}
\legende{Vue partielle d'une courroie crantée engagée sur une poulie dentée : les crans de la courroie pénètrent dans les rainures de la poulie, ce qui rend le glissement impossible.}
\end{popfigure}

\begin{savaistu}[Pourquoi dit-on « courroie synchrone » ?]
La courroie crantée est aussi appelée \textbf{courroie synchrone}, car les deux poulies tournent parfaitement \emph{en synchronisme}, sans aucun glissement : le rapport de transmission est exact et constant. C'est pour cela qu'on l'utilise pour la \textbf{distribution} des moteurs de voiture (la fameuse « courroie de distribution ») : les soupapes doivent s'ouvrir au bon moment, au millième de tour près ! On la retrouve aussi dans les imprimantes et les machines-outils de précision.
\end{savaistu}

\courssub{Méthode 6 : augmenter l'arc d'enroulement}

\textbf{Intuition.} Plus la courroie « embrasse » la poulie sur une grande portion de sa circonférence, plus la surface de contact est grande, et meilleure est l'adhérence. C'est comme tenir une corde : si tu l'enroules d'un demi-tour autour d'un poteau, elle tient mieux que si elle ne fait que le toucher.

\textbf{Principe.} L'\cle{arc d'enroulement} est la portion de circonférence de la poulie en contact avec la courroie. C'est surtout sur la \textbf{petite poulie} que cet arc est faible, donc que le glissement commence. Pour l'augmenter, on place un \textbf{galet presseur} qui pousse la courroie contre la petite poulie : la courroie enveloppe alors une plus grande partie de la poulie, et l'adhérence s'améliore.

\begin{aretenir}[Bilan des six méthodes de correction du glissement]
\begin{center}
\begin{tabularx}{\linewidth}{|l|X|}
\hline
\rowcolor{popblueL} \textbf{Cause du glissement} & \textbf{Méthode de correction} \\
\hline
Courroie détendue & Retendre en écartant les axes des poulies (moteur sur rails) \\
\hline
Tension qui se relâche souvent & Poser un galet tendeur (tension automatique) \\
\hline
Courroie usée, allongée, craquelée & Remplacer par une courroie neuve de bonne dimension \\
\hline
Huile ou graisse sur la transmission & Nettoyer poulies et courroie, réparer la fuite \\
\hline
Glissement absolument interdit & Utiliser une courroie crantée (synchrone) \\
\hline
Adhérence faible sur la petite poulie & Augmenter l'arc d'enroulement (galet presseur) \\
\hline
\end{tabularx}
\end{center}
\end{aretenir}

\begin{methode}[Démarche complète pour corriger un glissement]
Face à une courroie qui glisse, le bon technicien procède toujours dans l'ordre :
\begin{enumerate}
\item \textbf{Vérifier la tension} de la courroie (enfoncement au pouce) ;
\item \textbf{Nettoyer} les poulies et la courroie si elles sont grasses ;
\item \textbf{Retendre} la courroie ou \textbf{poser un galet tendeur} ;
\item \textbf{Remplacer} la courroie si elle est usée ;
\item \textbf{Passer à la courroie crantée} si le problème persiste ou si le glissement est interdit.
\end{enumerate}
\end{methode}

\begin{exoresolu}[La pompe à eau de quartier]
Une pompe à eau entraînée par un moteur à essence siffle au démarrage et débite mal. Le technicien constate que la courroie s'enfonce de \SI{4}{cm} sous le pouce et que les poulies sont couvertes d'huile provenant d'une fuite du moteur. Quelles corrections doit-il apporter, dans l'ordre ?
\tcblower
\textbf{Solution détaillée.}

\textbf{Diagnostic.} Deux causes de glissement sont identifiées :
\begin{itemize}
\item tension insuffisante : \SI{4}{cm} d'enfoncement au lieu d'environ \SI{1}{cm} ;
\item présence d'huile sur les poulies, qui détruit l'adhérence.
\end{itemize}

\textbf{Corrections, dans l'ordre de la démarche :}
\begin{enumerate}
\item \textbf{Réparer la fuite d'huile} du moteur (sinon la souillure reviendra) ;
\item \textbf{Nettoyer} les gorges des poulies et la courroie avec un chiffon et du dégraissant ;
\item \textbf{Retendre la courroie} en desserrant les boulons du moteur, en écartant celui-ci sur ses rails jusqu'à obtenir environ \SI{1}{cm} d'enfoncement au pouce, puis en resserrant ;
\item Si la courroie est imprégnée d'huile en profondeur ou usée, la \textbf{remplacer} par une courroie neuve de même dimension ;
\item Faire un \textbf{essai} : le sifflement doit disparaître et le débit redevenir normal.
\end{enumerate}

\textbf{Conclusion :} en appliquant la démarche méthodique (nettoyer, retendre, remplacer si besoin), le technicien supprime le glissement sans changer tout le système.
\end{exoresolu}

\begin{exercice}[À toi de jouer]
Le moulin à maïs du village patine quand on verse trop de grains d'un coup (surcharge), et sa courroie, vieille de cinq ans, est lisse et craquelée.
\begin{enumerate}
\item Cite les deux causes du glissement observé.
\item Propose une correction adaptée à chaque cause.
\item Pourquoi une courroie crantée réglerait-elle définitivement le problème de glissement ?
\end{enumerate}
\end{exercice}

\begin{corrige}[Exercice : le moulin à maïs]
\begin{enumerate}
\item Les deux causes sont : la \textbf{surcharge} (trop de grains d'un coup, l'effort dépasse l'adhérence) et l'\textbf{usure} de la courroie (flancs lisses et craquelés après cinq ans).
\item Pour la surcharge : verser les grains progressivement pour ne pas dépasser l'effort transmissible. Pour l'usure : \textbf{remplacer} la courroie par une courroie neuve de bonne dimension, puis vérifier sa tension.
\item La courroie crantée transmet le mouvement par \textbf{engrènement} de ses dents dans les rainures des poulies dentées, et non plus par adhérence : le glissement devient impossible, même en cas de forte charge.
\end{enumerate}
\end{corrige}

Tu sais désormais diagnostiquer un glissement et choisir la bonne méthode pour le corriger, du simple réglage de tension jusqu'à la courroie synchrone. Il reste un point essentiel pour utiliser ces machines sans danger : la sécurité et l'entretien, que nous étudierons dans la dernière section.

\coursec{Sécurité et entretien}

Dans la section précédente, tu as appris à corriger les glissements d'une courroie : la retendre, la remplacer, ou augmenter l'angle d'enroulement. Mais toutes ces opérations ont un point commun : elles obligent à \textbf{approcher les mains} d'une transmission en mouvement. Or une courroie qui tourne est un piège redoutable. Cette dernière section est peut-être la plus importante de la leçon : elle peut te sauver la main, voire la vie, à l'atelier ou au village.

\courssub{Pourquoi une courroie en mouvement est-elle dangereuse ?}

\courssub{Le phénomène de happement}

Observe une courroie qui tourne entre deux poulies. À l'endroit où la courroie \emph{entre} dans la poulie, il existe un point de pincement : la courroie et la gorge de la poulie se rapprochent l'une de l'autre en tournant. Si un doigt, un pan de chemise, une manche ou une mèche de cheveux touche ce point, il est \textbf{aspiré et entraîné} entre la courroie et la poulie. On dit que la personne est \cle{happée}.

\begin{definition}[Happement]
Le \cle{happement} est l'entraînement brutal d'une partie du corps (doigt, main, cheveux) ou d'un vêtement par un organe en mouvement, ici la courroie et la poulie. La force exercée est celle du moteur : elle est bien supérieure à la force d'un être humain, qui ne peut pas se dégager.
\end{definition}

\begin{experience}[Observation : le point de pincement]
\textbf{Matériel :} deux poulies montées sur un support, une courroie, une manivelle (démonstration par le professeur, moteur débranché).
\begin{enumerate}
\item Tourner lentement la manivelle et observer l'endroit où la courroie rejoint la poulie.
\item Approcher (sans toucher !) un crayon tenu par le professeur du point d'entrée de la courroie.
\item Lâcher légèrement le crayon au contact : il est violemment entraîné et coincé entre courroie et poulie.
\end{enumerate}
\textbf{Observation :} le crayon est happé instantanément, sans possibilité de le retenir. \textbf{Interprétation :} si un simple crayon ne peut être arraché, un doigt ne le pourra pas non plus. \textbf{Conclusion :} on ne touche jamais une transmission en mouvement, même « juste un peu ».
\end{experience}

\begin{attention}[Une courroie qui casse peut fouetter]
Une courroie usée, fissurée ou trop tendue peut \textbf{rompre en plein fonctionnement}. Sous l'effet de la tension accumulée, les brins fouettent violemment l'air, comme un fouet, et peuvent blesser les yeux ou le visage des personnes proches. C'est pourquoi le \cle{carter de protection} (capot qui recouvre la transmission) est \textbf{obligatoire} sur les machines d'atelier : il protège à la fois du happement et des projections en cas de rupture.
\end{attention}

\courssub{Les quatre consignes de sécurité}

\begin{aretenir}[Consignes de sécurité autour d'une transmission par courroie]
\begin{enumerate}
\item \textbf{Ne jamais approcher les mains} d'une courroie en rotation. Avant toute intervention (réglage, nettoyage, graissage), \textbf{arrêter le moteur} et attendre l'arrêt complet de la courroie.
\item \textbf{Installer un carter de protection} (capot) sur la transmission : il interdit tout contact accidentel avec la courroie et les poulies.
\item \textbf{Retendre ou remplacer une courroie uniquement machine arrêtée et débranchée} : couper l'alimentation électrique (disjoncteur, prise débranchée) ou le moteur thermique, pour empêcher tout redémarrage intempestif.
\item \textbf{Porter des vêtements ajustés} et \textbf{attacher les cheveux longs} à l'atelier ; retirer bagues, bracelets et foulards qui peuvent être happés.
\end{enumerate}
\end{aretenir}

\begin{popfigure}
\begin{tikzpicture}[scale=1.0, font=\small]
% Carter de protection (grand rectangle arrondi)
\fill[popmuted!25] (0.2,-1.6) rectangle (8.3,2.1);
\draw[line width=1.5pt, popdark, rounded corners=6pt] (0.2,-1.6) rectangle (8.3,2.1);
% Poulie motrice
\draw[fill=popblueL, draw=popblue, line width=1pt] (1.7,0.3) circle (0.75);
\draw[fill=white, draw=popblue] (1.7,0.3) circle (0.12);
\node[popdark] at (1.7,-1.15) {Poulie motrice};
% Poulie réceptrice
\draw[fill=poporangeL, draw=poporange, line width=1pt] (6.8,0.3) circle (0.5);
\draw[fill=white, draw=poporange] (6.8,0.3) circle (0.12);
\node[popdark] at (6.8,-1.15) {Poulie r\'eceptrice};
% Courroie
\draw[line width=2pt, popdark] (1.7,1.05) -- (6.8,0.8);
\draw[line width=2pt, popdark] (1.7,-0.45) -- (6.8,-0.2);
% Sens de rotation
\draw[-{Stealth[length=2.5mm]}, popblue, line width=1pt] (1.7,0.3) ++(115:0.95) arc (115:55:0.95);
\draw[-{Stealth[length=2.5mm]}, poporange, line width=1pt] (6.8,0.3) ++(115:0.7) arc (115:55:0.7);
% Fixations du carter
\foreach \x/\y in {0.45/1.85, 8.05/1.85, 0.45/-1.35, 8.05/-1.35}{
  \draw[fill=popdark] (\x,\y) circle (0.07);
}
\node[popdark, align=center] at (4.25,1.7) {\textbf{Carter de protection}};
% Pictogramme danger
\begin{scope}[shift={(4.25,-2.6)}]
\draw[line width=1.5pt, poporange, fill=poppinkL] (0,0.75) -- (0.85,-0.55) -- (-0.85,-0.55) -- cycle;
\node[font=\Large\bfseries, popdark] at (0,-0.12) {!};
\node[popdark, align=center, font=\footnotesize] at (0,-1.0) {Danger : organes en mouvement\\ Ne pas ouvrir en fonctionnement};
\end{scope}
% Flèche point de pincement
\draw[-{Stealth[length=2.5mm]}, poppink, line width=1.2pt] (3.1,-1.9) -- (2.35,-0.35);
\node[poppink, align=center, font=\footnotesize] at (3.3,-2.15) {Point de pincement\\ (zone de happement)};
\end{tikzpicture}
\legende{Transmission par poulies-courroie protégée par un carter. Le capot (grisé) recouvre entièrement les poulies et la courroie ; le pictogramme triangulaire signale le danger. Le point de pincement, à l'entrée de la courroie sur la poulie, est la zone la plus dangereuse.}
\end{popfigure}

\courssub{L'entretien courant de la transmission}

Une courroie bien entretenue dure longtemps, transmet bien le mouvement et ne casse pas. L'entretien comporte trois vérifications régulières, toujours réalisées \textbf{machine arrêtée et débranchée} :

\begin{itemize}
\item \textbf{La tension} : une courroie trop molle glisse ; une courroie trop tendue use les paliers et peut rompre.
\item \textbf{L'état} : rechercher des fissures, des craquelures, des bords effilochés ou une surface devenue lisse et brillante (signe d'usure). Une courroie fissurée doit être remplacée.
\item \textbf{La propreté} : enlever la poussière, l'huile et la graisse déposées sur la courroie et dans les gorges des poulies, car elles favorisent le glissement.
\end{itemize}

\begin{methode}[Vérifier la tension d'une courroie (test du pouce)]
\begin{enumerate}
\item Arrêter le moteur et \textbf{débrancher} la machine.
\item Repérer le \textbf{brin} le plus long (la partie de courroie tendue entre les deux poulies).
\item Appuyer franchement avec le \textbf{pouce} au milieu de ce brin, perpendiculairement à la courroie.
\item Mesurer l'enfoncement : il doit être \textbf{faible, environ 1 à \SI{2}{cm}} selon la longueur du brin.
\item Si la courroie s'enfonce trop, elle est détendue : \textbf{écarter le moteur} pour la retendre. Si elle ne s'enfonce presque pas, elle est trop tendue : \textbf{rapprocher légèrement le moteur}.
\end{enumerate}
\end{methode}

\begin{exemplebox}[Les moulins à maïs au village]
Dans beaucoup de villages du Cameroun, le moulin à maïs ou à manioc est entraîné par un moteur diesel relié à la meule par une courroie. Les accidents de \textbf{mains happées} y sont fréquents : un enfant joue près de la machine, une maman veut « aider » la courroie qui glisse en poussant dessus pendant que le moteur tourne... Deux mesures simples évitent presque tous ces drames : \textbf{poser un carter} (même une simple tôle pliée) autour de la transmission, et \textbf{arrêter le moteur avant toute intervention}. La sécurité n'est pas un luxe : c'est un réflexe.
\end{exemplebox}

\begin{exoresolu}[Que faut-il faire ?]
Énoncé. Au moulin du village, la courroie du moteur diesel glisse et le meunier veut la retendre pendant que le moteur tourne, « pour gagner du temps ». (1) Citer deux dangers auxquels il s'expose. (2) Décrire la bonne conduite à tenir.
\tcblower
\textbf{Solution.} (1) En intervenant moteur en marche, il risque d'abord le \textbf{happement} de la main ou d'un vêtement au point de pincement entre courroie et poulie ; il risque aussi d'être blessé par le \textbf{fouettement de la courroie} si celle-ci, déjà usée, casse pendant la manipulation. (2) La bonne conduite : arrêter le moteur diesel et attendre l'arrêt complet de la courroie ; s'assurer que personne ne peut le redémarrer ; vérifier l'état de la courroie (fissures) ; écarter le moteur pour retendre la courroie, puis contrôler la tension au test du pouce (enfoncement de 1 à \SI{2}{cm}) ; enfin remettre le carter de protection en place avant de redémarrer.
\end{exoresolu}

\begin{attention}[Pièges à éviter]
\begin{itemize}
\item Ne jamais « aider » une courroie qui glisse en appuyant dessus pendant la rotation : c'est la cause la plus fréquente de happement.
\item Ne jamais retirer le carter « parce qu'il gêne » : il doit être remis en place après chaque intervention.
\item Ne pas confondre \emph{arrêt} et \emph{débranchement} : une machine simplement arrêtée peut redémarrer par erreur ; il faut couper l'alimentation.
\end{itemize}
\end{attention}

\begin{exercice}[À toi de jouer]
(1) Explique pourquoi les cheveux longs doivent être attachés à l'atelier. (2) Lors du test du pouce, la courroie d'une pompe à eau s'enfonce de \SI{4}{cm}. Que conclus-tu et que fais-tu ? (3) Pourquoi le carter protège-t-il aussi en cas de rupture de la courroie ?
\end{exercice}

\begin{corrige}[Exercice]
(1) Les cheveux longs peuvent être happés par la courroie ou la poulie en rotation et entraîner le cuir chevelu ; on les attache pour les éloigner des organes en mouvement. (2) La courroie est détendue : elle risque de glisser. Machine arrêtée et débranchée, on écarte le moteur pour la retendre jusqu'à obtenir un enfoncement de 1 à \SI{2}{cm}. (3) En cas de rupture, les brins de la courroie fouettent violemment ; le carter les retient à l'intérieur et protège les personnes des projections.
\end{corrige}

\begin{aretenir}[L'essentiel]
Une courroie en mouvement peut \textbf{happer} doigts, vêtements et cheveux, et une courroie qui casse peut \textbf{fouetter}. Les quatre réflexes : mains éloignées, moteur arrêté et machine débranchée avant toute intervention, carter de protection en place, vêtements ajustés et cheveux attachés. L'entretien régulier (tension au test du pouce, état, propreté) garantit une transmission sûre et durable.
\end{aretenir}

\newpage

\coursec{Activité d'intégration}

\coursec{Système poulies-courroie}

\courssub{Objectifs de l'activité}

À la fin de cette activité, tu seras capable de :

\begin{itemize}
\item identifier sur un schéma réel les organes d'un système poulies-courroie : poulie menante, poulie menée, courroie, brin tendu et brin mou ;
\item calculer le rapport de transmission et la vitesse théorique de l'arbre mené à partir des diamètres des poulies ;
\item expliquer le phénomène de \cle{glissement} de la courroie et en citer les causes possibles ;
\item proposer des méthodes concrètes de correction du glissement ;
\item rédiger une consigne de sécurité adaptée à une machine entraînée par poulies et courroie.
\end{itemize}

\courssub{Situation}

\textbf{Le moulin à maïs de M. Etoundi qui patine.}

Au quartier Nkolndongo à Yaoundé, M. Etoundi exploite un petit moulin à maïs qui rend service à tout le voisinage : les mamans y apportent le maïs sec pour obtenir de la farine destinée au couscous et au bâton de manioc. Le moulin est entraîné par un moteur électrique relié à la meule par un système \cle{poulies-courroie}.

Les caractéristiques de l'installation sont les suivantes :

\begin{itemize}
\item le moteur électrique tourne à la vitesse $N_1 = \SI{1500}{tr/min}$ ;
\item la poulie fixée sur l'arbre du moteur a un diamètre $D_1 = \SI{80}{mm}$ ;
\item la poulie fixée sur l'arbre de la meule a un diamètre $D_2 = \SI{240}{mm}$ ;
\item la courroie est une courroie trapézoïdale en caoutchouc, installée il y a quatre ans.
\end{itemize}

Depuis quelques jours, les clientes se plaignent : la mouture est devenue très lente. M. Etoundi observe sa machine et témoigne :

\begin{quote}
\emph{« Depuis quelques jours, la meule tourne lentement, la courroie siffle et sent le brûlé ; le brin du bas pendouille. Pourtant, le moteur ronronne normalement. »}
\end{quote}

Le schéma de l'installation est représenté ci-dessous.

\begin{popfigure}
\begin{center}
\begin{tikzpicture}[scale=0.9]
% Poulie motrice
\draw[thick,fill=popblueL] (0,0) circle (0.8);
\draw[fill=popdark] (0,0) circle (0.08);
\node[below] at (0,-0.95) {\small Poulie 1 ($D_1 = \SI{80}{mm}$)};
\node[above] at (0,0.95) {\small Moteur $\SI{1500}{tr/min}$};
% Poulie menée
\draw[thick,fill=poporangeL] (6,0) circle (2.4);
\draw[fill=popdark] (6,0) circle (0.08);
\node[below] at (6,-2.55) {\small Poulie 2 ($D_2 = \SI{240}{mm}$)};
\node at (6,0.45) {\small Meule};
% Courroie brin tendu (haut)
\draw[very thick,popgreen] (0,0.8) -- (6,2.4);
% Courroie brin mou (bas) qui pendouille
\draw[very thick,poporange] (0,-0.8) .. controls (2,-1.9) and (4,-1.9) .. (6,-2.4);
% Flèches de rotation
\draw[->,thick,popblue] (-0.9,0.5) arc (140:40:0.9);
\draw[->,thick,poporange] (4.6,1.9) arc (140:40:0.9);
% Légendes fléchées
\draw[->,popgreen] (3,1.75) -- (3,1.05);
\node[above,popgreen] at (3,1.8) {\small Brin tendu};
\draw[->,poporange] (3,-1.15) -- (3,-1.75);
\node[below,poporange] at (3,-1.85) {\small Brin mou (pendouille)};
\end{tikzpicture}
\end{center}
\legende{Schéma du moulin à maïs de M. Etoundi : le moteur entraîne la meule par l'intermédiaire de deux poulies et d'une courroie. Le brin inférieur est détendu.}
\end{popfigure}

En tant que futur technicien, M. Etoundi te demande de l'aider à comprendre la panne et à proposer des solutions.

\courssub{Consignes}

\begin{enumerate}
\item \textbf{Identification des organes.} Sur le schéma de l'installation, identifie et nomme : la poulie menante, la poulie menée, la courroie, le brin tendu et le brin mou. Justifie tes choix.
\item \textbf{Calcul de la vitesse théorique.} Calcule le rapport de transmission $r = \dfrac{D_1}{D_2}$, puis la vitesse de rotation théorique $N_2$ de la meule. La meule devrait-elle tourner plus vite ou moins vite que le moteur ?
\item \textbf{Analyse de la panne.} La vitesse réelle de la meule est inférieure à la vitesse théorique calculée. À l'aide du témoignage de M. Etoundi (sifflement, odeur de brûlé, brin qui pendouille), nomme le phénomène responsable et explique en deux ou trois phrases pourquoi la vitesse réelle est inférieure à la vitesse théorique.
\item \textbf{Causes du glissement.} Cite deux causes possibles du glissement observé sur ce moulin.
\item \textbf{Solutions.} Propose trois solutions concrètes que M. Etoundi peut appliquer pour corriger le glissement et redonner à son moulin son rendement normal.
\item \textbf{Sécurité.} Rédige une consigne de sécurité courte et claire (une ou deux phrases) à afficher sur le moulin pour protéger l'opérateur et les clientes.
\end{enumerate}

\courssub{Ressources à mobiliser}

\begin{propriete}[Avantages et inconvénients du système poulies-courroie]
\begin{center}
\begin{tabularx}{\linewidth}{|l|X|X|}\hline
\rowcolor{popblueL}
\textbf{Critère} & \textbf{Avantages} & \textbf{Inconvénients} \\ \hline
Transmission & Transmet le mouvement entre deux axes éloignés (distance variable). & Glissement possible (perte de vitesse, non synchronisme). \\ \hline
Fonctionnement & Fonctionnement silencieux ; amortit les chocs et vibrations du moteur. & Usure de la courroie (vieillissement, élongation, glaçage). \\ \hline
Coût / Entretien & Coût faible ; montage et démontage faciles ; pas de lubrification. & Encombrement radial important ; tension à surveiller régulièrement. \\ \hline
Rapport & Rapport de transmission modifiable en changeant les poulies. & Rapport limité (généralement $< 1/7$ par étage) pour éviter le glissement. \\ \hline
\end{tabularx}
\end{center}
\end{propriete}

\begin{aretenir}[Ce qu'il faut se rappeler]
\begin{itemize}
\item Dans un système poulies-courroie, la \cle{poulie menante} est fixée sur l'arbre du moteur ; elle entraîne la \cle{courroie}, qui entraîne à son tour la \cle{poulie menée}.
\item Le \cle{brin tendu} est le brin de la courroie qui tire la poulie menée ; le \cle{brin mou} est l'autre brin, légèrement détendu.
\item Rapport de transmission : \[ r = \frac{N_2}{N_1} = \frac{D_1}{D_2} \qquad \text{donc} \qquad N_2 = N_1 \times \frac{D_1}{D_2} \]
\item Si $D_1 < D_2$, alors $N_2 < N_1$ : le système est un \cle{réducteur} de vitesse (augmentation du couple).
\item Si $D_1 > D_2$, alors $N_2 > N_1$ : le système est un \cle{multiplicateur} de vitesse.
\item Le \cle{glissement} est le patinage de la courroie sur les poulies : la vitesse réelle devient inférieure à la vitesse théorique.
\item Causes fréquentes du glissement : courroie détendue (tension insuffisante), courroie usée ou glacée, surcharge de la machine, huile ou graisse sur la courroie ou les gorges.
\item Corrections possibles : retendre la courroie en écartant le moteur, poser un \cle{galet tendeur}, remplacer la courroie usée, nettoyer les gorges et la courroie, éviter la surcharge.
\end{itemize}
\end{aretenir}

\begin{methode}[Corriger un glissement de courroie]
\begin{enumerate}
\item \textbf{Arrêter la machine} et couper l'alimentation électrique (consignation).
\item \textbf{Inspecter} : état de la courroie (fissures, glaçage, déchirures), propreté des gorges, tension du brin mou (flèche normale $\approx 10$ à $15$ mm par mètre d'entraxe).
\item \textbf{Agir} selon le diagnostic :
      \begin{itemize}
      \item Courroie saine mais détendue $\rightarrow$ écarter le moteur (ou le galet tendeur) pour retendre, puis bloquer.
      \item Courroie usée, glacée ou huilée $\rightarrow$ remplacer par une neuve de même profil et nettoyer les gorges.
      \item Surcharge $\rightarrow$ réduire la charge ou vérifier le dimensionnement.
      \end{itemize}
\item \textbf{Revérifier} la tension après quelques heures de fonctionnement (rodage).
\end{enumerate}
\end{methode}

\courssub{Démarche et corrigé commenté}

\begin{exoresolu}[Le moulin à maïs de M. Etoundi : corrigé complet]
Énoncé : répondre aux six consignes de l'activité à partir des données $N_1 = \SI{1500}{tr/min}$, $D_1 = \SI{80}{mm}$, $D_2 = \SI{240}{mm}$ et du témoignage de M. Etoundi.

\tcblower

\textbf{1. Identification des organes.}

\begin{itemize}
\item \textbf{Poulie menante} : la poulie 1 ($D_1 = \SI{80}{mm}$), car elle est fixée sur l'arbre du moteur électrique. C'est elle qui fournit le mouvement.
\item \textbf{Poulie menée} : la poulie 2 ($D_2 = \SI{240}{mm}$), car elle est fixée sur l'arbre de la meule et reçoit le mouvement.
\item \textbf{Courroie} : la bande de caoutchouc trapézoïdale qui relie les deux poulies et transmet le mouvement par adhérence (frottement) dans les gorges.
\item \textbf{Brin tendu} : le brin supérieur de la courroie, celui qui tire la poulie menée dans le sens de rotation (flèche verte sur le schéma).
\item \textbf{Brin mou} : le brin inférieur, qui « pendouille » d'après le témoignage. Un brin mou qui pend anormalement est un indice visuel fort de détente de la courroie.
\end{itemize}

\textbf{2. Calcul de la vitesse théorique de la meule.}

On applique la relation du rapport de transmission :
\begin{align*}
r &= \frac{D_1}{D_2} = \frac{80}{240} = \frac{1}{3} \approx 0{,}33 \\
N_2 &= N_1 \times \frac{D_1}{D_2} = 1500 \times \frac{80}{240} = 1500 \times \frac{1}{3} = \SI{500}{tr/min}
\end{align*}

La vitesse théorique de la meule est donc $N_2 = \SI{500}{tr/min}$. Comme $D_1 < D_2$, la meule tourne \textbf{moins vite} que le moteur : le système joue le rôle de réducteur de vitesse, ce qui est souhaitable pour un moulin (plus de force, plus de couple pour écraser les grains).

\textbf{3. Analyse de la panne.}

Le phénomène responsable est le \cle{glissement} (ou patinage) de la courroie sur les poulies. Normalement, la courroie adhère aux gorges des poulies et entraîne la poulie menée sans perte de vitesse (hypothèse du « sans glissement »). Ici, la courroie est détendue (le brin du bas pendouille) : elle n'appuie plus assez fort sur les parois des gorges, donc elle patine. Ce patinage produit un frottement anormal qui fait \emph{siffler} la courroie et dégage une \emph{odeur de brûlé} (le caoutchouc surchauffe par échauffement de frottement). Pendant ce temps, une partie du mouvement du moteur n'est plus transmise à la meule : la vitesse réelle de la meule devient inférieure à la vitesse théorique de $\SI{500}{tr/min}$, d'où la mouture lente.

\textbf{4. Deux causes possibles du glissement.}

\begin{itemize}
\item La courroie s'est \textbf{détendue} avec le temps (quatre ans de service, élongation naturelle du caoutchouc) : sa tension est devenue insuffisante, ce que confirme le brin mou qui pendouille anormalement.
\item La courroie est \textbf{usée ou glacée} : ses flancs, polis par quatre ans de frottements, sont devenus lisses et brillants ; ils n'adhèrent plus correctement aux gorges des poulies, même si la tension était correcte.
\end{itemize}

(On pourrait aussi citer une surcharge ponctuelle du moulin, si M. Etoundi verse trop de maïs d'un coup dans la trémie, ce qui demande un couple supérieur à ce que l'adhérence courroie-poulie peut transmettre.)

\textbf{5. Trois solutions concrètes.}

\begin{enumerate}
\item \textbf{Retendre la courroie} en écartant légèrement le moteur de la meule sur ses rails de fixation (ou en agissant sur le galet tendeur s'il existe), puis en resserrant fermement les boulons de fixation du moteur.
\item \textbf{Poser un galet tendeur} (roue libre montée sur ressort ou sur excentrique) qui appuie sur le brin mou pour maintenir la tension de façon permanente et compenser l'élongation future.
\item \textbf{Remplacer la courroie usée} par une courroie neuve de même profil (même référence trapézoïdale), car une courroie glacée ne retrouvera jamais son adhérence d'origine, même bien tendue.
\end{enumerate}

\textbf{6. Consigne de sécurité à afficher.}

\begin{quote}
\emph{« Arrêter le moteur et débrancher la prise avant toute intervention sur la courroie. Ne jamais approcher les mains, les vêtements amples ou les cheveux longs des poulies en rotation. »}
\end{quote}
\end{exoresolu}

\courssub{Bilan de l'activité}

Cette activité t'a permis de mobiliser l'ensemble des savoirs de la leçon dans une situation authentique de la vie courante camerounaise. Tu as montré que :

\begin{itemize}
\item un système poulies-courroie se décrit par ses organes (poulie menante, poulie menée, courroie, brins tendu et mou) et par son rapport de transmission $r = \dfrac{D_1}{D_2}$ ;
\item le calcul théorique ($N_2 = \SI{500}{tr/min}$) donne une valeur de référence : tout écart entre la vitesse réelle et la vitesse théorique révèle un dysfonctionnement, ici le glissement ;
\item les signes du glissement sont reconnaissables dans la pratique : sifflement, odeur de brûlé, brin mou détendu, baisse de rendement de la machine ;
\item le technicien dispose de méthodes simples de correction : retendre la courroie, installer un galet tendeur ou remplacer la courroie usée ;
\item toute intervention sur une machine tournante exige le respect de consignes de sécurité strictes (consignation, protection des parties tournantes).
\end{itemize}

Tu es maintenant capable de diagnostiquer et de corriger une panne courante sur les nombreuses machines à courroie utilisées au Cameroun : moulins, décortiqueuses, pompes à eau, groupes électrogènes ou alternateurs de moto-taxi.

\courssub{Exercices d'entraînement}

\begin{exercice}[Alternateur de moto-taxi]
Sur une moto-taxi type « bend skin » très répandue à Douala, l'alternateur est entraîné par le vilebrequin du moteur thermique via un système poulies-courroie. Le vilebrequin (poulie menante) a un diamètre $D_1 = \SI{60}{mm}$ et tourne à $N_1 = \SI{3000}{tr/min}$ au ralenti. L'alternateur (poulie menée) doit tourner à $N_2 = \SI{9000}{tr/min}$ pour charger correctement la batterie.
\begin{enumerate}
\item Calcule le diamètre $D_2$ nécessaire de la poulie d'alternateur.
\item Précise si ce système est un réducteur ou un multiplicateur de vitesse.
\item Le mécanicien constate que la batterie ne charge plus au ralenti, la courroie siffle. Donne une cause probable et la solution.
\end{enumerate}
\end{exercice}

\begin{corrige}[Exercice n°1 : Alternateur de moto-taxi]
\begin{enumerate}
\item On utilise la relation $N_2 = N_1 \times \dfrac{D_1}{D_2}$, d'où $D_2 = D_1 \times \dfrac{N_1}{N_2} = 60 \times \dfrac{3000}{9000} = 60 \times \dfrac{1}{3} = \SI{20}{mm}$.
\item Comme $N_2 > N_1$ (et $D_2 < D_1$), c'est un \textbf{multiplicateur} de vitesse.
\item Cause probable : courroie détendue ou glacée (usée) $\rightarrow$ glissement au ralenti où le couple résistif de l'alternateur (charge batterie) est maximal. Solution : retendre ou remplacer la courroie.
\end{enumerate}
\end{corrige}

\begin{exercice}[Pompe à eau de forage]
Une pompe à eau de forage à Bafoussam est entraînée par un moteur électrique ($N_1 = \SI{1450}{tr/min}$) via une courroie trapézoïdale. La poulie moteur mesure $D_1 = \SI{100}{mm}$, la poulie pompe $D_2 = \SI{300}{mm}$.
\begin{enumerate}
\item Calcule la vitesse théorique de la pompe $N_2$.
\item Après deux ans, le débit de la pompe baisse. On mesure $N_2^{\text{réel}} = \SI{420}{tr/min}$. Calcule le taux de glissement $g = \dfrac{N_2^{\text{théo}} - N_2^{\text{réel}}}{N_2^{\text{théo}}} \times 100$.
\item Propose deux actions de maintenance préventive pour limiter ce glissement.
\end{enumerate}
\end{exercice}

\begin{corrige}[Exercice n°2 : Pompe à eau de forage]
\begin{enumerate}
\item $N_2 = 1450 \times \dfrac{100}{300} = 1450 \times \dfrac{1}{3} \approx \SI{483}{tr/min}$ (arrondi à l'entier).
\item $N_2^{\text{théo}} = \SI{483}{tr/min}$. Taux de glissement $g = \dfrac{483 - 420}{483} \times 100 \approx \SI{13}{\%}$. C'est un glissement important (normalement $< 2\text{--}3\%$).
\item Actions préventives : (1) Vérifier et retendre la courroie tous les 6 mois (ou installer un galet tendeur automatique). (2) Nettoyer régulièrement les gorges des poulies (poussière, boue, graisse) et remplacer la courroie dès apparition de fissures ou de glaçage.
\end{enumerate}
\end{corrige}

\newpage

\coursec{Méthodes et exercices résolus}

\coursec{Système poulies-courroie}

\noindent Dans nos ateliers, nos maisons et sur nos routes, le mouvement doit souvent être transmis entre deux arbres \textbf{éloignés} l'un de l'autre : le moteur d'un groupe électrogène et son alternateur, le moteur d'une moto-taxi et sa roue arrière (par courroie ou chaîne), l'arbre d'un moulin à maïs et son moteur électrique. Les engrenages, qui nécessitent un contact direct, ne conviennent pas. C'est là qu'intervient le \cle{système poulies-courroie}, un organe de transmission souple, silencieux et économique, omniprésent au Cameroun.

\courssub{Observation et définition du système}

\begin{definition}[Système poulies-courroie]
Un \textbf{système poulies-courroie} est un dispositif de transmission de mouvement et de puissance entre deux arbres parallèles et \textbf{éloignés}. Il est constitué de trois éléments principaux :
\begin{itemize}
\item la \cle{poulie motrice} (ou émettrice), fixée sur l'arbre du moteur (source d'énergie) ;
\item la \cle{poulie menée} (ou réceptrice), fixée sur l'arbre de la machine à entraîner ;
\item la \cle{courroie}, organe souple (caoutchouc, cuir, matière synthétique, souvent trapézoïdale) qui relie les deux poulies et transmet le mouvement \textbf{par adhérence} (frottement).
\end{itemize}
\end{definition}

\begin{popfigure}
\begin{tikzpicture}[scale=0.75]
% Poulie motrice (gauche)
\draw[popblue, thick] (0,0) circle (2);
\draw[popblue, thick] (0,0) circle (0.4); % alésage
\draw[<->, thick] (0,2) -- (0,3.5) node[midway, right] {$D_1$};
\draw[->, thick] (0,2) arc (90:450:2);
\node at (0, -2.8) {\textbf{Poulie motrice}};
\node at (0, -3.3) {\small (Arbre moteur)};
\node at (0,0) {$N_1$};

% Courroie (brins sup et inf)
\draw[poporange, very thick] (2, 2) -- (7, 1.5);
\draw[poporange, very thick] (2, -2) -- (7, -1.5);
\node at (4.5, 2.6) {\textbf{Courroie}};

% Poulie menée (droite)
\draw[popgreen, thick] (9,0) circle (1.5);
\draw[popgreen, thick] (9,0) circle (0.4);
\draw[<->, thick] (9,1.5) -- (9,3) node[midway, right] {$D_2$};
\draw[->, thick] (9,1.5) arc (90:450:1.5);
\node at (9, -2.8) {\textbf{Poulie menée}};
\node at (9, -3.3) {\small (Machine entraînée)};
\node at (9,0) {$N_2$};

% Entraxe
\draw[<->, dashed, popmuted] (2, -3.5) -- (7.5, -3.5) node[midway, below] {Entraxe $E$};
\end{tikzpicture}
\legende{Schéma de principe d'un système poulies-courroie à courroie ouverte. $D_1$, $D_2$ : diamètres primitifs (en mm ou cm) ; $N_1$, $N_2$ : vitesses de rotation (en tr/min) ; $E$ : distance entre les axes (entraxe).}
\end{popfigure}

\begin{experience}[Observer un système réel : le moulin à maïs ou le groupe électrogène]
\textbf{Matériel :} moulin à maïs (ou groupe électrogène) accessible, mètre ruban, craie.
\textbf{Protocole :}
\begin{enumerate}
\item \textbf{Arrêt total de la machine} (sécurité). Identifier visuellement la poulie motrice (côté moteur) et la poulie menée (côté meules/alternateur).
\item Mesurer les diamètres $D_1$ et $D_2$ (au pied à coulisse ou au mètre sur la gorge de la courroie).
\item Noter le type de courroie (plate, trapézoïdale, crantée) et son état (neuve, usée, fissurée, graissée).
\item Vérifier la tension : appuyer sur le brin tendu au milieu de l'entraxe ; la flèche ne doit pas dépasser 10 à 15 mm.
\item Observer le sens de rotation des deux poulies (marquer une flèche à la craie sur chaque poulie, faire tourner le moteur à la main si possible).
\end{enumerate}
\textbf{Observations attendues :} La courroie relie deux arbres éloignés. Si la courroie n'est pas croisée, les flèches indiquent le même sens. La courroie trapézoïdale s'insère dans la gorge (meilleure adhérence).
\textbf{Interprétation :} La transmission se fait par frottement (adhérence). L'entraxe $E$ est grand (souvent > 50 cm), ce qui justifie le choix de ce système plutôt que des engrenages.
\end{experience}

\courssub{Fonctionnement et sens de rotation}

\begin{propriete}[Sens de rotation]
Le sens de rotation de la poulie menée par rapport à la poulie motrice dépend uniquement du montage de la courroie :
\begin{itemize}
\item \textbf{Courroie ouverte} (brins non croisés) : les deux poulies tournent dans le \textbf{même sens}.
\item \textbf{Courroie croisée} (brins croisés en un point) : les deux poulies tournent dans des \textbf{sens opposés}.
\end{itemize}
\emph{Remarque :} La courroie croisée augmente l'angle d'enroulement (meilleure adhérence) mais use la courroie plus vite au point de croisement. Elle est réservée aux faibles puissances.
\end{propriete}

\begin{popfigure}
\begin{tikzpicture}[scale=0.7]
% Courroie ouverte
\draw[popblue, thick] (0,0) circle (1.5);
\draw[popblue, thick] (5,0) circle (1);
\draw[poporange, thick] (1.5, 1.5) -- (4, 1);
\draw[poporange, thick] (1.5, -1.5) -- (4, -1);
\draw[->, thick] (0,1.5) arc (90:450:1.5);
\draw[->, thick] (5,1) arc (90:450:1);
\node at (0, -2.2) {\small Motrice};
\node at (5, -2.2) {\small Menée};
\node at (2.5, 2.0) {\textbf{Ouverte}};
\node at (2.5, -2.7) {\small Même sens};

% Courroie croisée
\begin{scope}[xshift=8.5cm]
\draw[popblue, thick] (0,0) circle (1.5);
\draw[popblue, thick] (5,0) circle (1);
\draw[poporange, thick] (1.5, 1.5) -- (4, -1);
\draw[poporange, thick] (1.5, -1.5) -- (4, 1);
\draw[->, thick] (0,1.5) arc (90:450:1.5);
\draw[->, thick] (5,-1) arc (-90:-450:1);
\node at (0, -2.2) {\small Motrice};
\node at (5, -2.2) {\small Menée};
\node at (2.5, 2.0) {\textbf{Croisée}};
\node at (2.5, -2.7) {\small Sens opposés};
\end{scope}
\end{tikzpicture}
\legende{Effet du montage de la courroie sur le sens de rotation. Flèches bleues = sens de rotation des poulies.}
\end{popfigure}

\courssub{Rapport de transmission et calculs de vitesse}

\begin{propriete}[Relation fondamentale (sans glissement)]
Si la courroie ne glisse pas sur les poulies (adhérence parfaite), la \textbf{vitesse linéaire} $v$ de la courroie est identique au contact des deux poulies :
\[ v = \pi \cdot D_1 \cdot N_1 = \pi \cdot D_2 \cdot N_2 \]
D'où la relation fondamentale, appelée \cle{rapport de transmission} $r$ :
\[ r = \frac{N_2}{N_1} = \frac{D_1}{D_2} \qquad \text{ou} \qquad D_1 \cdot N_1 = D_2 \cdot N_2 \]
\begin{itemize}
\item $D_1$, $D_2$ : diamètres primitifs des poulies (même unité : cm ou mm).
\item $N_1$, $N_2$ : vitesses de rotation (en tr/min ou tr/s).
\item $r < 1$ ($D_2 > D_1$) $\Rightarrow$ \textbf{réducteur de vitesse} (la menée tourne moins vite, couple multiplié).
\item $r > 1$ ($D_2 < D_1$) $\Rightarrow$ \textbf{multiplicateur de vitesse} (la menée tourne plus vite, couple divisé).
\end{itemize}
\end{propriete}

\begin{methode}[Calculer le rapport de transmission et la vitesse menée]
\begin{enumerate}
\item Identifier $D_1$ (motrice) et $D_2$ (menée) \textbf{dans la même unité}.
\item Calculer le rapport : $r = \dfrac{D_1}{D_2}$.
\item Calculer $N_2 = N_1 \times \dfrac{D_1}{D_2}$.
\item Interpréter : comparer $D_1$ et $D_2$ pour valider (réducteur ou multiplicateur).
\end{enumerate}
\end{methode}

\begin{methode}[Déterminer un diamètre inconnu]
\begin{enumerate}
\item Écrire $D_1 \cdot N_1 = D_2 \cdot N_2$.
\item Isoler l'inconnu : $D_2 = \dfrac{D_1 \cdot N_1}{N_2}$ (ou $D_1 = \dfrac{D_2 \cdot N_2}{N_1}$).
\item Vérifier la cohérence : « Plus vite $\Rightarrow$ plus petit » ; « Moins vite $\Rightarrow$ plus grand ».
\end{enumerate}
\end{methode}

\begin{exoresolu}[Le moulin à maïs de Bertoua]
Un moteur électrique tourne à $N_1 = \SI{1500}{tr/min}$ et porte une poulie motrice de diamètre $D_1 = \SI{8}{cm}$. Il entraîne, par courroie ouverte, la poulie d'un moulin à maïs de diamètre $D_2 = \SI{24}{cm}$.
\begin{enumerate}
\item Calculer le rapport de transmission $r$.
\item Calculer la vitesse de rotation $N_2$ du moulin.
\item S'agit-il d'un réducteur ou d'un multiplicateur ? Justifier. Les poulies tournent-elles dans le même sens ?
\end{enumerate}
\tcblower
\textbf{Solution détaillée.}
\begin{enumerate}
\item $r = \dfrac{D_1}{D_2} = \dfrac{8}{24} = \dfrac{1}{3} \approx 0,33$.
\item $N_2 = N_1 \times \dfrac{D_1}{D_2} = 1500 \times \dfrac{1}{3} = \SI{500}{tr/min}$.
\item $D_2 > D_1$ ($24 > 8$) $\Rightarrow$ $N_2 < N_1$ ($500 < 1500$) : c'est un \textbf{réducteur de vitesse}. Le moulin a besoin de force (couple) plus que de vitesse. Courroie ouverte $\Rightarrow$ \textbf{même sens de rotation}.
\end{enumerate}
\textbf{Conclusion :} $r = 1/3$, $N_2 = \SI{500}{tr/min}$, réducteur, même sens.
\end{exoresolu}

\begin{exoresolu}[La scie circulaire du menuisier]
Un menuisier de Yaoundé veut que la lame de sa scie circulaire tourne à $N_2 = \SI{3000}{tr/min}$. Son moteur tourne à $N_1 = \SI{1500}{tr/min}$ et porte une poulie motrice de diamètre $D_1 = \SI{20}{cm}$.
\begin{enumerate}
\item Quel diamètre $D_2$ doit avoir la poulie de la scie ?
\item Justifier qualitativement.
\end{enumerate}
\tcblower
\textbf{Solution détaillée.}
\begin{enumerate}
\item $D_2 = \dfrac{D_1 \cdot N_1}{N_2} = \dfrac{20 \times 1500}{3000} = \SI{10}{cm}$.
\item La scie doit tourner \textbf{2 fois plus vite} ($3000 = 2 \times 1500$). Pour multiplier la vitesse par 2, la poulie menée doit être \textbf{2 fois plus petite} : $D_2 = D_1 / 2 = \SI{10}{cm}$. C'est un \textbf{multiplicateur de vitesse}.
\end{enumerate}
\textbf{Conclusion :} $D_2 = \SI{10}{cm}$.
\end{exoresolu}

\courssub{Avantages et inconvénients}

\begin{methode}[Argumenter : avantages et inconvénients]
Pour comparer poulies-courroie et engrenages (ou analyser un choix technique) :
\begin{enumerate}
\item Citer \textbf{3 avantages} justifiés :
\begin{itemize}
\item Transmission entre arbres \textbf{éloignés} (entraxe grand) — impossible avec engrenages.
\item Fonctionnement \textbf{silencieux} (organe souple) — idéal en milieu urbain ou atelier.
\item \textbf{Absorption des chocs} et \textbf{protection} : en cas de blocage brutal (grain dans le moulin, bois dur dans la scie), la courroie \textbf{patine} (glisse) au lieu de casser l'arbre ou les dents d'engrenage. C'est un « fusible mécanique ».
\item Coût faible, montage et entretien simples.
\end{itemize}
\item Citer \textbf{2 inconvénients} :
\begin{itemize}
\item \textbf{Glissement} possible (perte de puissance, rapport non rigoureux, échauffement).
\item \textbf{Usure} de la courroie (vieillissement, chaleur, huile, poussière) $\Rightarrow$ remplacement périodique.
\item Encombrement au sol plus grand qu'un réducteur à engrenages compact.
\end{itemize}
\end{enumerate}
\end{methode}

\begin{exoresolu}[Choix technique à Bafoussam]
Un artisan hésite entre poulies-courroie et engrenages pour relier le moteur de sa décortiqueuse à café à l'arbre de la machine, distants de \SI{1,5}{m}.
\tcblower
\textbf{Solution.}
\begin{itemize}
\item \textbf{Avantage 1 (Distance) :} Seule la courroie permet de franchir \SI{1,5}{m} d'entraxe simplement.
\item \textbf{Avantage 2 (Sécurité) :} Si un grain bloque la machine, la courroie glisse et protège le moteur (pas de casse).
\item \textbf{Inconvénient :} Risque de glissement si la courroie se détend ou si de l'huile de moteur la souille $\Rightarrow$ perte de rendement, surveillance nécessaire.
\end{itemize}
\textbf{Conseil :} Choisir le système poulies-courroie (courroie trapézoïdale de préférence), prévoir un galet tendeur et un carter de protection.
\end{exoresolu}

\begin{savaistu}[Le saviez-vous ? La courroie au Cameroun]
\begin{itemize}
\item \textbf{Moto-taxi (\"benskin\") :} La plupart utilisent une \textbf{courroie crantée} (distribution) pour le moteur 4 temps, et une \textbf{courroie trapézoïdale} (variateur) pour la transmission vers la roue. La courroie crantée \textbf{ne glisse pas} (dents) : calage moteur précis garanti.
\item \textbf{Groupe électrogène :} La courroie relie le moteur thermique à l'alternateur. Un glissement fait chuter la tension (lumière qui faiblit). La tension de courroie est un entretien hebdomadaire.
\item \textbf{Gestion des déchets :} Les vieilles courroies en caoutchouc sont un déchet encombrant. Au Cameroun, des artisans les recyclent en semelles de sandales, joints d'étanchéité ou cordes pour l'agriculture. Économie circulaire !
\end{itemize}
\end{savaistu}

\courssub{Glissements : diagnostic et correction}

\begin{definition}[Glissement]
Le \textbf{glissement} est la perte d'adhérence entre la courroie et les poulies. La vitesse linéaire de la courroie devient différente de la vitesse périphérique des poulies. Conséquence : $N_2$ réel $< N_2$ théorique. Perte de puissance, échauffement, usure prématurée.
\end{definition}

\begin{methode}[Diagnostiquer et corriger le glissement]
\begin{enumerate}
\item \textbf{Signes :} Machine menée qui ralentit (moteur normal), \textbf{sifflement} aigu, odeur de caoutchouc brûlé, échauffement anormal des poulies, poudre noire (poussière de caoutchouc) sous le carter.
\item \textbf{Causes principales :}
\begin{itemize}
\item Courroie \textbf{détendue} (allongement naturel, entraxe mal réglé).
\item Courroie \textbf{usée} (gorges lisses, fond de gorge atteint, fissures).
\item \textbf{Graisse / huile} sur poulies ou courroie (fuite moteur, sur-graissage).
\item \textbf{Surcharge} de la machine (couple résistant > couple d'adhérence max).
\item Angle d'enroulement trop faible (poulie motrice trop petite, entraxe trop grand).
\end{itemize}
\item \textbf{Méthodes de correction (cause $\rightarrow$ remède) :}
\begin{itemize}
\item \textbf{Tendre la courroie} : déplacer le moteur sur ses glissières ou utiliser un \textbf{galet tendeur} (poulie folle sur brin détendu).
\item \textbf{Remplacer} la courroie usée/fissurée (même profil, même longueur).
\item \textbf{Nettoyer} poulies et courroie (dégraissant, essence, chiffon sec) et \textbf{réparer la fuite d'huile}.
\item Monter une \textbf{courroie trapézoïdale} (meilleure adhérence que plate) ou \textbf{crantée} (synchrone, glissement nul).
\item Augmenter l'angle d'enroulement (galet de renvoi sur brin détendu).
\end{itemize}
\end{enumerate}
\end{methode}

\begin{exoresolu}[Le moulin de Mme Ngo Bell]
Les meules tournent de moins en moins vite, moteur normal. Sifflement entendu. Traces d'huile sur poulie motrice.
\begin{enumerate}
\item Défaut observé ? Deux causes probables.
\item Deux méthodes de correction.
\end{enumerate}
\tcblower
\textbf{Solution.}
\begin{enumerate}
\item \textbf{Glissement de la courroie.} Causes : (1) Huile sur poulie motrice $\Rightarrow$ coefficient de frottement chute. (2) Courroie détendue ou usée (surface lisse).
\item (1) \textbf{Nettoyer} poulies et courroie à l'essence/dégraissant ; colmater la fuite d'huile du moteur. (2) \textbf{Tendre la courroie} (régler entraxe ou galet tendeur) ; si usée, \textbf{remplacer} par neuve même référence.
\end{enumerate}
\end{exoresolu}

\courssub{Sécurité et entretien}

\begin{methode}[Consignes de sécurité et d'entretien]
\begin{enumerate}
\item \textbf{Arrêt obligatoire :} Avant toute intervention (tension, nettoyage, remplacement), \textbf{arrêter le moteur} et \textbf{couper l'alimentation} (clé contact, disjoncteur, débrancher prise ENEO/groupe). Attendre l'arrêt complet.
\item \textbf{Carter de protection :} Ne \textbf{jamais} faire fonctionner la machine sans le carter (capot) en place. Il protège du happement et retient la courroie en cas de rupture.
\item \textbf{Tenue :} Pas de vêtements amples, pas de bijoux (bagues, bracelets, chaînes), cheveux longs attachés. Risque d'enroulement mortel.
\item \textbf{Entretien préventif (hebdo/mensuel) :} Vérifier tension (flèche 10-15 mm), état courroie (fissures, effilochage), propreté poulies, alignement des poulies (règle sur gorges), jeu des roulements.
\item \textbf{Remplacement :} Changer la courroie dès les premiers signes d'usure (ne pas attendre la rupture). Monter la neuve \textbf{sans forcer} (desserer l'entraxe au max) pour ne pas casser les fibres.
\end{enumerate}
\end{methode}

\begin{exoresolu}[L'apprenti de Douala]
Un apprenti veut retendre une courroie sur un moteur tournant. Le maître l'interdit.
\begin{enumerate}
\item Pourquoi c'est dangereux ?
\item Deux autres consignes de sécurité.
\end{enumerate}
\tcblower
\textbf{Solution.}
\begin{enumerate}
\item Risque d'\textbf{happement} : doigt, outil, vêtement pincé entre courroie et poulie $\rightarrow$ entraînement violent $\rightarrow$ écrasement, fracture, amputation. \textbf{Intervention interdite en mouvement.}
\item (1) Remettre le \textbf{carter de protection} avant redémarrage. (2) Porter une tenue ajustée, cheveux attachés, pas de bijoux. (3) Vérifier l'alignement des poulies (usure asymétrique = danger).
\end{enumerate}
\end{exoresolu}

\courssub{Entraînement : Exercices d'application}

\begin{exercice}[Application directe : Ventilateur industriel]
Un ventilateur d'atelier doit tourner à $N_2 = \SI{900}{tr/min}$. Le moteur électrique tourne à $N_1 = \SI{3000}{tr/min}$ et sa poulie motrice a un diamètre $D_1 = \SI{60}{mm}$.
\begin{enumerate}
\item Calculer le diamètre $D_2$ de la poulie du ventilateur.
\item Préciser si c'est un réducteur ou un multiplicateur.
\item Si la courroie est ouverte, les deux poulies tournent-elles dans le même sens ?
\end{enumerate}
\end{exercice}

\begin{exercice}[Diagnostic de panne : Pompe à eau]
Une pompe à eau entraînée par courroie trapézoïdale ne débite plus assez. On constate : courroie très chaude, odeur de brûlé, fine poudre noire sous le carter, poulie motrice brillante (lisse).
\begin{enumerate}
\item Quel est le défaut ? Justifier par deux observations.
\item Quelle est la cause probable de l'état de la poulie motrice ?
\item Proposer la procédure complète de remise en état (3 étapes).
\end{enumerate}
\end{exercice}

\begin{exercice}[Choix technologique : Scierie mobile]
Un menuisier itinérant veut motoriser une scie à ruban avec un moteur essence \SI{5,5}{ch} (\SI{3600}{tr/min}). La scie doit tourner à \SI{900}{tr/min}. L'entraxe est de \SI{80}{cm}.
\begin{enumerate}
\item Pourquoi le système poulies-courroie est-il le seul adapté ici ? (2 arguments).
\item Calculer le rapport de transmission $r$.
\item Si $D_1 = \SI{80}{mm}$, calculer $D_2$.
\item Citer un inconvénient majeur à surveiller sur ce chantier (poussière de bois, humidité).
\end{enumerate}
\end{exercice}

\begin{corrige}[Exercice 1 : Ventilateur]
\begin{enumerate}
\item $D_2 = \dfrac{D_1 \cdot N_1}{N_2} = \dfrac{60 \times 3000}{900} = \SI{200}{mm}$.
\item $D_2 > D_1$ ($200 > 60$) $\Rightarrow$ $N_2 < N_1$ : \textbf{Réducteur de vitesse}.
\item Oui, courroie ouverte $\Rightarrow$ \textbf{même sens}.
\end{enumerate}
\end{corrige}

\begin{corrige}[Exercice 2 : Pompe à eau]
\begin{enumerate}
\item \textbf{Glissement} de la courroie. Preuves : (1) Sifflement + odeur brûlé + poudre noire = échauffement par frottement glissant. (2) Poulie motrice brillante/lisse = usure par glissement répété (polissage).
\item Cause : Courroie \textbf{détendue} depuis longtemps (pas de retension) ou \textbf{surcharge} pompe (corps étranger) $\Rightarrow$ glissement continu $\rightarrow$ polissage de la gorge (perte d'adhérence définitive).
\item Procédure : (1) Arrêt machine, coupure alimentation, dépose carter. (2) \textbf{Remplacer la poulie motrice usée} (indispensable, une gorge polie fait glisser la neuve courroie) + remplacer la courroie. (3) Remonter, tendre correctement (flèche 10-15 mm), aligner poulies, remettre carter, test.
\end{enumerate}
\end{corrige}

\begin{corrige}[Exercice 3 : Scierie mobile]
\begin{enumerate}
\item (1) Entraxe \SI{80}{cm} trop grand pour engrenages. (2) Moteur essence = vibrations/chocs $\Rightarrow$ courroie absorbe les chocs (protection vilebrequin).
\item $r = \dfrac{N_2}{N_1} = \dfrac{900}{3600} = \dfrac{1}{4} = 0,25$.
\item $D_2 = \dfrac{D_1}{r} = \dfrac{80}{0,25} = \SI{320}{mm}$ (ou $D_2 = D_1 \times \frac{N_1}{N_2} = 80 \times 4 = 320$).
\item \textbf{Poussière de bois (sciure)} entre courroie et gorges $\Rightarrow$ agit comme abrasif + diminue frottement $\Rightarrow$ glissement + usure ultra-rapide. \textbf{Solution :} Carter fermé étanche + nettoyage quotidien + courroie trapézoïdale renforcée (type \"XP\" ou \"XPZ\").
\end{enumerate}
\end{corrige}

\begin{attention}[Les 5 erreurs qui coûtent des points au BEPC]
\begin{enumerate}
\item \textbf{Inverser le rapport de transmission.} $r = \dfrac{N_2}{N_1} = \dfrac{D_1}{D_2}$ (Diamètre \emph{Motrice} sur Diamètre \emph{Menee}). Écrire $r = D_2/D_1$ inverse le calcul. \textbf{Réflexe :} Grande poulie menée = tourne moins vite.
\item \textbf{Oublier « sans glissement ».} La formule $D_1 N_1 = D_2 N_2$ suppose adhérence parfaite. Si glissement : $N_2$ réel \textbf{inférieur} au calculé. Le signaler dans la conclusion.
\item \textbf{Confondre courroie ouverte / croisée.} Ouverte = même sens. Croisée = sens opposés. Regarder le schéma de l'énoncé, ne pas deviner.
\item \textbf{Mélanger les unités.} $D_1$ en cm et $D_2$ en mm $\Rightarrow$ erreur facteur 10. \textbf{Mettre tout dans la même unité} avant de calculer.
\item \textbf{Réponses « un mot » en technologie.} « Avantage : distance » $\rightarrow$ 0 point. Écrire : « L'avantage est de permettre la transmission entre deux arbres \textbf{éloignés} de \SI{1,5}{m}, ce que ne permettent pas les engrenages. » \textbf{Sécurité :} Ne jamais écrire « faire attention ». Écrire : « \textbf{Arrêter la machine et couper l'alimentation} avant toute intervention. »
\end{enumerate}
\end{attention}

\newpage

\coursec{Exercices}

\courssub{Je vérifie mes connaissances}

\begin{exercice}[QCM : choisir la bonne réponse]
Pour chaque question, une seule réponse est exacte. Recopie la lettre de la bonne réponse.
\begin{enumerate}
\item Dans un système poulies-courroie, la transmission du mouvement se fait :
\begin{enumerate}
\item[a)] par engrènement des dents ;
\item[b)] par adhérence de la courroie sur les poulies ;
\item[c)] par contact direct des deux poulies.
\end{enumerate}
\item Avec une courroie \cle{croisée}, la poulie menée tourne :
\begin{enumerate}
\item[a)] dans le même sens que la poulie menante ;
\item[b)] dans le sens inverse de la poulie menante ;
\item[c)] toujours plus vite que la poulie menante.
\end{enumerate}
\item Une poulie menante de petit diamètre entraîne une poulie menée de grand diamètre. Il s'agit :
\begin{enumerate}
\item[a)] d'une multiplication de vitesse ;
\item[b)] d'une réduction de vitesse ;
\item[c)] d'un blocage du système.
\end{enumerate}
\item Le système poulies-courroie peut jouer le rôle de \og fusible mécanique \fg{} car :
\begin{enumerate}
\item[a)] la courroie fond à la chaleur ;
\item[b)] la courroie peut glisser ou se rompre en cas de surcharge, protégeant le moteur ;
\item[c)] les poulies sont en plastique.
\end{enumerate}
\end{enumerate}
\end{exercice}

\begin{exercice}[Vrai ou faux, avec justification]
Réponds par \textbf{vrai} ou \textbf{faux}, puis justifie ta réponse en une phrase.
\begin{enumerate}
\item La courroie est l'organe qui relie les deux poulies et transmet le mouvement.
\item Avec une courroie non croisée (ouverte), les deux poulies tournent en sens inverse.
\item Le glissement de la courroie est un phénomène toujours souhaitable dans une machine.
\item Retendre une courroie détendue permet de réduire le glissement.
\end{enumerate}
\end{exercice}

\begin{exercice}[Définitions à compléter]
Recopie et complète les phrases suivantes avec les mots : \emph{menante}, \emph{menée}, \emph{courroie}, \emph{glissement}, \emph{tendeur}.
\begin{enumerate}
\item La poulie reliée au moteur est la poulie \dots ; celle qui reçoit le mouvement est la poulie \dots .
\item La \dots est l'organe flexible qui transmet le mouvement entre les deux poulies.
\item Le \dots est le défaut qui apparaît lorsque la courroie patine sur la poulie au lieu de l'entraîner.
\item Un galet \dots est un organe qui permet de régler la tension de la courroie.
\end{enumerate}
\end{exercice}

\begin{exercice}[Calcul direct de vitesse]
Sur le groupe électrogène d'un quartier de Douala, la poulie menante a un diamètre $D_1 = \SI{100}{mm}$ et tourne à $N_1 = \SI{900}{tr/min}$. Elle entraîne une poulie menée de diamètre $D_2 = \SI{300}{mm}$.
\begin{enumerate}
\item Rappelle la relation entre les vitesses de rotation et les diamètres des deux poulies (en l'absence de glissement).
\item Calcule la vitesse de rotation $N_2$ de la poulie menée.
\item S'agit-il d'une réduction ou d'une multiplication de vitesse ? Justifie.
\end{enumerate}
\end{exercice}

\courssub{Je m'entraîne}

\begin{exercice}[Lecture de schéma d'une transmission]
On donne le schéma d'une transmission par poulies et courroie \textbf{croisée} : la poulie 1 (menante) est reliée au moteur, la poulie 2 (menée) entraîne la machine, la courroie 3 relie les deux poulies.
\begin{enumerate}
\item Nomme les organes repérés par les numéros 1, 2 et 3.
\item La poulie menante tourne dans le sens des aiguilles d'une montre. Indique, en justifiant, le sens de rotation de la poulie menée.
\item Même question si la courroie est maintenant \textbf{non croisée}.
\end{enumerate}
\end{exercice}

\begin{exercice}[Moulin à maïs du village]
Le moteur d'un moulin à maïs de Bafoussam porte une poulie de diamètre $D_1 = \SI{80}{mm}$ tournant à $N_1 = \SI{1500}{tr/min}$. La meule doit tourner à $N_2 = \SI{500}{tr/min}$.
\begin{enumerate}
\item Calcule le diamètre $D_2$ de la poulie menée qu'il faut choisir.
\item Le mécanicien du village ne trouve qu'une poulie de $\SI{200}{mm}$. Quelle sera alors la vitesse réelle de la meule ?
\item Ce réglage convient-il ? Propose une solution technique pour obtenir la vitesse voulue.
\end{enumerate}
\end{exercice}

\begin{exercice}[Panne du ventilateur de voiture]
Au démarrage d'une voiture à Yaoundé, le conducteur entend un sifflement aigu : la courroie du ventilateur \cle{glisse} sur les poulies.
\begin{enumerate}
\item Cite deux causes possibles de ce glissement.
\item Pour chaque cause, propose une méthode de correction.
\item Explique pourquoi il est dangereux de rouler longtemps avec une courroie qui glisse.
\end{enumerate}
\end{exercice}

\begin{exercice}[Comparaison des systèmes de transmission]
Recopie et complète le tableau suivant en plaçant, pour chaque ligne, le système qui convient parmi : \emph{engrenages}, \emph{chaîne et pignons}, \emph{poulies et courroie}.
\begin{center}
\begin{tabularx}{\linewidth}{|X|l|}
\hline
\rowcolor{popblueL} \textbf{Caractéristique} & \textbf{Système} \\
\hline
Transmission par adhérence, pouvant glisser en cas de surcharge & \\
\hline
Transmission par engrènement de dents, sans glissement possible & \\
\hline
Transmission par maillons articulés, utilisée sur les moto-taxis & \\
\hline
Fonctionnement le plus silencieux & \\
\hline
\end{tabularx}
\end{center}
\end{exercice}

\courssub{Je me prépare au BEPC}

\begin{exercice}[Type BEPC : la machine à coudre de l'atelier]
Dans un atelier de couture de Bafang, une machine à coudre est entraînée par un moteur électrique grâce à un système poulies-courroie. Le schéma de la transmission comporte trois repères : 1 (poulie fixée sur l'arbre du moteur), 2 (poulie fixée sur l'arbre de la machine), 3 (organe flexible reliant 1 et 2). La courroie est \textbf{croisée}. La poulie 1 a un diamètre $D_1 = \SI{60}{mm}$ et tourne à $N_1 = \SI{1400}{tr/min}$ ; la poulie 2 a un diamètre $D_2 = \SI{210}{mm}$.
\begin{enumerate}
\item Nomme les organes repérés 1, 2 et 3, en précisant le rôle de chacun.
\item Le moteur tourne dans le sens des aiguilles d'une montre. Détermine, en justifiant, le sens de rotation de l'arbre de la machine à coudre.
\item Calcule la vitesse de rotation $N_2$ de l'arbre de la machine. Arrondis le résultat à l'unité.
\item Cite un avantage et un inconvénient de ce système de transmission par rapport à un système à engrenages.
\item Après plusieurs mois d'utilisation, la couturière remarque que la machine ralentit quand elle coud un tissu épais. Explique le phénomène et propose une méthode de correction.
\end{enumerate}
\end{exercice}

\begin{exercice}[Type BEPC : la pompe à eau du champ]
Pour irriguer son champ de tomates près de Nkongsamba, un agriculteur utilise une pompe à eau entraînée par un moteur à essence. La transmission se fait par deux poulies et une courroie non croisée. La poulie du moteur a un diamètre $D_1 = \SI{120}{mm}$ et tourne à $N_1 = \SI{1800}{tr/min}$. Pour fonctionner correctement, la pompe doit tourner entre $\SI{550}{tr/min}$ et $\SI{650}{tr/min}$.
\begin{enumerate}
\item Fais un schéma simple et légendé de cette transmission (poulies, courroie, sens de rotation).
\item Calcule le diamètre $D_2$ de la poulie de la pompe pour qu'elle tourne exactement à $N_2 = \SI{600}{tr/min}$.
\item L'agriculteur installe en réalité une poulie de $\SI{350}{mm}$. Calcule la vitesse obtenue et vérifie si elle se situe dans l'intervalle exigé.
\item En pleine saison sèche, la courroie s'est détendue et glisse. Cite deux méthodes de correction du glissement que l'agriculteur peut appliquer.
\item Rédige deux consignes de sécurité à respecter avant toute intervention sur cette transmission.
\end{enumerate}
\end{exercice}

\begin{exercice}[Type BEPC : sécurité et entretien en atelier scolaire]
L'atelier de technologie d'un collège de Bertoua possède une perceuse à colonne dont la transmission moteur-broche est assurée par un système poulies-courroie. Un élève constate que la courroie est détendue et doit être retendue.
\begin{enumerate}
\item Définis le système poulies-courroie et cite deux de ses avantages.
\item Explique ce qu'est le glissement d'une courroie et donne deux de ses causes.
\item Cite deux méthodes permettant de corriger le glissement, autres que le remplacement des poulies.
\item Rédige trois consignes de sécurité à respecter \textbf{avant} de retendre la courroie de cette machine d'atelier.
\item Le professeur hésite entre remplacer la transmission par courroie par une transmission à engrenages. Donne un argument pour chaque choix, puis conclus en justifiant ta proposition.
\end{enumerate}
\end{exercice}

\newpage

\coursec{Corrigés des exercices}

\begin{corrige}[Exercice 1 — QCM : choisir la bonne réponse]
\begin{enumerate}
\item \textbf{Réponse b)} : par adhérence de la courroie sur les poulies. La courroie est tendue sur les jantes des deux poulies ; c'est le frottement (l'adhérence) entre la courroie et les poulies qui transmet le mouvement. Il n'y a ni dents (engrenages), ni contact direct entre les poulies.
\item \textbf{Réponse b)} : dans le sens inverse de la poulie menante. Le croisement de la courroie inverse le sens de rotation entre les deux poulies.
\item \textbf{Réponse b)} : d'une réduction de vitesse. Quand la poulie menante est plus petite que la poulie menée ($D_1 < D_2$), la poulie menée tourne plus lentement : $N_2 = N_1 \times \dfrac{D_1}{D_2} < N_1$.
\item \textbf{Réponse b)} : la courroie peut glisser ou se rompre en cas de surcharge, protégeant le moteur. C'est cet effet \og fusible mécanique \fg{} qui évite la casse des organes coûteux de la machine.
\end{enumerate}
\end{corrige}

\begin{corrige}[Exercice 2 — Vrai ou faux, avec justification]
\begin{enumerate}
\item \textbf{Vrai.} La courroie est l'organe flexible tendu sur les deux poulies ; c'est elle qui relie les poulies et transmet le mouvement par adhérence.
\item \textbf{Faux.} Avec une courroie non croisée (ouverte), les deux poulies tournent dans le \cle{même sens} ; c'est la courroie croisée qui inverse le sens de rotation.
\item \textbf{Faux.} Le glissement est un défaut qui provoque une perte de vitesse, un échauffement et une usure rapide de la courroie ; il n'est souhaitable qu'en cas de surcharge, où il joue un rôle de protection (fusible mécanique).
\item \textbf{Vrai.} Retendre la courroie augmente son adhérence sur les poulies, ce qui réduit le patinage, donc le glissement.
\end{enumerate}
\end{corrige}

\begin{corrige}[Exercice 3 — Définitions à compléter]
\begin{enumerate}
\item La poulie reliée au moteur est la poulie \cle{menante} ; celle qui reçoit le mouvement est la poulie \cle{menée}.
\item La \cle{courroie} est l'organe flexible qui transmet le mouvement entre les deux poulies.
\item Le \cle{glissement} est le défaut qui apparaît lorsque la courroie patine sur la poulie au lieu de l'entraîner.
\item Un galet \cle{tendeur} est un organe qui permet de régler la tension de la courroie.
\end{enumerate}
\end{corrige}

\begin{corrige}[Exercice 4 — Calcul direct de vitesse]
\begin{enumerate}
\item En l'absence de glissement, les vitesses de rotation sont inversement proportionnelles aux diamètres :
\[ D_1 \times N_1 = D_2 \times N_2 \qquad \text{soit} \qquad N_2 = N_1 \times \frac{D_1}{D_2}. \]
\item Application numérique :
\[ N_2 = 900 \times \frac{100}{300} = 900 \times \frac{1}{3} = \SI{300}{tr/min}. \]
La poulie menée tourne à $\SI{300}{tr/min}$.
\item Il s'agit d'une \cle{réduction} de vitesse, car la poulie menée ($D_2 = \SI{300}{mm}$) est plus grande que la poulie menante ($D_1 = \SI{100}{mm}$) : elle tourne trois fois moins vite ($300 < 900$).
\end{enumerate}
\end{corrige}

\begin{corrige}[Exercice 5 — Lecture de schéma d'une transmission]
\begin{enumerate}
\item Repère 1 : poulie \cle{menante} (fixée sur l'arbre du moteur, elle donne le mouvement). Repère 2 : poulie \cle{menée} (fixée sur l'arbre de la machine, elle reçoit le mouvement). Repère 3 : la \cle{courroie} (organe flexible qui relie les deux poulies et transmet le mouvement).
\item La courroie est \textbf{croisée} : le croisement inverse le sens de rotation. La poulie menée tourne donc dans le sens \textbf{inverse} des aiguilles d'une montre.
\item Si la courroie est \textbf{non croisée} (ouverte), le sens de rotation est conservé : la poulie menée tourne dans le \textbf{même sens} que la poulie menante, c'est-à-dire dans le sens des aiguilles d'une montre.
\end{enumerate}
\end{corrige}

\begin{attention}[Point de vigilance]
Ne jamais répondre \og sens inverse \fg{} sans préciser la disposition de la courroie : c'est le croisement (ou non) de la courroie qui détermine le sens de rotation de la poulie menée.
\end{attention}

\begin{corrige}[Exercice 6 — Moulin à maïs du village]
\begin{enumerate}
\item On utilise la relation $D_1 \times N_1 = D_2 \times N_2$, d'où :
\[ D_2 = D_1 \times \frac{N_1}{N_2} = 80 \times \frac{1500}{500} = 80 \times 3 = \SI{240}{mm}. \]
Il faut une poulie menée de diamètre $\SI{240}{mm}$.
\item Avec $D_2 = \SI{200}{mm}$ :
\[ N_2 = N_1 \times \frac{D_1}{D_2} = 1500 \times \frac{80}{200} = 1500 \times 0{,}4 = \SI{600}{tr/min}. \]
La meule tournera à $\SI{600}{tr/min}$ au lieu de $\SI{500}{tr/min}$.
\item Ce réglage ne convient pas : la meule tourne trop vite ($600 > 500$), ce qui peut échauffer la farine, user la meule et provoquer des vibrations dangereuses. Solutions techniques possibles : trouver une poulie menée de $\SI{240}{mm}$, ou bien réduire la vitesse du moteur à $N_1 = N_2 \times \dfrac{D_2}{D_1} = 500 \times \dfrac{200}{80} = \SI{1250}{tr/min}$ avec la poulie de $\SI{200}{mm}$.
\end{enumerate}
\end{corrige}

\begin{corrige}[Exercice 7 — Panne du ventilateur de voiture]
\begin{enumerate}
\item Deux causes possibles du glissement :
\begin{itemize}
\item la courroie est \textbf{détendue} (tension insuffisante) ;
\item la courroie est \textbf{usée ou glacée} (flancs lisses, fissures), ou encrassée par de l'huile ou de la graisse.
\end{itemize}
\item Méthodes de correction :
\begin{itemize}
\item courroie détendue : \textbf{retendre la courroie} (en déplaçant le moteur ou l'alternateur sur ses supports, ou à l'aide d'un galet tendeur) ;
\item courroie usée ou encrassée : \textbf{nettoyer les poulies} et \textbf{remplacer la courroie} par une courroie neuve de même dimension.
\end{itemize}
\item Une courroie qui glisse patine sur les poulies : elle s'échauffe, s'use très vite et peut se rompre brutalement. Le ventilateur (ou l'alternateur, la pompe à eau selon les véhicules) n'est alors plus entraîné, ce qui peut provoquer la \textbf{surchauffe du moteur} et une panne coûteuse. Il faut donc intervenir rapidement.
\end{enumerate}
\end{corrige}

\begin{corrige}[Exercice 8 — Comparaison des systèmes de transmission]
\begin{center}
\begin{tabularx}{\linewidth}{|X|l|}
\hline
\rowcolor{popblueL} \textbf{Caractéristique} & \textbf{Système} \\
\hline
Transmission par adhérence, pouvant glisser en cas de surcharge & Poulies et courroie \\
\hline
Transmission par engrènement de dents, sans glissement possible & Engrenages \\
\hline
Transmission par maillons articulés, utilisée sur les moto-taxis & Chaîne et pignons \\
\hline
Fonctionnement le plus silencieux & Poulies et courroie \\
\hline
\end{tabularx}
\end{center}
Justifications : la courroie transmet par adhérence, donc elle peut glisser (rôle de fusible mécanique) ; les engrenages transmettent par engrènement de dents, sans glissement ; la chaîne de moto-taxi est formée de maillons articulés ; la courroie, souple et élastique, absorbe les vibrations et fonctionne sans bruit, contrairement aux engrenages et à la chaîne.
\end{corrige}

\begin{corrige}[Exercice 9 — Type BEPC : la machine à coudre de l'atelier]
\begin{enumerate}
\item Repère 1 : poulie \cle{menante}, fixée sur l'arbre du moteur ; elle reçoit le mouvement du moteur et le communique à la courroie. Repère 2 : poulie \cle{menée}, fixée sur l'arbre de la machine ; elle reçoit le mouvement et entraîne la machine à coudre. Repère 3 : la \cle{courroie}, organe flexible qui relie les deux poulies et transmet le mouvement par adhérence.
\item La courroie est \textbf{croisée} : le croisement inverse le sens de rotation. L'arbre de la machine tourne donc dans le sens \textbf{inverse} des aiguilles d'une montre.
\item Sans glissement : $N_2 = N_1 \times \dfrac{D_1}{D_2} = 1400 \times \dfrac{60}{210} = 1400 \times \dfrac{2}{7} = \dfrac{2800}{7} = \SI{400}{tr/min}$. L'arbre de la machine tourne à $\SI{400}{tr/min}$.
\item Avantage par rapport aux engrenages : fonctionnement \textbf{silencieux} et possibilité de glissement en cas de surcharge (protection du moteur, rôle de fusible mécanique) ; coût faible et entretien simple. Inconvénient : risque de \textbf{glissement} (perte de vitesse, usure de la courroie) et usure plus rapide de la courroie, qu'il faut remplacer périodiquement.
\item Quand la couturière coud un tissu épais, l'effort résistant augmente : la courroie, sans doute détendue ou usée, \textbf{glisse} sur les poulies, d'où le ralentissement de la machine. Correction : \textbf{retendre la courroie} (réglage de la position du moteur ou galet tendeur) et, si elle est usée, la \textbf{remplacer}.
\end{enumerate}
\textbf{Barème indicatif (sur 10)} : question 1 : 2 pts ; question 2 : 1,5 pt (sens + justification) ; question 3 : 2 pts (formule 1 pt, calcul 1 pt) ; question 4 : 2 pts ; question 5 : 2,5 pts (phénomène 1,5 pt, correction 1 pt).

\textbf{Points de vigilance} : citer la relation $D_1 N_1 = D_2 N_2$ avant le calcul ; justifier le sens de rotation par le croisement de la courroie ; pour la question 5, nommer explicitement le phénomène (\og glissement \fg).
\end{corrige}

\begin{corrige}[Exercice 10 — Type BEPC : la pompe à eau du champ]
\begin{enumerate}
\item Schéma de la transmission :
\begin{popfigure}
\begin{tikzpicture}[scale=0.9]
\draw[thick,fill=popblueL] (0,0) circle (0.8);
\draw[thick,fill=poporangeL] (5,0) circle (1.6);
\fill (0,0) circle (0.06); \fill (5,0) circle (0.06);
\draw[thick,popdark] (0,0.8) -- (5,1.6);
\draw[thick,popdark] (0,-0.8) -- (5,-1.6);
\draw[-{Stealth},thick,popblue] (0.8,0.5) arc (30:-30:1.05);
\draw[-{Stealth},thick,poporange] (6.6,0.9) arc (30:-30:1.95);
\node[below] at (0,-0.95) {\small Poulie menante (moteur)};
\node[below] at (5,-1.75) {\small Poulie men\'ee (pompe)};
\node[above,popdark] at (2.5,1.5) {\small Courroie non crois\'ee};
\end{tikzpicture}
\legende{Transmission poulies-courroie non croisée : la courroie relie les deux poulies ; les flèches indiquent que les deux poulies tournent dans le même sens.}
\end{popfigure}
\item Diamètre de la poulie de la pompe :
\[ D_2 = D_1 \times \frac{N_1}{N_2} = 120 \times \frac{1800}{600} = 120 \times 3 = \SI{360}{mm}. \]
\item Avec $D_2 = \SI{350}{mm}$ :
\[ N_2 = N_1 \times \frac{D_1}{D_2} = 1800 \times \frac{120}{350} = \frac{216000}{350} \approx \SI{617}{tr/min}. \]
Or $550 < 617 < 650$ : la vitesse obtenue se situe \textbf{dans l'intervalle exigé}, la pompe fonctionne correctement.
\item Deux méthodes de correction du glissement : \textbf{retendre la courroie} (rapprocher ou écarter les poulies en déplaçant le moteur sur ses supports, ou ajouter un galet tendeur) ; \textbf{nettoyer les poulies et la courroie} (éliminer huile, graisse et poussière) ou \textbf{remplacer la courroie} si elle est usée ou glacée.
\item Consignes de sécurité avant toute intervention :
\begin{itemize}
\item \textbf{Arrêter le moteur} et attendre l'arrêt complet des poulies ;
\item \textbf{Couper l'alimentation} (retirer la bougie ou la clé de contact du moteur à essence) pour empêcher tout redémarrage intempestif ;
\item ne jamais passer les mains près de la courroie en mouvement et porter des vêtements ajustés (gants retirés pendant la rotation, cheveux attachés).
\end{itemize}
\end{enumerate}
\textbf{Barème indicatif (sur 10)} : question 1 : 2 pts (schéma propre, légendé, sens de rotation) ; question 2 : 2 pts ; question 3 : 2 pts (calcul 1 pt, vérification de l'intervalle 1 pt) ; question 4 : 2 pts ; question 5 : 2 pts.

\textbf{Points de vigilance} : le schéma doit montrer une courroie \emph{non croisée} (brins parallèles) et deux flèches de même sens ; à la question 3, la conclusion doit comparer explicitement la valeur trouvée aux bornes $\SI{550}{tr/min}$ et $\SI{650}{tr/min}$ ; les consignes de sécurité doivent précéder l'intervention (arrêt et coupure de l'alimentation).
\end{corrige}

\begin{corrige}[Exercice 11 — Type BEPC : sécurité et entretien en atelier scolaire]
\begin{enumerate}
\item \textbf{Définition} : le système poulies-courroie est un système de transmission de mouvement de rotation entre deux arbres éloignés, composé de deux poulies (une menante, une menée) reliées par une courroie flexible tendue ; la transmission se fait par adhérence de la courroie sur les poulies. \textbf{Avantages} (deux au choix) : fonctionnement silencieux ; coût faible et entretien simple ; possibilité de relier des arbres éloignés ; absorption des vibrations ; protection de la machine en cas de surcharge (la courroie glisse ou se rompt : fusible mécanique).
\item Le \textbf{glissement} est le défaut qui apparaît lorsque la courroie patine sur la jante de la poulie au lieu de l'entraîner : la vitesse transmise diminue et la courroie chauffe et s'use. \textbf{Causes} : tension insuffisante de la courroie (courroie détendue) ; courroie usée, glacée ou fissurée ; présence d'huile ou de graisse sur les poulies ; surcharge de la machine.
\item Deux méthodes de correction (autres que le remplacement des poulies) : \textbf{retendre la courroie} (réglage de l'entraxe ou galet tendeur) ; \textbf{nettoyer les poulies et la courroie} ou \textbf{remplacer la courroie usée} par une courroie neuve.
\item Trois consignes de sécurité \textbf{avant} de retendre la courroie :
\begin{itemize}
\item arrêter la machine et \textbf{débrancher la prise électrique} (couper l'alimentation) ;
\item attendre l'\textbf{arrêt complet} des poulies et verrouiller l'interrupteur pour empêcher tout redémarrage ;
\item travailler avec des vêtements ajustés, cheveux attachés, et utiliser les outils adaptés (ne jamais tendre la courroie à la main pendant la rotation).
\end{itemize}
\item Argument pour la courroie : fonctionnement silencieux et protection du moteur en cas de surcharge (glissement-fusible), coût faible. Argument pour les engrenages : transmission \textbf{sans glissement}, vitesse exacte et durable, pas de courroie à remplacer. \textbf{Conclusion} : pour une perceuse à colonne d'atelier scolaire, on peut proposer de \textbf{conserver la courroie} : elle est silencieuse, protège le moteur en cas de blocage du foret (sécurité des élèves) et coûte peu cher à entretenir ; les engrenages seraient plus précis mais plus bruyants, plus coûteux et sans protection en cas de surcharge. (Toute conclusion argumentée et cohérente est acceptée.)
\end{enumerate}
\textbf{Barème indicatif (sur 10)} : question 1 : 2,5 pts (définition 1,5 pt, avantages 1 pt) ; question 2 : 2 pts ; question 3 : 1,5 pt ; question 4 : 2 pts ; question 5 : 2 pts (arguments 1 pt, conclusion justifiée 1 pt).

\textbf{Points de vigilance} : la définition doit mentionner les trois éléments (poulie menante, poulie menée, courroie) et le principe de l'adhérence ; les consignes de sécurité doivent être formulées \emph{avant} l'intervention (arrêt, coupure de l'alimentation, verrouillage) ; à la question 5, la conclusion doit être justifiée par les arguments cités.
\end{corrige}

\newpage

\coursec{Fiche bilan}

\begin{popfigure}
\centering
\begin{center}\tcbox[colback=poporange,colframe=poporange,coltext=white,arc=9pt,boxsep=4pt,left=6pt,right=6pt]{\bfseries\small Système poulies-courroie}\end{center}
\vspace{-2pt}
\begin{tcbraster}[raster columns=3,raster equal height=rows,raster column skip=6pt,raster row skip=6pt,enhanced,blankest]
\begin{tcolorbox}[enhanced,colback=white,colframe=poporange,boxrule=0.9pt,arc=3pt,title={Définition \& Composants},fonttitle=\bfseries\scriptsize,coltitle=white,colbacktitle=poporange,left=4pt,right=4pt,top=2pt,bottom=2pt,boxsep=1pt,toptitle=1.5pt,bottomtitle=1.5pt,halign=left]
\begin{itemize}[nosep,leftmargin=*,itemsep=1pt,font=\scriptsize]
\item Poulie motrice
\item Poulie réceptrice
\item Courroie
\end{itemize}
\end{tcolorbox}
\begin{tcolorbox}[enhanced,colback=white,colframe=popgreen,boxrule=0.9pt,arc=3pt,title={Fonctionnement},fonttitle=\bfseries\scriptsize,coltitle=white,colbacktitle=popgreen,left=4pt,right=4pt,top=2pt,bottom=2pt,boxsep=1pt,toptitle=1.5pt,bottomtitle=1.5pt,halign=left]
\begin{itemize}[nosep,leftmargin=*,itemsep=1pt,font=\scriptsize]
\item Sens rotation
\item Vitesse linéaire
\item Rapport $k$
\end{itemize}
\end{tcolorbox}
\begin{tcolorbox}[enhanced,colback=white,colframe=poppurple,boxrule=0.9pt,arc=3pt,title={Avantages},fonttitle=\bfseries\scriptsize,coltitle=white,colbacktitle=poppurple,left=4pt,right=4pt,top=2pt,bottom=2pt,boxsep=1pt,toptitle=1.5pt,bottomtitle=1.5pt,halign=left]
\begin{itemize}[nosep,leftmargin=*,itemsep=1pt,font=\scriptsize]
\item Distance > engrenages
\item Amortit chocs
\item Silencieux
\end{itemize}
\end{tcolorbox}
\begin{tcolorbox}[enhanced,colback=white,colframe=poppink,boxrule=0.9pt,arc=3pt,title={Inconvénients},fonttitle=\bfseries\scriptsize,coltitle=white,colbacktitle=poppink,left=4pt,right=4pt,top=2pt,bottom=2pt,boxsep=1pt,toptitle=1.5pt,bottomtitle=1.5pt,halign=left]
\begin{itemize}[nosep,leftmargin=*,itemsep=1pt,font=\scriptsize]
\item Glissement
\item Usure courroie
\item Encombrement
\end{itemize}
\end{tcolorbox}
\begin{tcolorbox}[enhanced,colback=white,colframe=popgold,boxrule=0.9pt,arc=3pt,title={Correction glissements},fonttitle=\bfseries\scriptsize,coltitle=white,colbacktitle=popgold,left=4pt,right=4pt,top=2pt,bottom=2pt,boxsep=1pt,toptitle=1.5pt,bottomtitle=1.5pt,halign=left]
\begin{itemize}[nosep,leftmargin=*,itemsep=1pt,font=\scriptsize]
\item Tendeur
\item Courroie crantée
\item Revêtement adhérant
\end{itemize}
\end{tcolorbox}
\begin{tcolorbox}[enhanced,colback=white,colframe=popblue!70!popdark,boxrule=0.9pt,arc=3pt,title={Sécurité \& Entretien},fonttitle=\bfseries\scriptsize,coltitle=white,colbacktitle=popblue!70!popdark,left=4pt,right=4pt,top=2pt,bottom=2pt,boxsep=1pt,toptitle=1.5pt,bottomtitle=1.5pt,halign=left]
\begin{itemize}[nosep,leftmargin=*,itemsep=1pt,font=\scriptsize]
\item Carter protection
\item Tension réglée
\item Arrêt avant intervention
\end{itemize}
\end{tcolorbox}
\end{tcbraster}

\legende{Carte mentale : Système poulies-courroie (3ème PCT)}
\end{popfigure}

\begin{bilan}
\textbf{Définitions clés}
\begin{itemize}
    \item \cle{Système poulies-courroie} : ensemble de deux poulies (motrice et réceptrice) reliées par une courroie flexible, transmettant un mouvement de rotation et une puissance entre deux axes parallèles éloignés.
    \item \cle{Poulie motrice} : poulie qui reçoit l'énergie (moteur, manivelle). Diamètre $D_1$, vitesse $N_1$ (tr/min).
    \item \cle{Poulie réceptrice} : poulie qui restitue l'énergie (machine entraînée). Diamètre $D_2$, vitesse $N_2$ (tr/min).
    \item \cle{Glissement} : différence de vitesse linéaire entre la courroie et la poulie (perte d'adhérence). Entraîne $N_2$ réel $<$ $N_2$ théorique.
\end{itemize}

\textbf{Formules à connaître par cœur (sans glissement)}
\begin{center}
\begin{tabularx}{\linewidth}{|>{\centering}X|>{\centering}X|>{\centering\arraybackslash}X|}
\hline
\rowcolor{popblueL}
\textbf{Grandeur} & \textbf{Formule} & \textbf{Unités} \\
\hline
Rapport de transmission & $k = \dfrac{D_1}{D_2} = \dfrac{N_2}{N_1}$ & Sans unité \\
\hline
Vitesse linéaire courroie & $v = \dfrac{\pi \cdot D_1 \cdot N_1}{60} = \dfrac{\pi \cdot D_2 \cdot N_2}{60}$ & m/s ($D$ en m, $N$ en tr/min) \\
\hline
Puissance transmise & $P = F_u \cdot v$ & W ($F_u$ en N, $v$ en m/s) \\
\hline
Couple / Puissance & $P = C \cdot \omega \quad (\omega = \frac{2\pi N}{60})$ & W ($C$ en N$\cdot$m, $\omega$ en rad/s) \\
\hline
\end{tabularx}
\end{center}

\textbf{Méthodes express}
\begin{enumerate}
    \item \textbf{Calculer $N_2$ :} Identifier $D_1, D_2, N_1$. Appliquer $N_2 = N_1 \times \frac{D_1}{D_2}$. Vérifier ordre de grandeur (poulie plus petite $\rightarrow$ plus rapide).
    \item \textbf{Sens de rotation :} Courroie \emph{non croisée} $\rightarrow$ même sens. Courroie \emph{croisée} $\rightarrow$ sens contraires (inversion).
    \item \textbf{Choisir courroie :} Puissance à transmettre, distance entre axes, vitesse, environnement (huile, chaleur, poussière).
\end{enumerate}

\textbf{Points de vigilance}
\begin{itemize}
    \item Unités : $D$ en \textbf{mètres} pour $v$ (m/s) ; $N$ en \textbf{tr/min} pour $k$.
    \item Glissement : toujours présent (sauf courroie crantée/synchrone). Le $k$ réel est légèrement inférieur au $k$ théorique côté récepteur.
    \item Tension : trop faible $\rightarrow$ glissement/usure ; trop forte $\rightarrow$ rupture/usure roulements.
\end{itemize}
\end{bilan}

\begin{aretenir}[Checklist avant le Bac]
\begin{itemize}
    \item $\square$ Je définis le système poulies-courroie et nomme ses 3 composants principaux.
    \item $\square$ Je distingue poulie motrice et poulie réceptrice sur un schéma.
    \item $\square$ Je calcule le rapport de transmission $k = D_1/D_2 = N_2/N_1$ (sans glissement).
    \item $\square$ Je calcule la vitesse de rotation $N_2$ connaissant $N_1, D_1, D_2$.
    \item $\square$ Je calcule la vitesse linéaire $v$ de la courroie ($v = \pi D N / 60$).
    \item $\square$ Je détermine le sens de rotation (croisée / non croisée).
    \item $\square$ Je cite 3 avantages (distance, amortissement, silence, coût) et 3 inconvénients (glissement, usure, encombrement, tension).
    \item $\square$ J'explique le phénomène de glissement et ses conséquences sur $N_2$.
    \item $\square$ Je propose 3 solutions pour limiter le glissement (tendeur, courroie crantée, augmenter adhérence/angle d'enroulement).
    \item $\square$ J'applique les consignes de sécurité : carter de protection, arrêt machine avant intervention, vérification tension/état courroie.
    \item $\square$ Je donne 2 exemples camerounais : alternateur groupe électrogène, compresseur climatisation voiture, pompe à eau moto-pompe.
\end{itemize}
\end{aretenir}

\findecours

\end{document}
